% Further Complex Numbers — Further Mathematics Upper Sixth — Cours signé OnBuch+
% Official MINESEC syllabus. Compiler avec : tectonic FMA02-further-complex-numbers.tex
\newcommand{\DOCMATIERE}{Further Mathematics}
\newcommand{\DOCNIVEAU}{Upper Sixth}
\newcommand{\DOCMODULE}{Upper Sixth — Further mathematics}
\newcommand{\DOCLECON}{2}
\newcommand{\DOCTITRE}{Further Complex Numbers}
% OnBuch+ — préambule « cours signés » ENRICHI (compilé avec tectonic / XeLaTeX)
% Système visuel « Pop » d'OnBuch+ : fond crème (jamais blanc), bordures encre
% épaisses, ombres « dures » décalées, titres Archivo Black en capitales.
% Version MATHÉMATIQUES : graphes (pgfplots/tikz), boîtes pédagogiques
% (définition, propriété, méthode, attention, le savais-tu, exercice résolu,
% exercice, TP, bilan), page de garde et signature OnBuch+ ancrée en bas.
%
% Macros attendues dans main.tex AVANT \input{preamble} :
%   \DOCMATIERE  \DOCNIVEAU  \DOCMODULE  \DOCLECON (numéro)  \DOCTITRE
\documentclass[11pt]{article}

\usepackage{fontspec}
\usepackage[english]{babel}
\usepackage[a4paper,margin=2cm,top=3.4cm,bottom=2.8cm,headheight=1.4cm,footskip=1.3cm]{geometry}
\usepackage{amsmath,amssymb,amsfonts}
\usepackage{mathtools} % \xmapsto, \coloneqq... souvent utilisés par les modèles
\usepackage{cancel} % simplifications barrées (bilans de forces)
% \abs{z} (module, valeur absolue) et \norm{u} (norme), utilisés par les modèles
\providecommand{\abs}{}\let\abs\relax\DeclarePairedDelimiter{\abs}{\lvert}{\rvert}
\providecommand{\norm}{}\let\norm\relax\DeclarePairedDelimiter{\norm}{\lVert}{\rVert}
\usepackage[table]{xcolor} % [table] : fournit aussi \rowcolors (alternance de lignes)
\usepackage{pagecolor}
\usepackage{tikz}
\usetikzlibrary{arrows.meta,positioning,calc,shapes.geometric,shapes.misc,shapes.arrows,decorations.pathmorphing,decorations.pathreplacing,decorations.markings,patterns,shadows,mindmap,trees,fit,backgrounds,matrix,intersections,angles,quotes}
\usepackage{pgfplots}
\usepackage[version=4]{mhchem} % équations nucléaires \ce{^{235}_{92}U}
\usepackage[european,siunitx]{circuitikz} % schémas électriques
\pgfplotsset{compat=1.18}
\usepgfplotslibrary{fillbetween,groupplots} % aires entre courbes (calcul intégral)
\usepackage{enumitem}
\usepackage{fancyhdr}
% [most] charge toutes les bibliothèques tcolorbox (listings, minted, raster,
% documentation...) et gonfle inutilement la mémoire TeX sur un document long
% (~300 boîtes) : on ne charge que ce qui est réellement utilisé.
\usepackage[skins,breakable]{tcolorbox}
\usepackage{graphicx}
\usepackage{colortbl}
\usepackage{booktabs}
\usepackage{tabularx}
\usepackage{diagbox} % case d'en-tête diagonale des tableaux à double entrée
% Colonnes extensibles centrée / gauche / droite (C, L, R), souvent utilisées par les modèles
\newcolumntype{C}{>{\centering\arraybackslash}X}
\newcolumntype{L}{>{\raggedright\arraybackslash}X}
\newcolumntype{R}{>{\raggedleft\arraybackslash}X}
\usepackage{array}
\usepackage{multicol}
\usepackage{siunitx}
% parse-numbers=false : accepte aussi \SI{2,5 \times 10^{-3}}{...}
\sisetup{locale=UK,output-decimal-marker={.},group-minimum-digits=4,parse-numbers=false}
\DeclareSIUnit\ohm{\ensuremath{\symup{\Omega}}} % U+2126 absent de la police
\DeclareSIUnit\division{div} % réglages d'oscilloscope (ms/div, V/div)
\DeclareSIUnit\tr{tr} % tours (vitesse de rotation, ex. \SI{1500}{\tr\per\minute})
\usepackage{qrcode}
% \degree : commande fréquemment utilisée par les modèles pour un angle ou une
% température en dehors de \SI{}{} (non fournie par les paquets chargés).
\providecommand{\degree}{^{\circ}}
% \qed (fin de démonstration), utilisé par les modèles sans amsthm
\providecommand{\qed}{\hfill$\square$}
% \barre : double barre des valeurs interdites dans un tableau de variations (tkz-tab)
\providecommand{\barre}{\ensuremath{\|}}
% environnement proof (amsthm non chargé) utilisé par les modèles
\ifdefined\proof\else\newenvironment{proof}[1][Proof]{\par\noindent\textit{#1.}\ }{\qed\par}\fi
\usepackage{lastpage}
\usepackage{hyperref}
\hypersetup{hidelinks}

% ============================================================
% Palette « Pop » OnBuch+ — mêmes hex que la classe P (plus_ui.dart)
% ============================================================
\definecolor{poporange}{HTML}{F59321}
\definecolor{popdark}{HTML}{E07A0C}
\definecolor{popcream}{HTML}{FFF6E8}
\definecolor{popink}{HTML}{1C1714}
\definecolor{poppurple}{HTML}{7A5AE0}
\definecolor{popgreen}{HTML}{2E9E6A}
\definecolor{poppink}{HTML}{E0457B}
\definecolor{popblue}{HTML}{2F7DE1}
\definecolor{popgold}{HTML}{C98A12}
\definecolor{popmuted}{HTML}{8A7F76}
\definecolor{popline}{HTML}{EADFD2}
\definecolor{popgray}{HTML}{8A7F76}
\colorlet{popgrey}{popgray}
% Alias tolérants (le modèle nomme parfois les couleurs différemment)
\colorlet{popwhite}{white}
\colorlet{popblack}{popink}
\colorlet{popred}{poppink}
\colorlet{popyellow}{popgold}
% Teintes claires pour fonds de tableaux / zones
\colorlet{popblueL}{popblue!10!white}
\colorlet{poporangeL}{poporange!14!white}
\colorlet{popgreenL}{popgreen!12!white}
\colorlet{poppurpleL}{poppurple!12!white}
\colorlet{poppinkL}{poppink!12!white}
\colorlet{popgoldL}{popgold!14!white}

\pagecolor{popcream}

% ============================================================
% Typographie
% ============================================================
\defaultfontfeatures{Ligatures=TeX}
\setmainfont{PlusJakartaSans-Medium.ttf}[Path=../../../fonts/,BoldFont=PlusJakartaSans-Bold.ttf]
% Maths en sans-serif (Fira Math) avec chiffres et lettres droites de Plus
% Jakarta, pour que les formules se fondent dans le texte.
\usepackage[math-style=ISO,bold-style=ISO]{unicode-math}
\setmathfont{FiraMath-Regular.otf}[Path=../../../fonts/]
\setmathfont{PlusJakartaSans-Medium.ttf}[Path=../../../fonts/,range={up/num,up/{Latin,latin},\mathup}]
% (icomma retiré : décimales avec point en anglais)
% Caractères mathématiques Unicode absents des polices de texte (Plus Jakarta,
% Archivo Black) : rendus par leur équivalent mathématique, en texte comme en
% formule — sinon ils s'affichent en carré vide (ex. « Arithmétique dans ℤ »).
\usepackage{newunicodechar}
\newunicodechar{ℝ}{\ensuremath{\mathbb{R}}}
\newunicodechar{ℤ}{\ensuremath{\mathbb{Z}}}
\newunicodechar{ℂ}{\ensuremath{\mathbb{C}}}
\newunicodechar{ℕ}{\ensuremath{\mathbb{N}}}
\newunicodechar{ℚ}{\ensuremath{\mathbb{Q}}}
\newunicodechar{𝔻}{\ensuremath{\mathbb{D}}}
\newunicodechar{α}{\ensuremath{\alpha}}
\newunicodechar{β}{\ensuremath{\beta}}
\newunicodechar{γ}{\ensuremath{\gamma}}
\newunicodechar{δ}{\ensuremath{\delta}}
\newunicodechar{ε}{\ensuremath{\varepsilon}}
\newunicodechar{θ}{\ensuremath{\theta}}
\newunicodechar{λ}{\ensuremath{\lambda}}
\newunicodechar{μ}{\ensuremath{\mu}}
\newunicodechar{σ}{\ensuremath{\sigma}}
\newunicodechar{τ}{\ensuremath{\tau}}
\newunicodechar{φ}{\ensuremath{\varphi}}
\newunicodechar{ψ}{\ensuremath{\psi}}
\newunicodechar{ω}{\ensuremath{\omega}}
\newunicodechar{Σ}{\ensuremath{\Sigma}}
% majuscules produites par \MakeUppercase dans les titres (α-aminés -> Α)
\newunicodechar{Α}{\ensuremath{\alpha}}
\newunicodechar{Β}{\ensuremath{\beta}}
\newunicodechar{Γ}{\ensuremath{\Gamma}}
\newunicodechar{Δ}{\ensuremath{\Delta}}
\newunicodechar{Ε}{\ensuremath{\varepsilon}}
\newunicodechar{Θ}{\ensuremath{\Theta}}
\newunicodechar{Λ}{\ensuremath{\Lambda}}
\newunicodechar{Μ}{\ensuremath{\mu}}
\newunicodechar{Τ}{\ensuremath{\tau}}
\newunicodechar{Φ}{\ensuremath{\Phi}}
\newunicodechar{Ψ}{\ensuremath{\Psi}}
\newunicodechar{Ω}{\ensuremath{\Omega}}
\newunicodechar{↦}{\ensuremath{\mapsto}}
\newunicodechar{∈}{\ensuremath{\in}}
\newunicodechar{∉}{\ensuremath{\notin}}
\newunicodechar{∧}{\ensuremath{\wedge}}
\newunicodechar{⇔}{\ensuremath{\Leftrightarrow}}
\newunicodechar{⟺}{\ensuremath{\Longleftrightarrow}}
\newunicodechar{⇒}{\ensuremath{\Rightarrow}}
\newunicodechar{⟹}{\ensuremath{\Longrightarrow}}
\newunicodechar{∘}{\ensuremath{\circ}}
\newunicodechar{∪}{\ensuremath{\cup}}
\newunicodechar{∩}{\ensuremath{\cap}}
\newunicodechar{≡}{\ensuremath{\equiv}}
\newunicodechar{⊂}{\ensuremath{\subset}}
\newunicodechar{⊥}{\ensuremath{\perp}}
\newunicodechar{∃}{\ensuremath{\exists}}
\newunicodechar{∀}{\ensuremath{\forall}}
\newunicodechar{∛}{\ensuremath{\sqrt[3]{\,}}}
\newunicodechar{∜}{\ensuremath{\sqrt[4]{\,}}}
\newunicodechar{⃗}{\ensuremath{{}^{\to}}}
\newunicodechar{ⁿ}{\ifmmode^{n}\else\textsuperscript{\ensuremath{n}}\fi}
\newunicodechar{ˣ}{\ifmmode^{x}\else\textsuperscript{\ensuremath{x}}\fi}
\newunicodechar{ᵇ}{\ifmmode^{b}\else\textsuperscript{\ensuremath{b}}\fi}
\newunicodechar{ʳ}{\ifmmode^{r}\else\textsuperscript{\ensuremath{r}}\fi}
\newunicodechar{ᵘ}{\ifmmode^{u}\else\textsuperscript{\ensuremath{u}}\fi}
\newunicodechar{ʸ}{\ifmmode^{y}\else\textsuperscript{\ensuremath{y}}\fi}
\newunicodechar{ᵃ}{\ifmmode^{a}\else\textsuperscript{\ensuremath{a}}\fi}
\newunicodechar{ᵉ}{\ifmmode^{e}\else\textsuperscript{\ensuremath{e}}\fi}
\newunicodechar{ᵏ}{\ifmmode^{k}\else\textsuperscript{\ensuremath{k}}\fi}
\newunicodechar{ᵖ}{\ifmmode^{p}\else\textsuperscript{\ensuremath{p}}\fi}
\newunicodechar{ᴺ}{\ifmmode^{N}\else\textsuperscript{\ensuremath{N}}\fi}
\newunicodechar{ᶜ}{\ifmmode^{c}\else\textsuperscript{\ensuremath{c}}\fi}
\newunicodechar{ʰ}{\ifmmode^{h}\else\textsuperscript{\ensuremath{h}}\fi}
\newunicodechar{ᵠ}{\ifmmode^{q}\else\textsuperscript{\ensuremath{q}}\fi}
\newunicodechar{ₙ}{\ifmmode_{n}\else\textsubscript{\ensuremath{n}}\fi}
\newunicodechar{ₐ}{\ifmmode_{a}\else\textsubscript{\ensuremath{a}}\fi}
\newunicodechar{ᵢ}{\ifmmode_{i}\else\textsubscript{\ensuremath{i}}\fi}
\newunicodechar{ᵣ}{\ifmmode_{r}\else\textsubscript{\ensuremath{r}}\fi}
\newunicodechar{ₖ}{\ifmmode_{k}\else\textsubscript{\ensuremath{k}}\fi}
\newunicodechar{ᵦ}{\ifmmode_{\beta}\else\textsubscript{\ensuremath{\beta}}\fi}
\newunicodechar{ₓ}{\ifmmode_{x}\else\textsubscript{\ensuremath{x}}\fi}
\newfontfamily\popdisplay{ArchivoBlack-Regular.ttf}[Path=../../../fonts/]
\newfontfamily\poplogo{SpaceGrotesk-Bold.ttf}[Path=../../../fonts/]
\newfontfamily\popsemi{PlusJakartaSans-SemiBold.ttf}[Path=../../../fonts/]
\newfontfamily\popxbold{PlusJakartaSans-ExtraBold.ttf}[Path=../../../fonts/]
\newfontfamily\popxboldls{PlusJakartaSans-ExtraBold.ttf}[Path=../../../fonts/,LetterSpace=3.0]

\setlist[enumerate]{leftmargin=1.7em,itemsep=5pt,topsep=4pt}
\setlist[itemize]{leftmargin=1.7em,itemsep=4pt,topsep=4pt,label={\color{poporange}\textbullet}}
\setlength{\parindent}{0pt}
\setlength{\parskip}{5pt}
\renewcommand{\arraystretch}{1.25}

% pgfplots : style Pop par défaut
\pgfplotsset{
  every axis/.append style={axis line style={popink,line width=1pt},tick style={popink},
    grid style={popline,line width=0.5pt},label style={font=\small},tick label style={font=\footnotesize}},
  popcurve/.style={poporange,line width=1.8pt,no markers,smooth},
}
% Même style disponible dans un \draw TikZ simple (hors pgfplots)
\tikzset{popcurve/.style={poporange,line width=1.8pt}}
\tikzset{popgrid/.style={popline,very thin}}
\colorlet{popcurve}{poporange} % le nom du style est parfois utilisé comme couleur

% ============================================================
% Pill / squiggle
% ============================================================
\newcommand{\poppill}[3]{%
  \tikz[baseline=-0.55ex]\node[draw=popink,line width=1.2pt,rounded corners=2.6pt,%
    fill=#2,inner xsep=6.5pt,inner ysep=3pt]{%
    \popxboldls\fontsize{7}{7}\selectfont\color{#3}\MakeUppercase{#1}};%
}
\newcommand{\popsquiggle}[1]{%
  \tikz[baseline=-0.4ex]{
    \draw[#1,line width=2.6pt,line cap=round]
      plot [smooth] coordinates {(0,0.09) (0.34,0.01) (0.68,0.21) (1.02,0.01) (1.36,0.10)};
  }%
}
% Mot-clé surligné (marqueur orange sous le texte)
\newcommand{\cle}[1]{{\popxbold\color{popdark}#1}}

% ============================================================
% En-tête / pied
% ============================================================
\pagestyle{fancy}
\fancyhf{}
\renewcommand{\headrulewidth}{0pt}
\renewcommand{\footrulewidth}{0pt}
\lhead{\raisebox{-0.15ex}{\poplogo\fontsize{13}{13}\selectfont\color{popink}OnBuch\color{poporange}+}}
\rhead{\poppill{\DOCMATIERE\ \textbullet\ \DOCNIVEAU\ \textbullet\ Chapter \DOCLECON}{white}{popink}}
\lfoot{\popxbold\fontsize{7.6}{7.6}\selectfont\color{popmuted}\MakeUppercase{Signed course OnBuch+}\ \textbullet\ {\popsemi\DOCTITRE}}
\rfoot{\tikz[baseline=-0.6ex]\node[rounded corners=8.5pt,fill=popink,minimum height=17pt,minimum width=17pt,inner xsep=5pt,inner sep=0]{\popxbold\fontsize{8}{8}\selectfont\color{popcream}\thepage\,/\,\pageref*{LastPage}};}

% ============================================================
% Titres
% ============================================================
\newcommand{\chaptitre}[2]{%
  \begin{tcolorbox}[enhanced,colback=poporange,colframe=popink,boxrule=2.6pt,arc=11pt,
      drop shadow=popink,left=15pt,right=15pt,top=12pt,bottom=11pt,before skip=3pt,after skip=18pt]
    \poppill{#2}{popink}{popcream}\par\vspace{9pt}
    {\raggedright\hyphenpenalty=10000\exhyphenpenalty=10000\popdisplay\fontsize{22}{24}\selectfont\color{popcream}\MakeUppercase{#1}\par}
  \end{tcolorbox}
}
% Section numérotée : pastille orange avec le numéro + titre Archivo
\newcounter{popsec}
\newcommand{\coursec}[1]{%
  \stepcounter{popsec}\setcounter{popsub}{0}%
  \par\vspace{16pt}\noindent
  \begin{minipage}{\linewidth}
  \tikz[baseline=-0.7ex]\node[draw=popink,line width=1.6pt,fill=poporange,rounded corners=4pt,minimum size=22pt,inner sep=0]{\popdisplay\fontsize{12}{12}\selectfont\color{popcream}\thepopsec};%
  \hspace{8pt}{\popdisplay\fontsize{15}{17}\selectfont\color{popink}\MakeUppercase{#1}}\par
  \vspace{4pt}\noindent{\color{poporange}\rule{36pt}{3.6pt}}
  \end{minipage}\par\nopagebreak\vspace{8pt}
}
\newcounter{popsub}
\newcommand{\courssub}[1]{%
  \stepcounter{popsub}%
  \par\vspace{9pt}\noindent{\popxbold\fontsize{11.5}{13}\selectfont\color{popink}{\color{poporange}\thepopsec.\thepopsub}\hspace{6pt}#1}\par\nopagebreak\vspace{4pt}
}

% ============================================================
% Boîtes pédagogiques — même recette : fond blanc, bordure épaisse colorée,
% ombre dure encre, badge plein. Titre optionnel [#1].
% ============================================================
\tcbset{popbase/.style={breakable,enhanced,colback=white,boxrule=2.3pt,arc=9pt,
  shadow={3pt}{-3pt}{0pt}{popink},left=14pt,right=14pt,top=11pt,bottom=11pt,before skip=11pt,after skip=15pt,
  fontupper=\popsemi\fontsize{10.6}{14.5}\selectfont\color{popink},
  fontlower=\popsemi\fontsize{10.6}{14.5}\selectfont\color{popink}}}
% Coche dessinée (le glyphe ✓ n'existe pas dans les polices du document)
\newcommand{\popcheck}[1]{\tikz[baseline=-0.5ex]\draw[#1,line width=1.6pt,line cap=round,line join=round](-0.13,0.0)--(-0.04,-0.09)--(0.13,0.1);}
\providecommand{\coche}{\popcheck{popgreen}}
\providecommand{\resultat}[1]{\par\smallskip\noindent\fcolorbox{popgreen}{popgreenL}{\parbox{\dimexpr\linewidth-2\fboxsep-2\fboxrule\relax}{\textbf{Result.} #1}}\par\smallskip}
\newcommand{\popboxhead}[3]{\poppill{#1}{#2}{white}\ifx\relax#3\relax\else\hspace{6pt}{\popxbold\fontsize{10.6}{12}\selectfont\color{popink}#3}\fi\par\vspace{7pt}}

\newtcolorbox{aretenir}[1][]{popbase,colframe=popgreen,before upper={\popboxhead{Key points}{popgreen}{#1}}}
\newtcolorbox{exemplebox}[1][]{popbase,colframe=poporange,before upper={\popboxhead{Exemple}{poporange}{#1}}}
\newtcolorbox{definition}[1][]{popbase,colframe=popblue,before upper={\popboxhead{Definition}{popblue}{#1}}}
\newtcolorbox{propriete}[1][]{popbase,colframe=poppurple,before upper={\popboxhead{Property}{poppurple}{#1}}}
\newtcolorbox{methode}[1][]{popbase,colframe=popgold,before upper={\popboxhead{Method}{popgold}{#1}}}
\newtcolorbox{attention}[1][]{popbase,colframe=poppink,before upper={\popboxhead{Warning}{poppink}{#1}}}
\newtcolorbox{experience}[1][]{popbase,colframe=popblue,colback=popblueL,before upper={\popboxhead{Experiment}{popblue}{#1}}}
% « Le savais-tu ? » : carte violette pleine (contexte, culture, Cameroun)
\newtcolorbox{savaistu}[1][]{popbase,colback=poppurple,colframe=popink,
  fontupper=\popsemi\fontsize{10.4}{14.2}\selectfont\color{white},
  before upper={\poppill{Did you know?}{white}{poppurple}\ifx\relax#1\relax\else\hspace{6pt}{\popxbold\color{white}#1}\fi\par\vspace{7pt}}}
% Exercice résolu : énoncé en haut, \tcblower puis la solution sur fond crème
\newcounter{popexo}\newcounter{popexr}
\newtcolorbox{exoresolu}[1][]{popbase,colframe=poporange,colbacklower=popcream,
  segmentation style={popink,line width=1.2pt,dashed},
  before upper={\stepcounter{popexr}\popboxhead{Worked example \thepopexr}{poporange}{#1}},
  before lower={\poppill{Solution}{popink}{popcream}\par\vspace{6pt}}}
% Exercice (non résolu en place) — la correction est dans la partie Corrigés
\newtcolorbox{exercice}[1][]{popbase,colframe=popink,
  before upper={\stepcounter{popexo}\popboxhead{Exercise \thepopexo}{popink}{#1}}}
% Corrigé
\newtcolorbox{corrige}[1][]{popbase,colframe=popgreen,colback=popgreenL,
  before upper={\popboxhead{Solution}{popgreen}{#1}}}
% Objectifs / prérequis / bilan
\newtcolorbox{objectifs}{popbase,colframe=popink,colback=poporangeL,
  before upper={\popboxhead{Chapter objectives}{popink}{}}}
\newtcolorbox{prerequis}{popbase,colframe=popink,colback=popblueL,
  before upper={\popboxhead{Prerequisites}{popblue}{}}}
\newtcolorbox{bilan}{popbase,colframe=popink,colback=white,boxrule=2.8pt,
  before upper={\popboxhead{Fiche bilan}{poporange}{L'essentiel à connaître pour l'examen}}}
% Figure encadrée avec légende
\newtcolorbox{popfigure}[1][]{enhanced,breakable=false,colback=white,colframe=popline,boxrule=1.4pt,arc=8pt,
  left=10pt,right=10pt,top=10pt,bottom=8pt,before skip=10pt,after skip=12pt,halign=center,
  lower separated=false,halign lower=center,
  fontlower=\popsemi\fontsize{9}{11.5}\selectfont\color{popmuted},
  #1}
\newcommand{\legende}[1]{\par\vspace{6pt}{\popsemi\fontsize{9}{11.5}\selectfont\color{popmuted}#1}}

% ============================================================
% Signature OnBuch+
% ============================================================
\newcommand{\poplogoinline}[1]{{\poplogo\fontsize{#1}{#1}\selectfont\color{popink}OnBuch\color{poporange}+}}
\newcommand{\signatureonbuch}{%
  \begin{tcolorbox}[enhanced,colback=popink,colframe=popink,boxrule=2.2pt,arc=10pt,
      drop shadow=poporange,left=14pt,right=14pt,top=10pt,bottom=10pt,before skip=0pt,after skip=0pt]
    \tikz[baseline=-0.7ex]\node[circle,fill=popgreen,draw=popcream,line width=1.3pt,minimum size=22pt,inner sep=0]{\popcheck{white}};%
    \hspace{9pt}\begin{minipage}[c]{0.62\linewidth}
      {\popxbold\fontsize{12}{13}\selectfont\color{popcream}Signed\ \textbullet\ {\poplogo OnBuch\color{poporange}+}}\par\vspace{2pt}
      {\raggedright\popsemi\fontsize{8}{10.5}\selectfont\color{popline}\MakeUppercase{OnBuch+ teaching team}\\\MakeUppercase{Follows the official MINESEC syllabus}\par}
    \end{minipage}\hfill
    {\poplogo\fontsize{22}{22}\selectfont\color{popcream}OnBuch\color{poporange}+}
  \end{tcolorbox}%
}

% ============================================================
% Page de garde
% #1 titre · #2 matière · #3 niveau · #4 module · #5 n° leçon · #6 durée ·
% #7 description / objectifs (texte ou itemize)
% ============================================================
\newcommand{\pagegarde}[7]{%
  \thispagestyle{empty}
  \newgeometry{a4paper,margin=1.7cm,top=1.6cm,bottom=1.4cm}%
  \begin{tikzpicture}[remember picture,overlay]
    \fill[poporange!18] ([xshift=-3.2cm,yshift=-3.6cm]current page.north east) circle (4.6cm);
    \fill[poppurple!14] ([xshift=2.4cm,yshift=7.6cm]current page.south west) circle (3.8cm);
  \end{tikzpicture}%
  {\poplogo\fontsize{30}{30}\selectfont\color{popink}OnBuch\color{poporange}+}\hfill
  \raisebox{0.6ex}{\poppill{Signed course \textbullet\ Premium}{popink}{popcream}}\par
  \vspace{1pt}\popsquiggle{poppurple}\par
  \vspace{8pt}
  \begin{tcolorbox}[enhanced,colback=poporange,colframe=popink,boxrule=2.8pt,arc=13pt,
      drop shadow=popink,left=18pt,right=18pt,top=12pt,bottom=13pt,after skip=10pt]
    \poppill{#2 \textbullet\ #3}{popink}{popcream}\hspace{5pt}\poppill{#4}{white}{popink}\par\vspace{8pt}
    {\popdisplay\fontsize{14}{14}\selectfont\color{popink}CHAPTER #5}\par\vspace{4pt}
    {\raggedright\hyphenpenalty=10000\exhyphenpenalty=10000\popdisplay\fontsize{25}{27}\selectfont\color{popcream}\MakeUppercase{#1}\par}
    \vspace{8pt}
    {\popxbold\fontsize{9.5}{9.5}\selectfont\color{popink}\MakeUppercase{Suggested time: #6}}
  \end{tcolorbox}
  \begin{tcolorbox}[enhanced,colback=white,colframe=popink,boxrule=2.2pt,arc=10pt,
      drop shadow=popink,left=16pt,right=16pt,top=10pt,bottom=10pt,after skip=8pt]
    \poppill{In this chapter}{poppurple}{white}\par\vspace{5pt}
    \popsemi\fontsize{9.3}{12.6}\selectfont\color{popink}#7
  \end{tcolorbox}
  \noindent\begin{minipage}[t]{0.487\linewidth}
    \begin{tcolorbox}[enhanced,colback=popgreen,colframe=popink,boxrule=2.2pt,arc=10pt,
        drop shadow=popink,left=11pt,right=11pt,top=10pt,bottom=10pt,height=3.1cm,valign=center]
      \tcbox[colback=white,colframe=popink,boxrule=1.3pt,arc=5pt,boxsep=0pt,nobeforeafter,
        left=3pt,right=3pt,top=3pt,bottom=3pt,valign=center]{\qrcode[height=1.9cm]{https://play.google.com/store/apps/details?id=com.luvvix.onbuch&hl=en}}%
      \hspace{8pt}\begin{minipage}[c]{0.5\linewidth}
        {\popdisplay\fontsize{11}{12.5}\selectfont\color{white}\MakeUppercase{Download OnBuch}}\par\vspace{4pt}
        {\raggedright\popsemi\fontsize{8.4}{11}\selectfont\color{white}Free \textbullet\ Google Play\\AI tutor, past papers, results\par}
      \end{minipage}
    \end{tcolorbox}
  \end{minipage}\hfill
  \begin{minipage}[t]{0.487\linewidth}
    \begin{tcolorbox}[enhanced,colback=poppurple,colframe=popink,boxrule=2.2pt,arc=10pt,
        drop shadow=popink,left=11pt,right=11pt,top=10pt,bottom=10pt,height=3.1cm,valign=center]
      \tikz[baseline=-0.7ex]\node[circle,fill=white,draw=popink,line width=1.3pt,minimum size=1.6cm,inner sep=0]{\popdisplay\fontsize{24}{24}\selectfont\color{poppurple}+};%
      \hspace{8pt}\begin{minipage}[c]{0.6\linewidth}
        {\popdisplay\fontsize{11}{12.5}\selectfont\color{white}\MakeUppercase{Go OnBuch+}}\par\vspace{4pt}
        {\raggedright\popsemi\fontsize{8.4}{11}\selectfont\color{white}All signed courses, unlimited AI tutor, PDF sheets — OnBuch+ tab in the app\par}
      \end{minipage}
    \end{tcolorbox}
  \end{minipage}
  \vspace{8pt}
  \popcontenu
  \vfill
  \signatureonbuch
  \restoregeometry
  \newpage
}

% Rangée « Ce cours contient » (page de garde)
\newcommand{\poptile}[3]{% #1 couleur #2 gros texte #3 légende
  \begin{tcolorbox}[enhanced,colback=#1,colframe=popink,boxrule=2pt,arc=9pt,shadow={2.5pt}{-2.5pt}{0pt}{popink},
      left=6pt,right=6pt,top=9pt,bottom=9pt,halign=center,valign=center,height=1.75cm,width=\linewidth,nobeforeafter]
    {\popdisplay\fontsize{12.5}{13}\selectfont\color{white}\MakeUppercase{#2}}\par\vspace{3pt}
    {\centering\popsemi\fontsize{7.8}{9.5}\selectfont\color{white}#3\par}
  \end{tcolorbox}}
\newcommand{\popcontenu}{%
  \par\noindent{\popxboldls\fontsize{7.5}{7.5}\selectfont\color{popmuted}THIS SIGNED COURSE CONTAINS}\par\vspace{7pt}
  \noindent\begin{minipage}[t]{0.235\linewidth}\poptile{popblue}{Course}{complete, illustrated, step by step}\end{minipage}\hfill
  \begin{minipage}[t]{0.235\linewidth}\poptile{poporange}{Methods}{and worked examples}\end{minipage}\hfill
  \begin{minipage}[t]{0.235\linewidth}\poptile{poppink}{Exercises}{exam style + solutions}\end{minipage}\hfill
  \begin{minipage}[t]{0.235\linewidth}\poptile{popgreen}{Summary sheet}{the essentials to remember}\end{minipage}\par}

% Sommaire
\newtcolorbox{sommaire}{enhanced,colback=white,colframe=popink,boxrule=2.2pt,arc=10pt,
  drop shadow=popink,left=18pt,right=18pt,top=13pt,bottom=13pt,before skip=6pt,after skip=20pt,
  before upper={\poppill{Contents}{poppurple}{white}\par\vspace{9pt}\popsemi\fontsize{10.6}{14.5}\selectfont\color{popink}}}

% Signature de fin de cours — poussée tout en bas de la dernière page
\newcommand{\findecours}{%
  \par\vfill\nopagebreak
  \begin{center}\popsquiggle{poporange}\end{center}\vspace{4pt}
  \signatureonbuch
}
\AtBeginDocument{\def\llbracket{\mathopen{[\mkern-3.5mu[}}\def\rrbracket{\mathclose{]\mkern-3.5mu]}}} % intervalles d'entiers (glyphe ⟦ absent de la police)
% Glyphes absents de Fira Math : redéfinis à partir de glyphes disponibles
\AtBeginDocument{%
  \def\setminus{\mathbin{\backslash}}\let\smallsetminus\setminus
  \def\vdots{\mathord{\vbox{\baselineskip3.2pt\lineskiplimit0pt\kern4pt\hbox{.}\hbox{.}\hbox{.}}}}%
  \def\ddots{\mathinner{\mkern1mu\raise6.5pt\hbox{.}\mkern2mu\raise3.6pt\hbox{.}\mkern2mu\raise0.7pt\hbox{.}\mkern1mu}}%
  \def\triangle{\mathord{\scalebox{0.9}{$\symup{\Delta}$}}}%
  \def\nabla{\mathord{\raisebox{\depth}{\scalebox{1}[-1]{$\symup{\Delta}$}}}}%
  \def\checkmark{\popcheck{popdark}}%
  \def\star{\mathbin{\ast}}\def\bigstar{\mathord{\ast}}%
  \def\diamond{\mathbin{\scalebox{0.6}{$\Diamond$}}}%
  \def\bigcup{\mathop{\mathchoice{\raisebox{-0.3em}{\scalebox{1.7}{$\cup$}}}{\scalebox{1.25}{$\cup$}}{\cup}{\cup}}\displaylimits}%
  \def\bigcap{\mathop{\mathchoice{\raisebox{-0.3em}{\scalebox{1.7}{$\cap$}}}{\scalebox{1.25}{$\cap$}}{\cap}{\cap}}\displaylimits}%
  \def\lesseqqgtr{\mathrel{\lessgtr}}\def\bigoplus{\mathop{\oplus}\displaylimits}%
}
\providecommand{\ssi}{if and only if} % abréviation fréquente
\AtBeginDocument{\providecommand{\wideparen}{\overparen}} % arc de cercle

%% FIGFIX-BEGIN v4 — correctifs globaux des figures (docs/onbuch-generation/tools/figfix)
\usepackage{adjustbox}\usepackage{etoolbox}
% toute figure TikZ/pgfplots plus large que la ligne est réduite à la largeur disponible
\BeforeBeginEnvironment{tikzpicture}{\begin{adjustbox}{max width=\linewidth}}
\AfterEndEnvironment{tikzpicture}{\end{adjustbox}}
% légendes pgfplots lisibles par-dessus les courbes
\pgfplotsset{every axis/.append style={legend style={fill=white,fill opacity=0.92,text opacity=1,draw=black!15,inner sep=3pt}}}
% étiquettes d'arêtes (syntaxe edge["..."]) avec fond blanc
\tikzset{every edge quotes/.append style={fill=white,inner sep=1.2pt}}
%% FIGFIX-END
\begin{document}

\pagegarde{Further Complex Numbers}{Further Mathematics}{Upper Sixth}{Upper Sixth — Further mathematics}{2}{14 periods}{This chapter deepens the study of complex numbers: their forms, De Moivre's theorem and its applications, loci in the Argand diagram, and transformations of the complex plane. It culminates in similitudes, modulus inequalities and the classification of rigid motions by orientation.
\vspace{4pt}
\begin{multicols}{2}\small
\begin{itemize}
\item By the end of this chapter I can convert a complex number between Cartesian, trigonometric and exponential forms and use each form appropriately.
\item By the end of this chapter I can state and apply De Moivre's theorem to compute powers and nth roots of complex numbers.
\item By the end of this chapter I can use De Moivre's theorem to express cos(nθ), sin(nθ) in terms of cos θ and sin θ and to linearise powers of cosine and sine.
\item By the end of this chapter I can interpret and sketch loci defined by modulus and argument conditions in the Argand diagram.
\item By the end of this chapter I can describe elementary transformations from the z-plane to the w-plane and find images of points and curves.
\item By the end of this chapter I can represent reflections, rotations, enlargements, stretches and shears by 2×2 matrices and compose them.
\end{itemize}\end{multicols}}

\begin{sommaire}
\begin{multicols}{2}
\begin{enumerate}
\item Forms of a Complex Number and De Moivre's Theorem
\item Applications of De Moivre's Theorem
\item Loci in the Argand Diagram
\item Transformations from the z-plane to the w-plane
\item Plane Transformations and Matrix Representation
\item Similitudes: Complex Form, Matrix Interpretation and Applications
\item Modulus Inequalities and Rigid Motions
\item Integration activity
\item Methods and worked examples
\item Exercises
\item Solutions to the exercises
\item Summary sheet
\end{enumerate}
\end{multicols}
\end{sommaire}

\begin{objectifs}
\begin{itemize}
\item By the end of this chapter I can convert a complex number between Cartesian, trigonometric and exponential forms and use each form appropriately.
\item By the end of this chapter I can state and apply De Moivre's theorem to compute powers and nth roots of complex numbers.
\item By the end of this chapter I can use De Moivre's theorem to express cos(nθ), sin(nθ) in terms of cos θ and sin θ and to linearise powers of cosine and sine.
\item By the end of this chapter I can interpret and sketch loci defined by modulus and argument conditions in the Argand diagram.
\item By the end of this chapter I can describe elementary transformations from the z-plane to the w-plane and find images of points and curves.
\item By the end of this chapter I can represent reflections, rotations, enlargements, stretches and shears by 2×2 matrices and compose them.
\item By the end of this chapter I can identify a similitude, write its complex form z' = az + b or z' = a z̄ + b, and determine its characteristic elements.
\item By the end of this chapter I can solve and interpret modulus inequalities in the complex plane and apply them to geometric problems.
\item By the end of this chapter I can classify rigid motions as orientation-preserving (translations, rotations) or orientation-reversing (reflections, glide reflections).
\end{itemize}
\end{objectifs}

\begin{prerequis}
\begin{itemize}
\item Basic algebra of complex numbers from Lower Sixth: modulus, argument, conjugate, operations in Cartesian form
\item Trigonometric identities: addition formulae, double-angle formulae, solving trigonometric equations
\item Equations of lines and circles in the Cartesian plane; distance between two points
\item Basic 2×2 matrix operations: addition, multiplication, determinant, inverse
\item Solving polynomial equations, including factorisation and the quadratic formula with complex roots
\end{itemize}
\end{prerequis}

\newpage

\coursec{Forms of a Complex Number and De Moivre's Theorem}

When an engineer at the Douala seaport rotates a crane arm through an angle, or when MTN engineers analyse alternating-current signals across the Littoral region, they are secretly multiplying complex numbers. A single complex number can describe a position on a map of Yaoundé, a rotation of a satellite dish, and a stretch of a bridge cable --- all at once. How can one algebraic object encode turning, scaling and sliding on the plane? This chapter reveals how complex numbers, through De Moivre's theorem and the Argand diagram, become a complete language for plane geometry and transformations. Our first task is to master the three ways of writing a complex number --- Cartesian, trigonometric and exponential --- and to discover the astonishing rule that powers of complex numbers are computed by simply multiplying angles.

\courssub{Review of the Cartesian Form}

\begin{definition}[Complex number, Cartesian form]
A \cle{complex number} is an expression $z = x + iy$ where $x, y \in \mathbb{R}$ and $i$ is a number satisfying $i^2 = -1$. We call $x = \mathrm{Re}(z)$ the \cle{real part} and $y = \mathrm{Im}(z)$ the \cle{imaginary part} of $z$. The \cle{conjugate} of $z$ is $\bar{z} = x - iy$, and its \cle{modulus} is
\[ |z| = \sqrt{x^2 + y^2} = \sqrt{z\bar{z}}. \]
\end{definition}

Recall the essential algebraic facts, all proved by direct expansion: $\overline{z_1 + z_2} = \bar{z}_1 + \bar{z}_2$, $\overline{z_1 z_2} = \bar{z}_1 \bar{z}_2$, $z + \bar{z} = 2\,\mathrm{Re}(z)$, $z - \bar{z} = 2i\,\mathrm{Im}(z)$, and $z\bar{z} = |z|^2$. The modulus measures a \emph{distance}: geometrically, $|z|$ is the distance from the origin to the point representing $z$. Note also that $|\bar{z}| = |z|$ and that $|z| = 0$ if and only if $z = 0$.

\courssub{The Argand Diagram: Affixes and Geometry}

The plane equipped with an orthonormal frame $(O; \vec{u}, \vec{v})$ is called the \cle{Argand diagram} (or complex plane). To each complex number $z = x + iy$ we associate the point $M(x, y)$, called the \emph{image} of $z$; conversely $z$ is called the \cle{affix} of $M$, written $z_M$. Similarly, the vector $\vec{w} = a\vec{u} + b\vec{v}$ has affix $z_{\vec{w}} = a + ib$. This bijection between $\mathbb{C}$ and the plane is the key that unlocks all the geometry of this chapter.

Two geometric facts make the Argand diagram powerful:
\begin{itemize}
\item \textbf{Addition.} If $A$ and $B$ have affixes $z_A$ and $z_B$, then $z_A + z_B$ is the affix of $\vec{OA} + \vec{OB}$: addition of complex numbers is addition of vectors (parallelogram rule).
\item \textbf{Subtraction.} The vector $\vec{AB}$ has affix $z_B - z_A$, and the distance $AB = |z_B - z_A|$. The modulus of a difference is a \emph{distance between two points}.
\end{itemize}

\begin{popfigure}
\begin{center}
\begin{tikzpicture}[scale=1.6,>=stealth]
\draw[->,popmuted] (-0.6,0) -- (3.6,0) node[below left] {$\mathrm{Re}$};
\draw[->,popmuted] (0,-0.6) -- (0,2.8) node[below left] {$\mathrm{Im}$};
\coordinate (O) at (0,0);
\coordinate (M) at (2.6,1.8);
\draw[thick,popblue] (O) -- (M) node[above right,popdark] {$M(z=x+iy)$};
\draw[thick,poporange,->] (0.9,0) arc (0:34.7:0.9) node[midway,right] {$\theta$};
\draw[dashed,popmuted] (2.6,0) node[below] {$x$} -- (M);
\draw[dashed,popmuted] (0,1.8) node[left] {$y$} -- (M);
\draw[popmuted] (2.6,0.15) -- (2.45,0.15) -- (2.45,0);
\node[popblue,below] at (1.25,1.05) {$r=|z|$};
\fill[popdark] (M) circle (1.2pt);
\fill[popdark] (O) circle (1.2pt) node[below left] {$O$};
\end{tikzpicture}
\end{center}
\legende{The Argand diagram: the affix $z = x + iy$, its modulus $r = |z|$ and its argument $\theta = \arg(z)$ form a right triangle with $x = r\cos\theta$, $y = r\sin\theta$.}
\end{popfigure}

\courssub{Argument of a Non-zero Complex Number}

The modulus tells us \emph{how far} $M$ is from $O$; to locate $M$ completely we also need a \emph{direction}.

\begin{definition}[Modulus and principal argument]
Let $z \neq 0$ with image $M$. An \cle{argument} of $z$ is any measure $\theta$ of the oriented angle $(\vec{u}, \vec{OM})$. The unique argument lying in $(-\pi, \pi]$ is called the \cle{principal argument}, denoted $\mathrm{Arg}(z)$. Every argument has the form
\[ \arg(z) = \mathrm{Arg}(z) + 2k\pi, \qquad k \in \mathbb{Z}. \]
\end{definition}

If $z = x + iy$ with $|z| = r$, then $\theta = \mathrm{Arg}(z)$ satisfies
\[ \cos\theta = \frac{x}{r}, \qquad \sin\theta = \frac{y}{r}, \qquad \tan\theta = \frac{y}{x}\ (x \neq 0). \]
The pair $(\cos\theta, \sin\theta)$ fixes $\theta$ uniquely in $(-\pi, \pi]$; the tangent alone does not, since $\tan\theta = \tan(\theta + \pi)$. Always use the signs of $x$ and $y$ to choose the correct quadrant. In particular, a positive real number has argument $0$, a negative real number has principal argument $\pi$, and a pure imaginary $iy$ has principal argument $\frac{\pi}{2}$ if $y > 0$ and $-\frac{\pi}{2}$ if $y < 0$.

\begin{attention}[The argument is defined only up to $2k\pi$]
A very common mistake is to write ``$\arg(z) = \frac{\pi}{3}$'' as if the argument were a single number. In fact $\arg(z)$ is a \emph{family} of values differing by multiples of $2\pi$. In every solution, state clearly which determination you use: either the principal argument $\mathrm{Arg}(z) \in (-\pi, \pi]$, or write $\arg(z) \equiv \theta \pmod{2\pi}$. Also remember: the number $0$ has \textbf{no argument} --- the direction of the zero vector is undefined.
\end{attention}

\courssub{Trigonometric and Exponential Forms}

Putting modulus and argument together gives the second great way of writing a complex number.

\begin{definition}[Trigonometric (polar) form]
Every non-zero complex number can be written
\[ z = r(\cos\theta + i\sin\theta), \qquad r = |z| > 0,\ \theta = \arg(z). \]
This is the \cle{trigonometric form} (or polar form) of $z$.
\end{definition}

\begin{methode}[Converting from Cartesian to trigonometric form]
Given $z = x + iy \neq 0$:
\begin{enumerate}
\item Compute the modulus $r = \sqrt{x^2 + y^2}$.
\item Factorise: $z = r\left(\dfrac{x}{r} + i\dfrac{y}{r}\right)$.
\item Find $\theta$ such that $\cos\theta = \dfrac{x}{r}$ and $\sin\theta = \dfrac{y}{r}$, using the signs of $x$ and $y$ to fix the quadrant.
\item Conclude $z = r(\cos\theta + i\sin\theta)$, stating $\mathrm{Arg}(z) = \theta$.
\end{enumerate}
\end{methode}

\begin{exemplebox}[From Cartesian to polar form]
Convert $z = -\sqrt{3} + i$ to trigonometric form. We have $r = \sqrt{3 + 1} = 2$, so
\[ z = 2\left(-\frac{\sqrt{3}}{2} + i\,\frac{1}{2}\right). \]
We need $\cos\theta = -\frac{\sqrt{3}}{2}$ and $\sin\theta = \frac{1}{2}$: this is $\theta = \frac{5\pi}{6}$ (second quadrant, since $x<0$, $y>0$). Hence
\[ z = 2\left(\cos\frac{5\pi}{6} + i\sin\frac{5\pi}{6}\right), \qquad \mathrm{Arg}(z) = \frac{5\pi}{6}. \]
\end{exemplebox}

Euler's brilliant idea was to compress $\cos\theta + i\sin\theta$ into a genuine exponential. One defines, for real $\theta$,
\[ e^{i\theta} = \cos\theta + i\sin\theta, \]
so that every non-zero complex number has the \cle{exponential form} $z = re^{i\theta}$. This definition is consistent because $e^{i\theta}$ obeys the usual exponential law $e^{i\theta_1}e^{i\theta_2} = e^{i(\theta_1+\theta_2)}$ --- which is exactly the addition formula for sines and cosines, as the proof of the product rule below shows. From $e^{i\theta} = \cos\theta + i\sin\theta$ and $e^{-i\theta} = \cos\theta - i\sin\theta$, adding and subtracting gives \cle{Euler's formulae}:
\[ \cos\theta = \frac{e^{i\theta} + e^{-i\theta}}{2}, \qquad \sin\theta = \frac{e^{i\theta} - e^{-i\theta}}{2i}. \]
These will be indispensable when we linearise trigonometric powers in the next section. Note also the useful special values $e^{i\pi} = -1$, $e^{i\pi/2} = i$ and $e^{-i\pi/2} = -i$.

\courssub{Multiplication and Division in Polar Form}

Here is the heart of the matter: polar form turns multiplication into something geometrically transparent.

\begin{propriete}[Modulus and argument of a product and quotient]
For non-zero complex numbers $z_1 = r_1 e^{i\theta_1}$ and $z_2 = r_2 e^{i\theta_2}$:
\[ |z_1 z_2| = |z_1|\,|z_2|, \qquad \arg(z_1 z_2) \equiv \arg(z_1) + \arg(z_2) \pmod{2\pi}, \]
\[ \left|\frac{z_1}{z_2}\right| = \frac{|z_1|}{|z_2|}, \qquad \arg\!\left(\frac{z_1}{z_2}\right) \equiv \arg(z_1) - \arg(z_2) \pmod{2\pi}. \]
\emph{Proof (examinable).} Write $z_1 z_2 = r_1 r_2(\cos\theta_1 + i\sin\theta_1)(\cos\theta_2 + i\sin\theta_2)$. Expanding and using the addition formulae $\cos(\theta_1+\theta_2) = \cos\theta_1\cos\theta_2 - \sin\theta_1\sin\theta_2$ and $\sin(\theta_1+\theta_2) = \sin\theta_1\cos\theta_2 + \cos\theta_1\sin\theta_2$, we obtain
\[ z_1 z_2 = r_1 r_2\bigl(\cos(\theta_1+\theta_2) + i\sin(\theta_1+\theta_2)\bigr), \]
which is the trigonometric form of $z_1z_2$ with modulus $r_1 r_2$ and argument $\theta_1 + \theta_2$. For the quotient, apply the product rule to $z_1 = \dfrac{z_1}{z_2}\cdot z_2$: this gives $|z_1| = \left|\dfrac{z_1}{z_2}\right||z_2|$ and $\arg(z_1) \equiv \arg\!\left(\dfrac{z_1}{z_2}\right) + \arg(z_2) \pmod{2\pi}$, from which both claims follow. \hfill$\blacksquare$
\end{propriete}

Geometrically, multiplying by $z_2 = r_2 e^{i\theta_2}$ \emph{rotates} the plane about $O$ through the angle $\theta_2$ and \emph{scales} distances by the factor $r_2$. In particular, multiplication by $e^{i\theta}$ (modulus $1$) is a pure rotation of angle $\theta$ about the origin --- this is exactly the crane arm of our opening hook. Since $\bar{z} = re^{-i\theta}$, we also read off $\arg(\bar{z}) \equiv -\arg(z) \pmod{2\pi}$ and $\arg\!\left(\dfrac{1}{z}\right) \equiv -\arg(z) \pmod{2\pi}$.

\begin{popfigure}
\begin{center}
\begin{tikzpicture}[scale=1.7,>=stealth]
\draw[->,popmuted] (-0.5,0) -- (3.4,0) node[below left] {$\mathrm{Re}$};
\draw[->,popmuted] (0,-0.5) -- (0,3.0) node[below left] {$\mathrm{Im}$};
\coordinate (O) at (0,0);
\coordinate (Z) at (2.5,0.9);
\coordinate (W) at (0.9,2.5);
\draw[thick,popblue,->] (O) -- (Z) node[midway,below right] {$z$};
\draw[thick,poporange,->] (O) -- (W) node[midway,above left] {$w = e^{i\theta}z$};
\draw[thick,popgreen,->] (1.1,0.4) arc (20:70:1.17) node[midway,above right] {$\theta$};
\draw[dashed,popmuted] (O) circle (2.64);
\fill[popdark] (Z) circle (1.2pt);
\fill[popdark] (W) circle (1.2pt);
\fill[popdark] (O) circle (1.2pt) node[below left] {$O$};
\node[popmuted] at (2.9,2.2) {$|w|=|z|$};
\end{tikzpicture}
\end{center}
\legende{Multiplication by $e^{i\theta}$ rotates the image of $z$ about $O$ through the angle $\theta$; the modulus is preserved, so the point moves along its circle centred at $O$.}
\end{popfigure}

\courssub{De Moivre's Theorem}

What happens when we multiply a complex number by itself repeatedly? The product rule says the modulus is raised to a power and the arguments \emph{accumulate}. This observation becomes one of the most useful theorems of the whole course.

\begin{propriete}[De Moivre's theorem --- proof examinable]
For every real number $\theta$ and every integer $n \in \mathbb{Z}$,
\[ \bigl(\cos\theta + i\sin\theta\bigr)^n = \cos(n\theta) + i\sin(n\theta), \]
equivalently $\bigl(e^{i\theta}\bigr)^n = e^{in\theta}$.
\end{propriete}

\textbf{Proof.} \emph{Case $n \geq 0$: induction.} Let $P(n)$ be the statement $(\cos\theta + i\sin\theta)^n = \cos(n\theta) + i\sin(n\theta)$.

\emph{Base case.} For $n = 0$: $(\cos\theta + i\sin\theta)^0 = 1 = \cos 0 + i\sin 0$, so $P(0)$ holds.

\emph{Inductive step.} Assume $P(n)$ true for some $n \geq 0$. Then
\begin{align*}
(\cos\theta + i\sin\theta)^{n+1}
&= (\cos\theta + i\sin\theta)^n (\cos\theta + i\sin\theta) \\
&= \bigl(\cos(n\theta) + i\sin(n\theta)\bigr)(\cos\theta + i\sin\theta) \\
&= \cos(n\theta + \theta) + i\sin(n\theta + \theta) \quad \text{(product rule)} \\
&= \cos\bigl((n+1)\theta\bigr) + i\sin\bigl((n+1)\theta\bigr),
\end{align*}
so $P(n+1)$ holds. By induction, $P(n)$ is true for all $n \in \mathbb{N}$.

\emph{Case $n < 0$.} Write $n = -m$ with $m \in \mathbb{N}^*$. Since $\cos\theta + i\sin\theta \neq 0$ (its modulus is $1$),
\[ (\cos\theta + i\sin\theta)^{-m} = \frac{1}{(\cos\theta + i\sin\theta)^m} = \frac{1}{\cos(m\theta) + i\sin(m\theta)}. \]
Multiplying numerator and denominator by the conjugate $\cos(m\theta) - i\sin(m\theta)$, whose product with the denominator is $\cos^2(m\theta) + \sin^2(m\theta) = 1$, gives
\[ (\cos\theta + i\sin\theta)^{-m} = \cos(m\theta) - i\sin(m\theta) = \cos(-m\theta) + i\sin(-m\theta) = \cos(n\theta) + i\sin(n\theta), \]
using the parity of cosine and sine. Hence the theorem holds for all $n \in \mathbb{Z}$. \hfill$\blacksquare$

\begin{attention}[De Moivre is stated for integer exponents]
The identity $(\cos\theta + i\sin\theta)^n = \cos(n\theta) + i\sin(n\theta)$ is proved here for $n \in \mathbb{Z}$. For \emph{fractional} exponents the formula becomes ambiguous: a complex number has $n$ distinct $n$-th roots, so $\bigl(\cos\theta + i\sin\theta\bigr)^{1/n}$ takes $n$ different values, not one. Writing $\cos(\theta/n) + i\sin(\theta/n)$ gives only \emph{one} of them. We shall exploit this properly when studying roots of unity; for now, never apply De Moivre blindly to fractions.
\end{attention}

\courssub{Computing Powers Efficiently}

De Moivre's theorem converts the painful task of expanding $(x + iy)^n$ into a one-line computation: raise the modulus to the power $n$ and multiply the argument by $n$. Explicitly, if $z = re^{i\theta}$ then
\[ z^n = r^n e^{in\theta} = r^n\bigl(\cos(n\theta) + i\sin(n\theta)\bigr). \]

\begin{exoresolu}[Powers via polar form]
Compute $(1 + i)^{12}$ and $(\sqrt{3} + i)^9$, giving answers in Cartesian form.
\tcblower
\textbf{For $1 + i$:} $r = \sqrt{2}$, and $\cos\theta = \sin\theta = \frac{1}{\sqrt{2}}$ gives $\mathrm{Arg}(1+i) = \frac{\pi}{4}$. Hence $1 + i = \sqrt{2}\,e^{i\pi/4}$ and by De Moivre,
\[ (1+i)^{12} = (\sqrt{2})^{12}\,e^{i \cdot 12\pi/4} = 2^6 e^{3i\pi} = 64\bigl(\cos 3\pi + i\sin 3\pi\bigr) = -64, \]
since $\cos 3\pi = \cos\pi = -1$ and $\sin 3\pi = 0$.

\textbf{For $\sqrt{3} + i$:} $r = \sqrt{3+1} = 2$, $\cos\theta = \frac{\sqrt{3}}{2}$, $\sin\theta = \frac{1}{2}$, so $\mathrm{Arg} = \frac{\pi}{6}$. Thus
\[ (\sqrt{3}+i)^9 = 2^9 e^{i \cdot 9\pi/6} = 512\,e^{3i\pi/2} = 512\left(\cos\frac{3\pi}{2} + i\sin\frac{3\pi}{2}\right) = -512\,i. \]
Notice how reducing the final angle modulo $2\pi$ (here $3\pi \equiv \pi$, while $\frac{3\pi}{2}$ is already principal) makes the Cartesian answer immediate.
\end{exoresolu}

\begin{exemplebox}[A negative power]
Compute $(1 + i)^{-8}$. With $1 + i = \sqrt{2}\,e^{i\pi/4}$:
\[ (1+i)^{-8} = (\sqrt{2})^{-8} e^{-2i\pi} = \frac{1}{16}\bigl(\cos(-2\pi) + i\sin(-2\pi)\bigr) = \frac{1}{16}. \]
\end{exemplebox}

\begin{aretenir}[Key points of this section]
\begin{itemize}
\item Three forms of $z \neq 0$: Cartesian $x + iy$, trigonometric $r(\cos\theta + i\sin\theta)$, exponential $re^{i\theta}$, with $r = |z|$, $\theta = \arg(z)$.
\item The principal argument $\mathrm{Arg}(z) \in (-\pi, \pi]$ is unique; general arguments differ by $2k\pi$. The number $0$ has no argument.
\item Moduli multiply, arguments add: $|z_1z_2| = |z_1||z_2|$, $\arg(z_1z_2) \equiv \arg(z_1) + \arg(z_2) \pmod{2\pi}$. Multiplication by $e^{i\theta}$ is a rotation of angle $\theta$ about $O$.
\item Euler's formulae: $\cos\theta = \frac{e^{i\theta}+e^{-i\theta}}{2}$, $\sin\theta = \frac{e^{i\theta}-e^{-i\theta}}{2i}$.
\item De Moivre: $(\cos\theta + i\sin\theta)^n = \cos(n\theta) + i\sin(n\theta)$ for all $n \in \mathbb{Z}$ --- proved by induction for $n \geq 0$, extended to negatives via the conjugate.
\item To compute $z^n$: convert to polar form, then $z^n = r^n\bigl(\cos(n\theta) + i\sin(n\theta)\bigr)$.
\end{itemize}
\end{aretenir}

\begin{exercice}[Practice]
\begin{enumerate}
\item Write in trigonometric and exponential form: (a) $-1 + i\sqrt{3}$; (b) $-4i$; (c) $3 - 3i$.
\item Given $z_1 = 2e^{i\pi/6}$ and $z_2 = 3e^{-i\pi/4}$, find the modulus and principal argument of $z_1 z_2$ and of $\dfrac{z_1}{z_2}$.
\item Compute $(1 - i)^{10}$ and $(\sqrt{3} - i)^{-6}$ in Cartesian form.
\item Show that $(1 + i)^n + (1 - i)^n$ is real for every $n \in \mathbb{N}$, and find its value in terms of $n$.
\end{enumerate}
\end{exercice}

\begin{corrige}[Exercise 1]
\textbf{1.} (a) $r = 2$, $\cos\theta = -\frac12$, $\sin\theta = \frac{\sqrt3}{2}$, so $-1 + i\sqrt3 = 2e^{2i\pi/3}$. (b) $-4i = 4e^{-i\pi/2}$. (c) $r = 3\sqrt2$, $\theta = -\frac{\pi}{4}$, so $3 - 3i = 3\sqrt2\,e^{-i\pi/4}$.

\textbf{2.} $z_1z_2 = 6e^{-i\pi/12}$: modulus $6$, $\mathrm{Arg} = -\frac{\pi}{12}$. $\dfrac{z_1}{z_2} = \frac{2}{3}e^{5i\pi/12}$: modulus $\frac23$, $\mathrm{Arg} = \frac{5\pi}{12}$.

\textbf{3.} $1 - i = \sqrt2\,e^{-i\pi/4}$, so $(1-i)^{10} = 2^5 e^{-10i\pi/4} = 32e^{-i\pi/2} = -32i$. Also $\sqrt3 - i = 2e^{-i\pi/6}$, so $(\sqrt3 - i)^{-6} = 2^{-6}e^{i\pi} = -\frac{1}{64}$.

\textbf{4.} $1 + i = \sqrt2\,e^{i\pi/4}$ and $1 - i = \sqrt2\,e^{-i\pi/4}$. By De Moivre, $(1+i)^n + (1-i)^n = 2^{n/2}\bigl(e^{in\pi/4} + e^{-in\pi/4}\bigr) = 2^{n/2}\cdot 2\cos\frac{n\pi}{4} = 2^{(n+2)/2}\cos\frac{n\pi}{4} \in \mathbb{R}$ by Euler's formula.
\end{corrige}

\begin{savaistu}[Abraham de Moivre]
Abraham de Moivre (1667--1754) was a French mathematician who fled to England and became a friend of Isaac Newton. His famous formula, published in the early eighteenth century, later became the bridge that led Euler to the identity $e^{i\pi} + 1 = 0$, often called the most beautiful equation in mathematics because it links the five fundamental constants $0, 1, i, e, \pi$. Today, the same mathematics you have just learnt drives the signal processing inside every mobile phone transmitting across Douala and Yaoundé.
\end{savaistu}

\coursec{Applications of De Moivre's Theorem}

In the previous section we established De Moivre's theorem, $(\cos\theta+i\sin\theta)^n=\cos n\theta+i\sin n\theta$, and its exponential form $(re^{i\theta})^n=r^ne^{in\theta}$. That single identity is a remarkably powerful machine. In this section we put it to work in four directions: extracting \emph{all} the $n$th roots of a complex number, expanding $\cos n\theta$ and $\sin n\theta$ as polynomials, \emph{linearising} powers $\cos^n\theta$ and $\sin^n\theta$ (a gift for integration), and summing trigonometric series. These are among the most frequently examined techniques in the GCE Advanced Level Further Mathematics paper.

\courssub{The $n$th Roots of a Complex Number}

\textbf{Motivation.} In $\mathbb{R}$, the equation $x^2=9$ has two solutions, $x=\pm3$, and $x^3=8$ has only one, $x=2$. In $\mathbb{C}$ the picture is far richer: the fundamental theorem of algebra guarantees that $z^n=w$ (with $w\neq0$) has \emph{exactly $n$ solutions}. De Moivre's theorem lets us find them all explicitly.

\begin{definition}[$n$th roots of a complex number]
Let $w\in\mathbb{C}$, $w\neq0$, and let $n\in\mathbb{N}$, $n\geq2$. An \cle{$n$th root} of $w$ is any complex number $z$ such that $z^n=w$. A solution of $z^n=1$ is called an \cle{$n$th root of unity}.
\end{definition}

\begin{propriete}[Formula for the $n$th roots]
Let $w=re^{i\theta}$ with $r>0$. Then $z^n=w$ has exactly $n$ distinct solutions, given by
\[
z_k=r^{1/n}\,e^{\,i(\theta+2k\pi)/n},\qquad k=0,1,2,\dots,n-1.
\]
\end{propriete}

\textbf{Derivation (examinable).} Write the unknown $z$ in polar form, $z=\rho e^{i\varphi}$ with $\rho>0$. By De Moivre, $z^n=\rho^n e^{in\varphi}$. The equation $z^n=re^{i\theta}$ then forces
\[
\rho^n=r\quad\text{and}\quad n\varphi=\theta+2k\pi,\ k\in\mathbb{Z},
\]
because two complex numbers in polar form are equal exactly when their moduli are equal and their arguments differ by a multiple of $2\pi$. Since $\rho>0$ is real, $\rho=r^{1/n}$ (the ordinary positive real $n$th root), and $\varphi=\dfrac{\theta+2k\pi}{n}$. Taking $k=0,1,\dots,n-1$ gives $n$ distinct values of $\varphi$ in an interval of length $2\pi$; any other integer $k$ merely repeats one of these (adding $2\pi$ to an argument changes nothing). Hence there are exactly $n$ distinct roots. $\blacksquare$

\begin{propriete}[Geometric distribution of the $n$th roots]
The $n$ roots of $z^n=w$ all have the same modulus $r^{1/n}$ and their arguments differ successively by $\dfrac{2\pi}{n}$. Hence their images in the Argand diagram are the vertices of a \cle{regular $n$-gon} inscribed in the circle of centre $O$ and radius $r^{1/n}$.
\end{propriete}

\begin{attention}[There are exactly $n$ roots — not one!]
A very common error is to give only the ``obvious'' root, for instance writing $z^3=8\Rightarrow z=2$ and stopping. In $\mathbb{C}$, $z^3=8$ has \textbf{three} solutions: $2$, $2e^{2i\pi/3}=-1+i\sqrt3$ and $2e^{-2i\pi/3}=-1-i\sqrt3$. Always count your answers: for $z^n=w$ you must list $n$ distinct roots.
\end{attention}

\begin{methode}[Solving $z^n=w$ step by step]
\begin{enumerate}
\item Write $w$ in polar form: $w=re^{i\theta}$ (find $r=|w|$ and an argument $\theta$).
\item The common modulus of all roots is $r^{1/n}$.
\item The arguments are $\dfrac{\theta+2k\pi}{n}$ for $k=0,1,\dots,n-1$.
\item List the $n$ roots $z_k=r^{1/n}e^{i(\theta+2k\pi)/n}$, and if required convert each to Cartesian form $a+ib$.
\item Check: you must have exactly $n$ roots, equally spaced by $\dfrac{2\pi}{n}$ around the circle of radius $r^{1/n}$.
\end{enumerate}
\end{methode}

\begin{exoresolu}[Solving $z^4=-16$]
Solve $z^4=-16$ in $\mathbb{C}$, giving the roots in Cartesian form.
\tcblower
\textbf{Solution.} Write $-16$ in polar form: $|-16|=16$ and $\arg(-16)=\pi$, so $-16=16e^{i\pi}$. The common modulus of the roots is $16^{1/4}=2$, and the arguments are
\[
\varphi_k=\frac{\pi+2k\pi}{4},\qquad k=0,1,2,3,
\]
giving $\varphi_0=\dfrac{\pi}{4}$, $\varphi_1=\dfrac{3\pi}{4}$, $\varphi_2=\dfrac{5\pi}{4}$, $\varphi_3=\dfrac{7\pi}{4}$. Hence
\begin{align*}
z_0&=2e^{i\pi/4}=\sqrt2+i\sqrt2, &
z_1&=2e^{3i\pi/4}=-\sqrt2+i\sqrt2,\\
z_2&=2e^{5i\pi/4}=-\sqrt2-i\sqrt2, &
z_3&=2e^{7i\pi/4}=\sqrt2-i\sqrt2.
\end{align*}
The four roots are the vertices of a square inscribed in the circle of radius $2$.
\end{exoresolu}

\begin{exoresolu}[Solving $z^3=8i$]
Solve $z^3=8i$ in $\mathbb{C}$.
\tcblower
\textbf{Solution.} We have $|8i|=8$ and $\arg(8i)=\dfrac{\pi}{2}$, so $8i=8e^{i\pi/2}$. The cube roots have modulus $8^{1/3}=2$ and arguments
\[
\varphi_k=\frac{\frac{\pi}{2}+2k\pi}{3}=\frac{\pi}{6}+\frac{2k\pi}{3},\qquad k=0,1,2,
\]
that is $\dfrac{\pi}{6}$, $\dfrac{5\pi}{6}$, $\dfrac{3\pi}{2}$. Therefore
\[
z_0=2e^{i\pi/6}=\sqrt3+i,\qquad
z_1=2e^{5i\pi/6}=-\sqrt3+i,\qquad
z_2=2e^{3i\pi/2}=-2i.
\]
The three roots are spaced at $120^\circ$ on the circle of radius $2$, as shown below.
\end{exoresolu}

\begin{popfigure}
\begin{center}
\begin{tikzpicture}[scale=1.5]
\draw[popline,->] (-2.6,0)--(2.8,0) node[below left]{\small Re};
\draw[popline,->] (0,-2.6)--(0,2.8) node[below left]{\small Im};
\draw[popblue,thick] (0,0) circle (2);
\foreach \a/\name/\pos in {30/{z_0=\sqrt3+i}/{above right},150/{z_1=-\sqrt3+i}/{above left},270/{z_2=-2i}/{below}}{
  \fill[poppurple] (\a:2) circle (1.6pt);
  \node[\pos,popdark] at (\a:2.15) {\small $\name$};
  \draw[popmuted,dashed] (0,0)--(\a:2);
}
\draw[poporange,->] (0.9,0) arc (0:30:0.9) node[midway,right]{\small $\frac{\pi}{6}$};
\draw[popgreen,->] (30:1.15) arc (30:150:1.15) node[midway,above]{\small $\frac{2\pi}{3}$};
\fill (0,0) circle (1.2pt) node[below left]{\small $O$};
\end{tikzpicture}
\end{center}
\legende{The three cube roots of $8i$: equally spaced at $120^\circ$ on the circle of radius $2$.}
\end{popfigure}

\begin{attention}[Degrees and radians: do not mix them]
When listing the arguments $\dfrac{\theta+2k\pi}{n}$, work \textbf{entirely in radians} (or entirely in degrees, but radians are the GCE standard). A classic blunder is to write $\theta=\frac{\pi}{2}$ and then compute $\frac{\pi/2+360k}{3}$. Mixing the two systems produces wrong angles and wrong roots. Convert first, then compute.
\end{attention}

\courssub{Roots of Unity}

The equation $z^n=1$ is the special case $w=1=1\cdot e^{i\cdot0}$. Its solutions, the \cle{$n$th roots of unity}, are
\[
z_k=e^{2ik\pi/n},\qquad k=0,1,\dots,n-1,
\]
the vertices of a regular $n$-gon inscribed in the \emph{unit circle}, one vertex being $z_0=1$.

\begin{popfigure}
\begin{center}
\begin{tikzpicture}[scale=2.2]
\draw[popline,->] (-1.5,0)--(1.6,0) node[below left]{\small Re};
\draw[popline,->] (0,-1.4)--(0,1.5) node[below left]{\small Im};
\draw[popblue,thick] (0,0) circle (1);
\draw[poppink,thick] (0:1)--(72:1)--(144:1)--(216:1)--(288:1)--cycle;
\foreach \i/\pos in {0/{right},1/{above right},2/{above left},3/{below left},4/{below right}}{
  \pgfmathsetmacro{\angle}{\i*72}
  \fill[poporange] (\angle:1) circle (1.4pt);
  \node[\pos,popdark] at (\angle:1.14) {\small $z_{\i}$};
  \draw[popmuted,dashed] (0,0)--(\angle:1);
}
\draw[popgreen,->] (0:0.55) arc (0:72:0.55) node[midway,right]{\small $\frac{2\pi}{5}$};
\fill (0,0) circle (1pt) node[below left]{\small $O$};
\end{tikzpicture}
\end{center}
\legende{The five fifth roots of unity: vertices of a regular pentagon on the unit circle.}
\end{popfigure}

\begin{propriete}[Sum and product of the $n$th roots of unity]
Let $z_0,z_1,\dots,z_{n-1}$ be the $n$th roots of unity, with $n\geq2$. Then
\[
z_0+z_1+\cdots+z_{n-1}=0
\qquad\text{and}\qquad
z_0z_1\cdots z_{n-1}=(-1)^{n-1}.
\]
\end{propriete}

\textbf{Justification.} The roots are the solutions of $z^n-1=0$. In a polynomial $z^n+0\cdot z^{n-1}+\cdots-1$, the sum of the roots is minus the coefficient of $z^{n-1}$, which is $0$; the product is $(-1)^n\times(\text{constant term})=(-1)^n(-1)=(-1)^{n-1}$. Alternatively, writing $\omega=e^{2i\pi/n}$, the roots are $1,\omega,\omega^2,\dots,\omega^{n-1}$, a geometric progression whose sum is $\dfrac{\omega^n-1}{\omega-1}=0$ since $\omega^n=1$. $\blacksquare$

\begin{aretenir}[Key facts about roots of unity]
\begin{itemize}
\item Setting $\omega=e^{2i\pi/n}$, the roots are $1,\omega,\omega^2,\dots,\omega^{n-1}$ and $\omega^n=1$.
\item Their sum is $0$ (for $n\geq2$); geometrically, the polygon is ``balanced'' about $O$.
\item Non-real roots occur in conjugate pairs, since the coefficients of $z^n-1$ are real.
\end{itemize}
\end{aretenir}

\courssub{Expanding $\cos n\theta$ and $\sin n\theta$}

De Moivre gives $\cos n\theta+i\sin n\theta=(\cos\theta+i\sin\theta)^n$. Expanding the right-hand side with the \textbf{binomial theorem} and then equating real and imaginary parts produces polynomial identities.

\begin{exemplebox}[The triple-angle formulae]
With $c=\cos\theta$, $s=\sin\theta$:
\[
\cos3\theta+i\sin3\theta=(c+is)^3=c^3+3c^2(is)+3c(is)^2+(is)^3
=(c^3-3cs^2)+i(3c^2s-s^3).
\]
Equating real and imaginary parts:
\[
\cos3\theta=\cos^3\theta-3\cos\theta\sin^2\theta=4\cos^3\theta-3\cos\theta,
\]
\[
\sin3\theta=3\cos^2\theta\sin\theta-\sin^3\theta=3\sin\theta-4\sin^3\theta,
\]
where we used $\sin^2\theta=1-\cos^2\theta$ and $\cos^2\theta=1-\sin^2\theta$.
\end{exemplebox}

\begin{methode}[Expanding $\cos n\theta$ and $\sin n\theta$]
\begin{enumerate}
\item Write $\cos n\theta+i\sin n\theta=(c+is)^n$ (De Moivre).
\item Expand $(c+is)^n$ by the binomial theorem; use $i^2=-1$, $i^3=-i$, $i^4=1$, \dots
\item Equate real parts to get $\cos n\theta$, imaginary parts to get $\sin n\theta$.
\item If a polynomial in $\cos\theta$ alone (or $\sin\theta$ alone) is wanted, eliminate the other function via $c^2+s^2=1$.
\end{enumerate}
\end{methode}

\textbf{Expressing $\tan n\theta$ in terms of $t=\tan\theta$.} Divide the imaginary part by the real part and then divide numerator and denominator by $\cos^n\theta$. For example, with $t=\tan\theta$:
\[
\tan3\theta=\frac{3c^2s-s^3}{c^3-3cs^2}=\frac{3t-t^3}{1-3t^2}.
\]
This is the pattern of the classical ``$\tan n\theta$'' formulae, and the same technique works for any $n$.

\courssub{Linearisation of $\cos^n\theta$ and $\sin^n\theta$}

The reverse problem is just as important, especially for integration: express $\cos^n\theta$ or $\sin^n\theta$ as a \emph{sum} of cosines or sines of multiple angles. The tool is now Euler's formulae.

\begin{methode}[Linearising $\cos^n\theta$ using Euler's formulae]
\begin{enumerate}
\item Write $\cos\theta=\dfrac{e^{i\theta}+e^{-i\theta}}{2}$, so $\cos^n\theta=\dfrac{1}{2^n}\left(e^{i\theta}+e^{-i\theta}\right)^n$.
\item Expand with the binomial theorem.
\item Group symmetric terms $e^{ik\theta}+e^{-ik\theta}=2\cos k\theta$.
\item Simplify to obtain a sum of cosines of multiple angles. (For $\sin^n\theta$, start from $\sin\theta=\dfrac{e^{i\theta}-e^{-i\theta}}{2i}$; sines of multiple angles appear when $n$ is odd.)
\end{enumerate}
\end{methode}

\begin{exoresolu}[Linearising $\cos^4\theta$]
Express $\cos^4\theta$ as a sum of cosines of multiple angles. Hence evaluate $\displaystyle\int_0^{\pi/2}\cos^4\theta\,d\theta$.
\tcblower
\textbf{Solution.} By Euler's formula and the binomial theorem,
\begin{align*}
\cos^4\theta&=\frac{1}{16}\left(e^{i\theta}+e^{-i\theta}\right)^4\\
&=\frac{1}{16}\left(e^{4i\theta}+4e^{2i\theta}+6+4e^{-2i\theta}+e^{-4i\theta}\right)\\
&=\frac{1}{16}\bigl(2\cos4\theta+8\cos2\theta+6\bigr)
=\frac{1}{8}\cos4\theta+\frac{1}{2}\cos2\theta+\frac{3}{8}.
\end{align*}
Then
\[
\int_0^{\pi/2}\cos^4\theta\,d\theta
=\left[\frac{1}{32}\sin4\theta+\frac{1}{4}\sin2\theta+\frac{3\theta}{8}\right]_0^{\pi/2}
=0+0+\frac{3\pi}{16}=\frac{3\pi}{16}.
\]
Without linearisation this integral would be painful; with it, it is immediate.
\end{exoresolu}

\begin{savaistu}[Why examiners love linearisation]
Integrals such as $\displaystyle\int\cos^6x\,dx$ or $\displaystyle\int\sin^3x\cos^2x\,dx$ appear regularly in GCE Further Mathematics and in first-year university courses (including at the University of Buea and Yaound\'e I). Linearisation via Euler's formulae converts them into sums of easy integrals $\int\cos kx\,dx$. It is one of the highest-value techniques of this chapter.
\end{savaistu}

\courssub{Summing Trigonometric Series}

A final beautiful application: sums such as
\[
C=\sum_{k=0}^{n-1}\cos(\theta+k\alpha),\qquad S=\sum_{k=0}^{n-1}\sin(\theta+k\alpha)
\]
are handled by recognising them as the real and imaginary parts of a \textbf{geometric series} of complex exponentials.

\begin{exemplebox}[The standard summation]
Consider $C+iS=\displaystyle\sum_{k=0}^{n-1}e^{i(\theta+k\alpha)}=e^{i\theta}\sum_{k=0}^{n-1}\left(e^{i\alpha}\right)^k$.
For $\alpha\not\equiv0\pmod{2\pi}$, the geometric sum gives
\[
C+iS=e^{i\theta}\cdot\frac{e^{in\alpha}-1}{e^{i\alpha}-1}.
\]
Using the half-angle trick, $e^{ix}-1=e^{ix/2}\cdot2i\sin\frac{x}{2}$:
\[
C+iS=e^{i\theta}\cdot\frac{e^{in\alpha/2}\cdot2i\sin\frac{n\alpha}{2}}{e^{i\alpha/2}\cdot2i\sin\frac{\alpha}{2}}
=\frac{\sin\frac{n\alpha}{2}}{\sin\frac{\alpha}{2}}\,e^{i\left(\theta+\frac{(n-1)\alpha}{2}\right)}.
\]
Taking real and imaginary parts:
\[
\boxed{\;C=\frac{\sin\frac{n\alpha}{2}}{\sin\frac{\alpha}{2}}\cos\!\left(\theta+\tfrac{(n-1)\alpha}{2}\right),\qquad
S=\frac{\sin\frac{n\alpha}{2}}{\sin\frac{\alpha}{2}}\sin\!\left(\theta+\tfrac{(n-1)\alpha}{2}\right).\;}
\]
Note the elegance: the sum behaves like a single cosine (or sine) taken at the \emph{average} angle $\theta+\frac{(n-1)\alpha}{2}$, scaled by the factor $\dfrac{\sin(n\alpha/2)}{\sin(\alpha/2)}$.
\end{exemplebox}

\begin{attention}[Check the ratio before summing]
The geometric-series formula requires the common ratio $e^{i\alpha}\neq1$, i.e. $\alpha$ not a multiple of $2\pi$. If $\alpha=2m\pi$, every term equals $\cos\theta$ (or $\sin\theta$) and the sums are simply $C=n\cos\theta$, $S=n\sin\theta$. Always treat this degenerate case separately.
\end{attention}

\begin{exercice}[Practice]
\begin{enumerate}
\item Solve $z^5=32$ in $\mathbb{C}$ and represent the roots on an Argand diagram.
\item Show that $\cos5\theta=16\cos^5\theta-20\cos^3\theta+5\cos\theta$.
\item Linearise $\sin^3\theta$ and hence find $\displaystyle\int_0^{\pi}\sin^3\theta\,d\theta$.
\item Using $\omega=e^{2i\pi/3}$, verify directly that $1+\omega+\omega^2=0$.
\item Sum $\displaystyle\sum_{k=0}^{n-1}\cos k\theta$ for $\theta\neq2m\pi$.
\end{enumerate}
\end{exercice}

\begin{corrige}[Exercises 1--5]
\textbf{1.} $32=32e^{i\cdot0}$, so $z_k=2e^{2ik\pi/5}$, $k=0,\dots,4$: the fifth roots are $2$, $2e^{2i\pi/5}$, $2e^{4i\pi/5}$, $2e^{6i\pi/5}$, $2e^{8i\pi/5}$, vertices of a regular pentagon on the circle of radius $2$.
\textbf{2.} Expand $(c+is)^5$: real part $c^5-10c^3s^2+5cs^4$; replace $s^2=1-c^2$: $c^5-10c^3(1-c^2)+5c(1-c^2)^2=16c^5-20c^3+5c$.
\textbf{3.} $\sin^3\theta=\frac{1}{(2i)^3}(e^{i\theta}-e^{-i\theta})^3=\frac{3\sin\theta-\sin3\theta}{4}$; the integral is $\left[-\frac{3\cos\theta}{4}+\frac{\cos3\theta}{12}\right]_0^\pi=\frac{4}{3}$.
\textbf{4.} $\omega=-\frac12+i\frac{\sqrt3}{2}$, $\omega^2=-\frac12-i\frac{\sqrt3}{2}$; sum $=1-1+0i=0$.
\textbf{5.} With $\theta$ playing the role of $\alpha$ and initial angle $0$: $\displaystyle\sum_{k=0}^{n-1}\cos k\theta=\frac{\sin\frac{n\theta}{2}}{\sin\frac{\theta}{2}}\cos\frac{(n-1)\theta}{2}$.
\end{corrige}

\begin{aretenir}[Section summary]
\begin{itemize}
\item $z^n=w=re^{i\theta}$ has exactly $n$ roots: $z_k=r^{1/n}e^{i(\theta+2k\pi)/n}$, $k=0,\dots,n-1$ — vertices of a regular $n$-gon.
\item The $n$th roots of unity are $1,\omega,\dots,\omega^{n-1}$ with $\omega=e^{2i\pi/n}$; their sum is $0$.
\item De Moivre $+$ binomial theorem $\Rightarrow$ polynomial forms of $\cos n\theta$, $\sin n\theta$, $\tan n\theta$.
\item Euler's formulae $+$ binomial theorem $\Rightarrow$ linearisation of $\cos^n\theta$, $\sin^n\theta$ (ideal for integration).
\item Trigonometric sums are real/imaginary parts of geometric series of exponentials.
\end{itemize}
\end{aretenir}

\coursec{Loci in the Argand Diagram}

In the previous section we exploited De Moivre's theorem to extract roots, linearise trigonometric expressions and sum series. In every one of those applications the complex number $z$ was treated as an \emph{algebraic} object. We now return to the geometric point of view: a complex number is also a \emph{point} of the plane. The central question of this section is the following: \emph{if a complex number $z$ is required to satisfy a condition involving a modulus or an argument, what is the set of points $M(z)$ that satisfy it?} Such a set is called a \cle{locus}. Locus questions are among the most frequently examined items in the GCE Advanced Level paper, and they are also the gateway to the transformations of the $z$-plane that we study in the next section.

\courssub{The Affine Plane and Affixes}

Recall that to every complex number $z = x + iy$ (with $x, y \in \mathbb{R}$) we associate the point $M(x, y)$ of the plane equipped with an orthonormal frame $(O; \vec{u}, \vec{v})$. This plane, viewed as the set of affixes of points, is called the \cle{affine complex plane} (or Argand plane).

\begin{definition}[Affix of a point and of a vector]
Let the plane be referred to an orthonormal frame $(O; \vec{u}, \vec{v})$.
\begin{itemize}
\item The point $M(x, y)$ is said to have \cle{affix} $z_M = x + iy$. We write $M(z_M)$.
\item The vector $\vec{w}\begin{pmatrix} a \\ b \end{pmatrix}$ has affix $z_{\vec{w}} = a + ib$.
\item If $A(z_A)$ and $B(z_B)$ are two points, then the vector $\overrightarrow{AB}$ has affix
\[ z_{\overrightarrow{AB}} = z_B - z_A. \]
\item The midpoint $I$ of the segment $[AB]$ has affix
\[ z_I = \frac{z_A + z_B}{2}. \]
\end{itemize}
\end{definition}

These formulas are immediate consequences of the coordinate formulas $\overrightarrow{AB}(x_B - x_A,\; y_B - y_A)$ and $I\left(\tfrac{x_A + x_B}{2}, \tfrac{y_A + y_B}{2}\right)$, but they are far more powerful in complex form because they combine with the modulus and argument.

\begin{propriete}[Modulus and argument of a difference]
Let $A(z_A)$ and $B(z_B)$. Then
\[ AB = |z_B - z_A| \qquad \text{and} \qquad (\vec{u}, \overrightarrow{AB}) \equiv \arg(z_B - z_A) \pmod{2\pi}. \]
\end{propriete}

This single property is the key to the whole section: \textbf{a condition on $|z - a|$ is a condition on a distance, and a condition on $\arg(z - a)$ is a condition on a direction.}

\courssub{Circles: the locus $|z - a| = r$}

\begin{propriete}[Circle]
Let $A(a)$ be a fixed point and $r > 0$. The locus of points $M(z)$ such that
\[ |z - a| = r \]
is the \cle{circle} of centre $A$ and radius $r$. The condition $|z - a| \leq r$ describes the closed disc, and $|z - a| < r$ the open disc.
\end{propriete}

The proof is one line: $|z - a| = |z_M - z_A| = AM$, so the condition says exactly $AM = r$.

\begin{exemplebox}
The equation $|z - 2 - i| = 3$ represents the circle of centre $C(2 + i)$, i.e.\ $C(2, 1)$, and radius $3$. Writing $z = x + iy$, we recover the Cartesian equation $(x - 2)^2 + (y - 1)^2 = 9$.
\end{exemplebox}

\begin{popfigure}
\begin{center}
\begin{tikzpicture}[scale=0.85]
\draw[->, gray] (-1.5,0) -- (6,0) node[below left] {Re};
\draw[->, gray] (0,-2.5) -- (0,4.5) node[below left] {Im};
\draw[thick, popblue] (2,1) circle (3);
\fill[popdark] (2,1) circle (1.5pt) node[above right, black] {$C(2+i)$};
\fill[poporange] (4.6,2.8) circle (1.5pt);
\draw[poporange, thick] (2,1) -- (4.6,2.8) node[midway, above, sloped] {$r=3$};
\node[black] at (4.9,3.1) {$M(z)$};
\node[popblue] at (2,-2.1) {$|z-2-i|=3$};
\end{tikzpicture}
\end{center}
\legende{The locus $|z - 2 - i| = 3$: circle of centre $C(2,1)$ and radius $3$.}
\end{popfigure}

\courssub{The Perpendicular Bisector: $|z - a| = |z - b|$}

Suppose a condition says that $M(z)$ is \emph{as far from $A$ as from $B$}. Geometric intuition (think of two radio transmitters in Yaounde and Douala: the set of positions receiving both equally strongly) tells us the answer is a straight line.

\begin{propriete}[Perpendicular bisector]
Let $A(z_A)$ and $B(z_B)$ be distinct points. The locus of points $M(z)$ such that
\[ |z - z_A| = |z - z_B| \]
is the \cle{perpendicular bisector} of the segment $[AB]$: the line through the midpoint $I\left(\tfrac{z_A + z_B}{2}\right)$ perpendicular to $(AB)$.
\end{propriete}

\textbf{Proof (examinable).} The condition $|z - z_A| = |z - z_B|$ means $MA = MB$. Squaring with $z = x + iy$, $z_A = x_A + iy_A$, $z_B = x_B + iy_B$:
\[ (x - x_A)^2 + (y - y_A)^2 = (x - x_B)^2 + (y - y_B)^2, \]
which simplifies to
\[ 2(x_B - x_A)x + 2(y_B - y_A)y = x_B^2 + y_B^2 - x_A^2 - y_A^2, \]
the equation of a line whose normal vector is $\overrightarrow{AB}(x_B - x_A, y_B - y_A)$ and which passes through $I$ (substitute the coordinates of $I$: both sides agree). Hence it is the perpendicular bisector of $[AB]$. $\blacksquare$

\begin{exoresolu}[A perpendicular bisector locus]
Describe and sketch the locus of $M(z)$ such that $|z - 1| = |z + 3i|$.
\tcblower
The condition is $|z - 1| = |z - (-3i)|$, i.e.\ $MA = MB$ with $A(1, 0)$ and $B(0, -3)$. The locus is the perpendicular bisector of $[AB]$. The midpoint is $I\left(\tfrac{1}{2}, -\tfrac{3}{2}\right)$. Since $\overrightarrow{AB}(-1, -3)$, a point $M(x,y)$ of the bisector satisfies $\overrightarrow{IM} \cdot \overrightarrow{AB} = 0$:
\[ -\left(x - \tfrac{1}{2}\right) - 3\left(y + \tfrac{3}{2}\right) = 0 \quad\Longleftrightarrow\quad x + 3y + 4 = 0. \]
The locus is the line $x + 3y + 4 = 0$, passing through $I(0.5, -1.5)$, perpendicular to $(AB)$.
\end{exoresolu}

\begin{popfigure}
\begin{center}
\begin{tikzpicture}[scale=0.9]
\draw[->, gray] (-4.5,0) -- (3.5,0) node[below left] {Re};
\draw[->, gray] (0,-4.2) -- (0,2.5) node[below left] {Im};
\fill[popdark] (1,0) circle (1.5pt) node[above right, black] {$A(1)$};
\fill[popdark] (0,-3) circle (1.5pt) node[below right, black] {$B(-3i)$};
\draw[popmuted, dashed] (1,0) -- (0,-3);
\fill[popgreen] (0.5,-1.5) circle (1.5pt) node[right, black] {$I$};
\draw[thick, popblue] (-4.17,-0.28) -- (2.83,-2.28);
\node[popblue] at (-3.4,-0.75) {$|z-1|=|z+3i|$};
\end{tikzpicture}
\end{center}
\legende{The perpendicular bisector of $[AB]$ with $A(1)$, $B(-3i)$.}
\end{popfigure}

\courssub{Circles of Apollonius: $|z - a| = k|z - b|$}

What happens when the distances are in a fixed ratio $k \neq 1$? The locus is no longer a line but a circle — the \cle{circle of Apollonius}. This is a GCE-useful enrichment: the question appears regularly as a structured exercise.

\begin{propriete}[Circle of Apollonius]
Let $A(z_A)$, $B(z_B)$ distinct and $k > 0$, $k \neq 1$. The locus of $M(z)$ such that
\[ |z - z_A| = k\,|z - z_B| \]
is a circle whose diameter is the segment joining the two points of the line $(AB)$ that divide $[AB]$ in the ratio $k$ (one internally, one externally). Its centre lies on the line $(AB)$.
\end{propriete}

\textbf{Derivation (examinable in outline).} Squaring $MA^2 = k^2 MB^2$ with $z = x + iy$ gives
\[ (x - x_A)^2 + (y - y_A)^2 = k^2\big[(x - x_B)^2 + (y - y_B)^2\big]. \]
Expanding and grouping terms:
\[ (1 - k^2)(x^2 + y^2) - 2(x_A - k^2 x_B)x - 2(y_A - k^2 y_B)y + (x_A^2 + y_A^2) - k^2(x_B^2 + y_B^2) = 0. \]
Since $k \neq 1$, we divide by $1 - k^2$ to obtain an equation of the form $x^2 + y^2 + ux + vy + w = 0$, which is a circle (provided the radius squared is positive, which it is for distinct $A,B$). The centre is $\left(\frac{x_A - k^2 x_B}{1 - k^2}, \frac{y_A - k^2 y_B}{1 - k^2}\right)$, which lies on line $AB$. $\blacksquare$

\begin{exemplebox}
The locus $|z - 4| = 2|z - 1|$: squaring with $z = x + iy$,
\[ (x-4)^2 + y^2 = 4\big[(x-1)^2 + y^2\big] \Longleftrightarrow 3x^2 + 3y^2 - 12 = 0 \Longleftrightarrow x^2 + y^2 = 4. \]
The locus is the circle centred at $O$ of radius $2$. (Here $A(4,0)$, $B(1,0)$, $k=2$; the internal division point is $2$, the external is $-2$; the circle has diameter $[-2,2]$.)
\end{exemplebox}

\courssub{Argument Loci: Half-Lines and Circles}

\begin{propriete}[Half-line]
Let $A(a)$ and $\theta \in \mathbb{R}$. The locus of points $M(z)$, $M \neq A$, such that
\[ \arg(z - a) \equiv \theta \pmod{2\pi} \]
is the \cle{half-line} starting at $A$ (the point $A$ itself \textbf{excluded}) making the angle $\theta$ with the positive real axis.
\end{propriete}

\begin{attention}[The half-line excludes its origin]
A very common mistake is to draw the whole ray \emph{including} $A$, or even the full line. But $\arg(z - a)$ is \textbf{not defined} when $z = a$, since $\arg(0)$ does not exist. Mark $A$ with an open (hollow) dot. Similarly, $\arg(z - a) = \theta$ is only \emph{half} of the line through $A$ in direction $\theta$: the opposite half corresponds to $\arg(z - a) = \theta + \pi$.
\end{attention}

For loci involving \emph{two} fixed points, the key tool is the argument of a quotient.

\begin{propriete}[Angle between vectors via a quotient of affixes]
Let $A(z_A)$, $B(z_B)$, $C(z_C)$ be distinct points. Then
\[ \arg\!\left(\frac{z_C - z_A}{z_C - z_B}\right) \equiv (\overrightarrow{CB}, \overrightarrow{CA}) \pmod{2\pi}, \]
i.e.\ this argument measures the angle $\widehat{BCA}$ at the vertex $C$. In particular, with $M(z)$ a variable point,
\[ \arg\!\left(\frac{z - z_A}{z - z_B}\right) \equiv (\overrightarrow{MB}, \overrightarrow{MA}) \pmod{2\pi}. \]
\end{propriete}

\textbf{Proof.} $\arg\!\left(\dfrac{z - z_A}{z - z_B}\right) = \arg(z - z_A) - \arg(z - z_B) = (\vec{u}, \overrightarrow{MA}) - (\vec{u}, \overrightarrow{MB}) = (\overrightarrow{MB}, \overrightarrow{MA})$ by the Chasles relation on angles. $\blacksquare$

\begin{propriete}[Thales' circle and arcs of circles]
Let $A(z_A)$, $B(z_B)$ be distinct.
\begin{itemize}
\item The locus of $M(z)$, $M \notin \{A, B\}$, such that
\[ \arg\!\left(\frac{z - z_A}{z - z_B}\right) \equiv \pm\frac{\pi}{2} \pmod{\pi} \]
is the \cle{circle of diameter $[AB]$} (Thales' circle), with $A$ and $B$ excluded.
\item More generally, for a constant $\alpha$ with $0 < \alpha < \pi$, the condition
\[ \arg\!\left(\frac{z - z_A}{z - z_B}\right) \equiv \alpha \pmod{2\pi} \]
describes \textbf{one arc} of the circle through $A$ and $B$ (the arc on which the segment $[AB]$ subtends an angle $\alpha$). The condition $\equiv \alpha \pmod{\pi}$ gives the \textbf{whole circle} (both arcs) minus $A$ and $B$.
\end{itemize}
\end{propriete}

Indeed, $(\overrightarrow{MB}, \overrightarrow{MA}) = \pm \tfrac{\pi}{2}$ means $MA \perp MB$, and the set of points from which a segment $[AB]$ is seen at a right angle is exactly the circle of diameter $[AB]$. The inscribed angle theorem guarantees that a constant angle $\alpha$ (mod $2\pi$) corresponds to one arc.

\begin{exemplebox}
The locus $\arg\!\left(\dfrac{z - 2}{z + 2}\right) = \dfrac{\pi}{2}$ is the circle of diameter $[AB]$ with $A(2, 0)$, $B(-2, 0)$: the circle centred at $O$ of radius $2$, with the points $(\pm 2, 0)$ excluded.
\end{exemplebox}

\courssub{Regions Defined by Inequalities}

Inequalities define \emph{regions} of the plane rather than curves. The standard vocabulary:

\begin{aretenir}[Standard regions]
\begin{itemize}
\item $|z - a| \leq r$: \cle{closed disc}, centre $A(a)$, radius $r$ (boundary circle \textbf{included}: solid line).
\item $|z - a| < r$: \cle{open disc} (boundary \textbf{excluded}: dashed line).
\item $r_1 \leq |z| \leq r_2$: \cle{annulus} centred at $O$, between the circles of radii $r_1$ and $r_2$.
\item $\alpha \leq \arg(z) \leq \beta$: \cle{sector} (wedge) of the plane between the half-lines of arguments $\alpha$ and $\beta$, with $O$ excluded.
\end{itemize}
\end{aretenir}

\begin{attention}[Open versus closed boundaries]
Do not confuse $|z - a| < r$ (open disc: shade the interior and draw the boundary circle \textbf{dashed}) with $|z - a| \leq r$ (closed disc: boundary \textbf{solid}). The same convention applies to strict inequalities on arguments. In the GCE, an incorrectly drawn boundary costs marks even when the region is correctly identified.
\end{attention}

\begin{methode}[Identifying and sketching a locus]
\begin{enumerate}
\item \textbf{Translate} the condition geometrically: $|z - a|$ is the distance $AM$; $\arg(z - a)$ is the direction of $\overrightarrow{AM}$; a quotient $\tfrac{z - a}{z - b}$ encodes the angle at $M$.
\item \textbf{Identify} the type: modulus equal to a constant $\to$ circle; two equal moduli $\to$ perpendicular bisector; ratio of moduli $\to$ Apollonius circle; constant argument $\to$ half-line; right-angle argument $\to$ Thales' circle; constant argument (mod $2\pi$) of a quotient $\to$ arc of circle; inequality $\to$ region.
\item \textbf{Determine} the characteristic elements: centre, radius, midpoint, angle, division points for Apollonius.
\item \textbf{Sketch} on an Argand diagram: mark fixed points, draw the locus, respect boundary conventions (solid/dashed, hollow dots for excluded points), and \textbf{shade} the required region for inequalities — shade a test point (e.g.\ $z = 0$) to check which side of a boundary is wanted.
\item \textbf{Verify} algebraically with $z = x + iy$ if in doubt.
\end{enumerate}
\end{methode}

\begin{exoresolu}[Intersection of loci: annulus and sector]
Shade the region of the Argand diagram defined by
\[ 1 \leq |z| \leq 3 \qquad \text{and} \qquad 0 \leq \arg(z) \leq \frac{\pi}{3}. \]
\tcblower
The first condition is the closed annulus between the circles centred at $O$ of radii $1$ and $3$ (both boundaries included, drawn solid). The second is the sector between the positive real axis and the half-line from $O$ at angle $\tfrac{\pi}{3}$. The required region is the \emph{intersection}: the portion of the annulus lying inside the sector. Both boundary circles are solid; both bounding half-lines are solid; the origin is excluded (as $\arg(0)$ is undefined) but this does not affect the shaded region.
\end{exoresolu}

\begin{popfigure}
\begin{center}
\begin{tikzpicture}[scale=1.0]
\fill[popblueL] (0:1) arc (0:60:1) -- (60:3) arc (60:0:3) -- cycle;
\draw[thick, popblue] (0,0) circle (1);
\draw[thick, popblue] (0,0) circle (3);
\draw[thick, popblue] (0,0) -- (3.4,0);
\draw[thick, popblue] (0,0) -- (60:3.4);
\draw[->, gray] (-3.8,0) -- (4,0) node[below left] {Re};
\draw[->, gray] (0,-1.2) -- (0,4) node[below left] {Im};
\draw[poporange] (0.55,0) arc (0:60:0.55) node[midway, right, poporange] {$\frac{\pi}{3}$};
\node[black] at (1.5,0.55) {$1$};
\node[black] at (2.6,-0.28) {$3$};
\node[black] at (1.9,1.35) {$R$};
\end{tikzpicture}
\end{center}
\legende{The region $R$: $1 \leq |z| \leq 3$ and $0 \leq \arg(z) \leq \tfrac{\pi}{3}$ (annular sector).}
\end{popfigure}

\begin{exercice}
On the same Argand diagram, sketch and shade the set of points $M(z)$ satisfying simultaneously $|z - 2i| \leq 2$ and $|z| > 1$. State precisely which parts of the boundary are included.
\end{exercice}

\begin{corrige}[Exercise 1]
$|z - 2i| \leq 2$ is the closed disc of centre $(0, 2)$ and radius $2$ (it passes through $O$); $|z| > 1$ is the exterior of the open disc of radius $1$. The required set is the part of the first disc lying outside or on the unit circle: the circle of centre $(0,2)$ is drawn solid, the unit circle dashed (strict inequality), and the region between them inside the disc is shaded.
\end{corrige}

\begin{savaistu}[Why loci matter beyond the exam]
Locus conditions are the language of many technologies. The set of points where a mobile phone signal from two relay antennas arrives with equal strength is a perpendicular bisector; points where the \emph{ratio} of signal powers is constant lie on Apollonius circles — a fact used in positioning and in the design of cellular networks. The same mathematics will reappear in the next section, where we ask: what becomes of these circles and lines when the plane is transformed by a mapping $z \mapsto w = f(z)$?
\end{savaistu}

\coursec{Transformations from the z-plane to the w-plane}

In the previous section we studied \emph{loci}: given a condition on $z$ (for example $|z - 2i| = 3$), we described the set of points of the Argand plane satisfying it. We now change our point of view. Instead of asking \emph{where} a point may lie, we ask \emph{where it goes}: we feed each point $z$ of the plane into a rule, and a new point $w$ comes out. The rule is a \cle{transformation} of the plane, and complex numbers make such transformations astonishingly easy to handle, because a single formula like $w = (1+i)z + 3$ encodes a rotation, an enlargement and a translation all at once.

\begin{definition}[Transformation from the z-plane to the w-plane]
A \cle{transformation from the $z$-plane to the $w$-plane} is a rule $T$ which associates to each complex number $z = x + iy$ (in its domain) a unique complex number
\[ w = T(z) = u + iv, \]
where $u$ and $v$ are real functions of $x$ and $y$. We picture $z$ in one Argand diagram (the \cle{$z$-plane}, with axes $x$ and $y$) and its image $w$ in a second Argand diagram (the \cle{$w$-plane}, with axes $u$ and $v$). If $C$ is a curve in the $z$-plane, its \cle{image} is the set $T(C) = \{T(z) : z \in C\}$ in the $w$-plane.
\end{definition}

The central question of this section is: \emph{given a curve (line, circle, ray) in the $z$-plane and a map $T$, what is the equation of the image curve in the $w$-plane?}

\courssub{The three elementary maps}

\textbf{Translation: $w = z + b$.}\par
Let $b = p + iq$ be fixed. Writing $z = x + iy$, we get $w = (x + p) + i(y + q)$: every point is shifted by the vector $\begin{pmatrix} p \\ q \end{pmatrix}$. Geometrically, $w - z = b$ means the displacement from $z$ to $w$ is the same everywhere. Translations preserve distances, angles, shapes and sizes.

\textbf{Rotation composed with enlargement: $w = az$, with $a = re^{i\alpha}$ ($r>0$).}\par
This is where the exponential form shines. If $z = \rho e^{i\theta}$, then
\[ w = az = r\rho\, e^{i(\theta + \alpha)}. \]
So $|w| = r|z|$ and $\arg w = \arg z + \alpha$: the map \emph{enlarges distances from the origin by the factor $r = |a|$} and \emph{rotates the plane through the angle $\alpha = \arg a$} about $O$. This is exactly a \cle{direct similitude} centred at $O$, which we shall study formally later in this chapter. (If $a$ is real and positive, it is a pure enlargement; if $|a|=1$, it is a pure rotation.)

\textbf{Reflection in the real axis: $w = \bar{z}$.}\par
If $z = \rho e^{i\theta}$ then $\bar{z} = \rho e^{-i\theta}$: the modulus is kept and the argument is negated. Every point is mirrored across the $x$-axis. Note that this map \emph{reverses orientation}: a triangle traversed anticlockwise is sent to a triangle traversed clockwise.

\begin{exemplebox}[Image of a point under $w = (1+i)z + 3$]
Take $z_0 = 2 - i$. Since $1 + i = \sqrt{2}\,e^{i\pi/4}$, the map is: enlarge by $\sqrt{2}$, rotate by $\dfrac{\pi}{4}$, then translate by $3$. Directly:
\[ w_0 = (1+i)(2 - i) + 3 = (2 - i + 2i - i^2) + 3 = (3 + i) + 3 = 6 + i. \]
Check with the geometric description: $|z_0| = \sqrt{5}$, so $|w_0 - 3| = \sqrt{2}\cdot\sqrt{5} = \sqrt{10}$, and indeed $|3 + i| = \sqrt{10}$. \checkmark
\end{exemplebox}

\courssub{Image of a line and of a circle under $w = az + b$}

Because $w = az + b$ is a similitude (rotation + enlargement + translation), it preserves shapes: \textbf{lines map to lines and circles map to circles}. Algebraically this is easy to see because the map is invertible: $z = \dfrac{w - b}{a}$ (provided $a \neq 0$).

\begin{exoresolu}[Image of the circle $|z| = 2$ under $w = (1+i)z + 3$]
Find the image of the circle $|z| = 2$ under the map $w = (1 + i)z + 3$, and describe it geometrically.
\tcblower
\textbf{Step 1: invert the map.} $z = \dfrac{w - 3}{1 + i}$.

\textbf{Step 2: substitute into the locus condition.}
\[ |z| = 2 \iff \left|\frac{w - 3}{1 + i}\right| = 2 \iff \frac{|w - 3|}{|1 + i|} = 2 \iff |w - 3| = 2\sqrt{2}. \]

\textbf{Step 3: interpret.} The image is the circle of centre $3 + 0i$ (the image of the centre $0$) and radius $2\sqrt{2} = 2|1+i|$ (the original radius multiplied by the scale factor $\sqrt{2}$). Geometrically: enlarge by $\sqrt{2}$ (radius $2 \to 2\sqrt{2}$), rotate by $\pi/4$ (centre $0$ unchanged), translate by $3$ (centre $0 \to 3$).
\end{exoresolu}

\begin{exemplebox}[Image of a line under $w = az + b$]
Find the image of the line $\operatorname{Re}(z) = 1$ under $w = 2iz + 1$. Invert: $z = \dfrac{w-1}{2i}$. Write $w = u + iv$: $z = \dfrac{(u-1) + iv}{2i} = \dfrac{v}{2} - i\,\dfrac{u-1}{2}$. Hence $\operatorname{Re}(z) = \dfrac{v}{2} = 1$, i.e.\ $v = 2$: the image is the horizontal line $\operatorname{Im}(w) = 2$. Geometrically, the vertical line was rotated by $\pi/2$ (becoming horizontal), stretched by $2$, then shifted up by $1$.
\end{exemplebox}

\begin{popfigure}
\begin{center}
\begin{tikzpicture}[scale=0.85]
% z-plane
\begin{scope}[xshift=-4.2cm]
\draw[->,gray] (-2.2,0) -- (2.2,0) node[right]{$x$};
\draw[->,gray] (0,-2.2) -- (0,2.2) node[above]{$y$};
\draw[popblue,thick] (0,0) circle (1.2);
\fill[popblue] (0.85,0.85) circle (1.5pt) node[above right,popblue]{$z$};
\fill (0,0) circle (1.2pt) node[below left]{$O$};
\node[popblue] at (0,-2.7) {$z$-plane: $|z| = r$};
\end{scope}
% arrow
\draw[->,very thick,poporange] (-1.6,0) -- (1.2,0) node[midway,above,align=center,text width=2.6cm]{$w = 2e^{i\pi/4}z$};
% w-plane
\begin{scope}[xshift=4.2cm]
\draw[->,gray] (-3,0) -- (3,0) node[right]{$u$};
\draw[->,gray] (0,-3) -- (0,3) node[above]{$v$};
\draw[poppurple,thick] (0,0) circle (2.4);
\fill[poppurple] (1.2,2.08) circle (1.5pt) node[above right,poppurple]{$w$};
\draw[->,popmuted,dashed] (0,0) -- (0.85,0.85);
\draw[->,popmuted,dashed] (0,0) -- (1.2,2.08);
\draw[popgold,thick] (0.55,0.55) arc (45:60:0.78);
\node[popgold] at (1.15,0.72) {$\frac{\pi}{4}$};
\fill (0,0) circle (1.2pt) node[below left]{$O$};
\node[poppurple] at (0,-3.3) {$w$-plane: $|w| = 2r$};
\end{scope}
\end{tikzpicture}
\end{center}
\legende{Under $w = 2e^{i\pi/4}z$, the circle $|z| = r$ is enlarged by the factor $2$ and rotated through $\dfrac{\pi}{4}$; since a circle centred at $O$ is invariant under rotation about $O$, only the enlargement is visible: the image is $|w| = 2r$.}
\end{popfigure}

\courssub{The reciprocal map $w = \dfrac{1}{z}$}

The map $z \mapsto \dfrac{1}{z}$ is defined on $\mathbb{C} \setminus \{0\}$ and is called \cle{inversion} (composed with reflection in the real axis). Its behaviour is completely described by the modulus–argument analysis.

\begin{propriete}[Modulus and argument under $w = 1/z$]
For $z \neq 0$, writing $z = \rho e^{i\theta}$,
\[ w = \frac{1}{z} = \frac{1}{\rho}e^{-i\theta}, \qquad\text{so}\qquad |w| = \frac{1}{|z|} \quad\text{and}\quad \arg w = -\arg z. \]
In particular: points far from $O$ map close to $O$ and vice versa; the upper half-plane maps to the lower half-plane; and the \textbf{unit circle is invariant} ($|z| = 1 \Rightarrow |w| = 1$), each point $e^{i\theta}$ being sent to $e^{-i\theta}$.
\end{propriete}

\textbf{Image of the line $\operatorname{Re}(z) = c$ ($c \neq 0$).}\par
Let $z = c + iy$ with $y \in \mathbb{R}$, and write $w = u + iv = \dfrac{1}{c + iy}$. Then
\[ u = \frac{c}{c^2 + y^2}, \qquad v = \frac{-y}{c^2 + y^2}. \]
Compute:
\[ u^2 + v^2 = \frac{c^2 + y^2}{(c^2 + y^2)^2} = \frac{1}{c^2 + y^2} = \frac{u}{c}. \]
Hence $u^2 + v^2 - \dfrac{u}{c} = 0$, i.e.
\[ \left(u - \frac{1}{2c}\right)^2 + v^2 = \left(\frac{1}{2c}\right)^2 : \]
a \textbf{circle through the origin}, of centre $\dfrac{1}{2c}$ on the real axis and radius $\dfrac{1}{2|c|}$. The origin itself is excluded (it would correspond to $z = \infty$).

\begin{attention}[$w = 1/z$ is undefined at $z = 0$]
The reciprocal map is only defined for $z \neq 0$. Consequently: a line \emph{through} the origin maps to a line through the origin (minus $O$), since $\arg w = -\arg z$ keeps the points on a ray; but a line \emph{not} through the origin maps to a \emph{circle through the origin} (with $O$ excluded). Never write ``the image of a line is a line'' without checking whether the line passes through $O$. Similarly, a circle through $O$ maps to a line, and a circle not through $O$ maps to a circle.
\end{attention}

\begin{exoresolu}[Image of $\operatorname{Re}(z) = 1$ under $w = 1/z$]
Find the image of the line $\operatorname{Re}(z) = 1$ under $w = \dfrac{1}{z}$.
\tcblower
Parametrise: $z = 1 + iy$, $y \in \mathbb{R}$. Then with $w = u + iv$:
\[ w = \frac{1}{1 + iy} = \frac{1 - iy}{1 + y^2} \implies u = \frac{1}{1+y^2},\quad v = \frac{-y}{1+y^2}. \]
Eliminate $y$:
\[ u^2 + v^2 = \frac{1 + y^2}{(1+y^2)^2} = \frac{1}{1+y^2} = u. \]
So $u^2 - u + v^2 = 0$, i.e.\ $\left(u - \tfrac12\right)^2 + v^2 = \tfrac14$. The image is the circle of centre $\dfrac{1}{2}$ and radius $\dfrac{1}{2}$, passing through the origin (the origin itself excluded, since $u = \frac{1}{1+y^2} > 0$ for all finite $y$).
\end{exoresolu}

\begin{popfigure}
\begin{center}
\begin{tikzpicture}[scale=1.15]
% z-plane
\begin{scope}[xshift=-3.6cm]
\draw[->,gray] (-1.6,0) -- (2.4,0) node[right]{$x$};
\draw[->,gray] (0,-2.2) -- (0,2.2) node[above]{$y$};
\draw[popblue,thick] (1,-2.2) -- (1,2.2);
\fill[popblue] (1,0.9) circle (1.5pt) node[right,popblue]{$z=1+iy$};
\node[popblue] at (0.5,-2.7) {$z$-plane: $\operatorname{Re}(z)=1$};
\end{scope}
\draw[->,very thick,poporange] (-1.3,0) -- (0.9,0) node[midway,above]{$w=\dfrac{1}{z}$};
% w-plane
\begin{scope}[xshift=3.6cm]
\draw[->,gray] (-1.6,0) -- (2.4,0) node[right]{$u$};
\draw[->,gray] (0,-2.2) -- (0,2.2) node[above]{$v$};
\draw[poppurple,thick] (0.5,0) circle (0.5);
\fill[poppurple] (0.79,-0.31) circle (1.5pt) node[right,poppurple]{$w$};
\draw[popmuted] (0,0) circle (1.5pt);
\node[below left] at (0,0) {$O$ (excluded)};
\fill (0.5,0) circle (1pt) node[above]{$\frac12$};
\node[poppurple] at (0.4,-2.7) {$w$-plane: $\left(u-\frac12\right)^2+v^2=\frac14$};
\end{scope}
\end{tikzpicture}
\end{center}
\legende{The line $\operatorname{Re}(z) = 1$ (which avoids the origin) is mapped by $w = 1/z$ to the circle of centre $\tfrac12$ and radius $\tfrac12$, passing through $O$. The point $z = 1 + iy$ shown maps to $w$ in the lower half-plane, since $\arg w = -\arg z$.}
\end{popfigure}

\courssub{The map $w = z^2$}

If $z = \rho e^{i\theta}$, then $w = \rho^2 e^{2i\theta}$: \textbf{moduli are squared, arguments are doubled}.
\begin{itemize}
\item \textbf{Unit circle:} $|z| = 1 \Rightarrow |w| = 1$; the image is again the unit circle, but \emph{described twice} as $\theta$ runs from $0$ to $2\pi$ (the map is two-to-one: $z$ and $-z$ have the same image).
\item \textbf{Rays from $O$:} the ray $\arg z = \theta_0$ maps to the ray $\arg w = 2\theta_0$. In particular the positive real axis is fixed (as a set), and the positive imaginary axis ($\theta_0 = \pi/2$) maps to the \emph{negative} real axis.
\item \textbf{Circle $|z| = r$:} maps to $|w| = r^2$, described twice.
\end{itemize}

\begin{attention}[Not every map sends circles to circles]
The maps $w = az + b$, $w = \bar z$ and $w = 1/z$ send every line or circle to a line or circle. This is \emph{false} for $w = z^2$ in general. For instance, the line $\operatorname{Re}(z) = 1$ maps under $w = z^2$ to a \emph{parabola}: with $z = 1 + iy$, $w = (1 - y^2) + 2iy$, so $u = 1 - y^2$, $v = 2y$, and eliminating $y$ gives $u = 1 - \dfrac{v^2}{4}$ — certainly not a circle. Always compute; never assume.
\end{attention}

\courssub{General method: parametrise, substitute, eliminate}

\begin{methode}[Finding the image of a curve under a complex map]
\begin{enumerate}
\item \textbf{Parametrise} the given locus in the $z$-plane (e.g.\ $z = c + iy$ for a vertical line, $z = z_0 + re^{it}$ for a circle), or write its equation in terms of $z$ and $\bar z$.
\item \textbf{Substitute} into $w = T(z)$, or invert the map and substitute $z = T^{-1}(w)$ into the locus condition — whichever is simpler.
\item \textbf{Eliminate} the parameter (or simplify the condition) to obtain a relation between $u = \operatorname{Re}(w)$ and $v = \operatorname{Im}(w)$.
\item \textbf{Identify} the image curve and state any excluded points (e.g.\ the origin for $w = 1/z$).
\end{enumerate}
\end{methode}

\courssub{Möbius maps: a powerful class of transformations}

\begin{definition}[Linear fractional (Möbius) map]
A \cle{Möbius map} is a transformation of the form
\[ w = \frac{az + b}{cz + d}, \qquad ad - bc \neq 0, \]
defined for $z \neq -\dfrac{d}{c}$ (when $c \neq 0$). The condition $ad - bc \neq 0$ guarantees the map is not constant.
\end{definition}

Every Möbius map is a composition of translations, inversions $z \mapsto 1/z$ and maps $z \mapsto az$: indeed, for $c \neq 0$,
\[ \frac{az+b}{cz+d} = \frac{a}{c} + \frac{bc - ad}{c}\cdot\frac{1}{cz + d}. \]
Since each ingredient sends lines and circles to lines or circles, so does every Möbius map. This is a powerful exam shortcut: to find the image of a circle, it is enough to find the images of \emph{three} of its points.

\begin{exemplebox}[A Möbius image]
Under $w = \dfrac{z}{z + 1}$, find the image of the circle $|z| = 1$ (note $z = -1$ is on the circle, so expect a line). Invert: $z = \dfrac{w}{1 - w}$. Then $|z| = 1 \iff |w| = |1 - w| \iff |w| = |w - 1|$: the set of points equidistant from $0$ and $1$, i.e.\ the vertical line $\operatorname{Re}(w) = \dfrac{1}{2}$.
\end{exemplebox}

\begin{savaistu}[Möbius transformations and the Riemann sphere]
Möbius maps are the only conformal (angle-preserving) automorphisms of the Riemann sphere $\mathbb{C} \cup \{\infty\}$. They appear in physics (relativity, fluid dynamics), in cartography (stereographic projection), and in the study of fractals. The fact that they map ``circles to circles'' (where lines are circles through $\infty$) makes them indispensable in geometry.
\end{savaistu}

\begin{exercice}[Practice]
\begin{enumerate}
\item Find the image of the line $\operatorname{Im}(z) = 2$ under $w = (1 - i)z + 2i$.
\item Find the image of the circle $|z - 1| = 1$ under $w = \dfrac{1}{z}$. Why is the image a line?
\item Under $w = z^2$, find the image of the ray $\arg z = \dfrac{\pi}{6}$ and of the circle $|z| = 3$.
\item Find the image of the line $\operatorname{Re}(z) = \dfrac{1}{2}$ under $w = \dfrac{1}{z}$, giving the centre and radius of the image circle.
\end{enumerate}
\end{exercice}

\begin{corrige}[Practice]
\textbf{1.} Invert: $z = \dfrac{w - 2i}{1 - i}$. With $w = u + iv$: $z = \dfrac{u + i(v-2)}{1-i} = \dfrac{(u + i(v-2))(1+i)}{2} = \dfrac{u - (v-2)}{2} + i\,\dfrac{u + v - 2}{2}$. Then $\operatorname{Im}(z) = 2 \iff u + v - 2 = 4 \iff v = 6 - u$: the line $u + v = 6$.

\textbf{2.} $|z - 1| = 1 \iff \left|\dfrac{1}{w} - 1\right| = 1 \iff |1 - w| = |w|$: the perpendicular bisector of $0$ and $1$, i.e.\ $\operatorname{Re}(w) = \tfrac12$. The image is a line because the original circle passes through the origin.

\textbf{3.} The ray $\arg z = \pi/6$ maps to the ray $\arg w = \pi/3$ (origin excluded). The circle $|z| = 3$ maps to $|w| = 9$, described twice.

\textbf{4.} With $z = \tfrac12 + iy$: $u = \dfrac{1/2}{1/4 + y^2}$, $v = \dfrac{-y}{1/4 + y^2}$, and $u^2 + v^2 = \dfrac{1}{1/4 + y^2} = 2u$. Hence $(u - 1)^2 + v^2 = 1$: circle of centre $1$ and radius $1$, through $O$ (excluded).
\end{corrige}

\begin{aretenir}[Key points]
\begin{itemize}
\item $w = z + b$: translation; $w = az$ with $a = re^{i\alpha}$: enlargement of ratio $r$ composed with rotation of angle $\alpha$; $w = \bar z$: reflection in the real axis.
\item Under $w = az + b$: lines $\to$ lines, circles $\to$ circles (radius multiplied by $|a|$).
\item Under $w = \dfrac{1}{z}$: $|w| = \dfrac{1}{|z|}$, $\arg w = -\arg z$; the unit circle is invariant; a line avoiding $O$ maps to a circle through $O$.
\item Under $w = z^2$: moduli squared, arguments doubled; circles and lines are \emph{not} preserved in general.
\item Möbius maps $w = \dfrac{az+b}{cz+d}$ send lines and circles to lines or circles.
\item Method: \cle{parametrise} $\to$ \cle{substitute} $\to$ \cle{eliminate} $\to$ identify, stating excluded points.
\end{itemize}
\end{aretenir}

These transformations are not only geometric objects in their own right: in the next sections we shall see that the maps $w = az + b$ are exactly the \emph{direct similitudes}, and that their action on points $(x, y)$ can be written with $2 \times 2$ matrices — a bridge between complex numbers and the matrix transformations of the plane.

\coursec{Plane Transformations and Matrix Representation}

In the previous section we studied how a complex function $w = f(z)$ maps points from the $z$-plane to the $w$-plane. Many of those mappings — translations, rotations, dilations, and even the Joukowski map — are built from elementary geometric operations on the plane. In this section we place those operations on a firm algebraic footing: every \cle{linear transformation} of the plane can be represented by a $2 \times 2$ matrix acting on column vectors $\begin{pmatrix} x \\ y \end{pmatrix}$. This matrix viewpoint not only unifies the treatment of reflections, rotations, stretches and shears, but also gives us powerful tools (determinant, inverse, composition) that are indispensable in the GCE Advanced Level examination.

\courssub{Linear Transformations of the Plane and Matrix Representation}

\begin{definition}[Linear plane transformation]
A transformation $T : \mathbb{R}^2 \to \mathbb{R}^2$ is \textbf{linear} if for all vectors $\mathbf{u}, \mathbf{v} \in \mathbb{R}^2$ and all scalars $\lambda \in \mathbb{R}$:
\[
T(\mathbf{u} + \mathbf{v}) = T(\mathbf{u}) + T(\mathbf{v}), \qquad
T(\lambda \mathbf{u}) = \lambda T(\mathbf{u}).
\]
An immediate consequence is $T(\mathbf{0}) = \mathbf{0}$ (the origin is always fixed). Every linear transformation can be written as
\[
T\!\begin{pmatrix} x \\ y \end{pmatrix}
= \begin{pmatrix} a & b \\ c & d \end{pmatrix}
\begin{pmatrix} x \\ y \end{pmatrix}
= \begin{pmatrix} ax + by \\ cx + dy \end{pmatrix},
\]
where the matrix $M = \begin{pmatrix} a & b \\ c & d \end{pmatrix}$ is called the \cle{matrix of the transformation} (relative to the standard basis $\{\mathbf{e}_1, \mathbf{e}_2\}$). The columns of $M$ are precisely the images of the basis vectors:
\[
M\mathbf{e}_1 = \begin{pmatrix} a \\ c \end{pmatrix} = T(1,0), \qquad
M\mathbf{e}_2 = \begin{pmatrix} b \\ d \end{pmatrix} = T(0,1).
\]
\end{definition}

\begin{exemplebox}[Finding the matrix of a linear transformation]
Let $T$ be the transformation that maps $(x,y)$ to $(2x - y,\; x + 3y)$. Write $T$ in matrix form.
\par
The coefficients of $x$ and $y$ in the first component give the first row $(2,\,-1)$; the second component gives the second row $(1,\,3)$. Hence
\[
M = \begin{pmatrix} 2 & -1 \\ 1 & 3 \end{pmatrix}.
\]
Check: $M\begin{pmatrix} x \\ y \end{pmatrix} = \begin{pmatrix} 2x - y \\ x + 3y \end{pmatrix}$.
\end{exemplebox}

\begin{attention}[Translation is not linear]
A translation $T(x,y) = (x + a,\; y + b)$ with $(a,b) \neq (0,0)$ does \emph{not} satisfy $T(\mathbf{0}) = \mathbf{0}$, hence it is not linear and cannot be represented by a $2 \times 2$ matrix acting on $\begin{pmatrix} x \\ y \end{pmatrix}$ alone. (In homogeneous coordinates one uses a $3 \times 3$ matrix, but that is beyond the current syllabus.)
\end{attention}

\begin{savaistu}[Link with complex numbers]
A complex number $z = x+iy$ can be identified with the vector $\begin{pmatrix} x \\ y \end{pmatrix}$. Multiplication by a fixed complex number $w = a+ib$ corresponds to the linear transformation
\[
(x+iy) \mapsto (a+ib)(x+iy) = (ax-by) + i(bx+ay)
\quad\longleftrightarrow\quad
\begin{pmatrix} a & -b \\ b & a \end{pmatrix}
\begin{pmatrix} x \\ y \end{pmatrix}.
\]
Matrices of the form $\begin{pmatrix} a & -b \\ b & a \end{pmatrix}$ represent \cle{similitudes} (rotation combined with enlargement), a key concept we will develop in the next section.
\end{savaistu}

\courssub{Reflections}

A reflection is a linear transformation that flips the plane across a line through the origin. The most common reflection matrices are:

\begin{itemize}
\item Reflection in the $x$-axis: $\displaystyle M_x = \begin{pmatrix} 1 & 0 \\ 0 & -1 \end{pmatrix}$.
\item Reflection in the $y$-axis: $\displaystyle M_y = \begin{pmatrix} -1 & 0 \\ 0 & 1 \end{pmatrix}$.
\item Reflection in the line $y = x$: $\displaystyle M_{y=x} = \begin{pmatrix} 0 & 1 \\ 1 & 0 \end{pmatrix}$.
\item Reflection in the line $y = (\tan\theta)\,x$ (making an angle $\theta$ with the positive $x$-axis):
\[
M_\theta = \begin{pmatrix} \cos 2\theta & \sin 2\theta \\ \sin 2\theta & -\cos 2\theta \end{pmatrix}.
\]
\end{itemize}

\begin{experience}[Deriving the general reflection matrix]
To reflect in the line $L_\theta : y = (\tan\theta)x$:
\begin{enumerate}
\item Rotate by $-\theta$ to align $L_\theta$ with the $x$-axis: $R(-\theta)$.
\item Reflect in the $x$-axis: $M_x$.
\item Rotate back by $\theta$: $R(\theta)$.
\end{enumerate}
The composite matrix is $M_\theta = R(\theta) M_x R(-\theta)$. Using
$R(\theta) = \begin{pmatrix} \cos\theta & -\sin\theta \\ \sin\theta & \cos\theta \end{pmatrix}$ and $M_x = \begin{pmatrix} 1 & 0 \\ 0 & -1 \end{pmatrix}$, compute:
\[
M_\theta = \begin{pmatrix} \cos\theta & -\sin\theta \\ \sin\theta & \cos\theta \end{pmatrix}
\begin{pmatrix} 1 & 0 \\ 0 & -1 \end{pmatrix}
\begin{pmatrix} \cos\theta & \sin\theta \\ -\sin\theta & \cos\theta \end{pmatrix}
= \begin{pmatrix} \cos 2\theta & \sin 2\theta \\ \sin 2\theta & -\cos 2\theta \end{pmatrix}.
\]
This derivation is examinable and explains the double-angle formulas.
\end{experience}

\begin{propriete}[Properties of reflection matrices]
For any reflection matrix $M$:
\begin{enumerate}
\item $M^2 = I$ (applying the reflection twice returns the original point).
\item $\det M = -1$ (the orientation is reversed; area is preserved).
\item $M$ is symmetric ($M^T = M$) and orthogonal ($M^T M = I$).
\end{enumerate}
These properties are examinable; you may be asked to verify $M^2 = I$ or to find $\theta$ given a matrix of the form $\begin{pmatrix} a & b \\ b & -a \end{pmatrix}$ with $a^2+b^2=1$.
\end{propriete}

\begin{exemplebox}[Verifying $M^2 = I$ for reflection in $y = x$]
$M = \begin{pmatrix} 0 & 1 \\ 1 & 0 \end{pmatrix}$.
\[
M^2 = \begin{pmatrix} 0 & 1 \\ 1 & 0 \end{pmatrix}
\begin{pmatrix} 0 & 1 \\ 1 & 0 \end{pmatrix}
= \begin{pmatrix} 1 & 0 \\ 0 & 1 \end{pmatrix} = I.
\]
Also $\det M = -1$, confirming the orientation reversal.
\end{exemplebox}

\begin{popfigure}
\begin{tikzpicture}[scale=1.2, >=stealth]
% Axes
\draw[->] (-2.5,0) -- (2.5,0) node[right] {$x$};
\draw[->] (0,-2.5) -- (0,2.5) node[above] {$y$};
% Line y=x
\draw[popblue, thick, dashed] (-2,-2) -- (2,2) node[above right] {$y=x$};
% Original triangle
\draw[poporange, thick, fill=poporange!20] (1,0.5) -- (2,0.5) -- (1.5,1.5) -- cycle;
\node[poporange] at (1.6,0.8) {$T$};
% Reflected triangle
\draw[popgreen, thick, fill=popgreen!20] (0.5,1) -- (0.5,2) -- (1.5,1.5) -- cycle;
\node[popgreen] at (0.8,1.5) {$T'$};
% Matrix label
\node[align=center, text width=4.5cm] at (0,-2.8) {
Matrix of reflection in $y=x$:
$\begin{pmatrix} 0 & 1 \\ 1 & 0 \end{pmatrix}$.
};
\end{tikzpicture}
\legende{Reflection of a triangle in the line $y=x$. The matrix swaps the coordinates.}
\end{popfigure}

\courssub{Rotations about the Origin}

A rotation through an angle $\theta$ (anticlockwise) about the origin is a linear transformation with matrix
\[
R(\theta) = \begin{pmatrix} \cos\theta & -\sin\theta \\ \sin\theta & \cos\theta \end{pmatrix}.
\]

\begin{propriete}[Properties of rotation matrices]
For any rotation matrix $R(\theta)$:
\begin{enumerate}
\item $\det R(\theta) = \cos^2\theta + \sin^2\theta = 1$ (area and orientation are preserved).
\item $R(\theta)^{-1} = R(-\theta) = R(\theta)^T$ (the inverse is the transpose).
\item $R(\alpha)R(\beta) = R(\alpha+\beta)$ (rotations add angles).
\end{enumerate}
\end{propriete}

\begin{exemplebox}[Rotation by $60^\circ$]
$\theta = 60^\circ = \frac{\pi}{3}$, $\cos\frac{\pi}{3} = \frac12$, $\sin\frac{\pi}{3} = \frac{\sqrt3}{2}$.
\[
R\!\left(\frac{\pi}{3}\right) = \begin{pmatrix} \frac12 & -\frac{\sqrt3}{2} \\[2pt] \frac{\sqrt3}{2} & \frac12 \end{pmatrix}.
\]
The image of $(1,0)$ is $\bigl(\frac12,\frac{\sqrt3}{2}\bigr)$; the image of $(0,1)$ is $\bigl(-\frac{\sqrt3}{2},\frac12\bigr)$.
\end{exemplebox}

\courssub{Enlargements (Dilations) and Stretches}

An \cle{enlargement} (or dilation) centred at the origin with scale factor $k$ multiplies every vector by $k$. Its matrix is $kI = \begin{pmatrix} k & 0 \\ 0 & k \end{pmatrix}$.

A \cle{stretch} parallel to the coordinate axes uses different scale factors $a$ (in the $x$-direction) and $b$ (in the $y$-direction):
\[
S = \begin{pmatrix} a & 0 \\ 0 & b \end{pmatrix}.
\]
If $a,b > 0$ the transformation preserves orientation ($\det S = ab > 0$); if one factor is negative, a reflection is combined with the stretch. If $a=b$, the stretch reduces to an enlargement.

\begin{exemplebox}[Stretch with $a=2$, $b=1$]
$S = \begin{pmatrix} 2 & 0 \\ 0 & 1 \end{pmatrix}$. The unit square $[0,1]\times[0,1]$ becomes a $2\times1$ rectangle. Area is multiplied by $\det S = 2$.
\end{exemplebox}

\courssub{Shears}

A \cle{shear} parallel to the $x$-axis with factor $k$ fixes the $x$-axis and shifts points horizontally by an amount proportional to their $y$-coordinate:
\[
H_x(k) = \begin{pmatrix} 1 & k \\ 0 & 1 \end{pmatrix},\qquad
\begin{pmatrix} x \\ y \end{pmatrix} \mapsto \begin{pmatrix} x + ky \\ y \end{pmatrix}.
\]
A shear parallel to the $y$-axis with factor $k$ fixes the $y$-axis and shifts points vertically by an amount proportional to their $x$-coordinate:
\[
H_y(k) = \begin{pmatrix} 1 & 0 \\ k & 1 \end{pmatrix},\qquad
\begin{pmatrix} x \\ y \end{pmatrix} \mapsto \begin{pmatrix} x \\ kx + y \end{pmatrix}.
\]

\begin{propriete}[Effect of a shear]
\begin{itemize}
\item $\det H_x(k) = \det H_y(k) = 1$; area is preserved.
\item The unit square becomes a parallelogram of the same area.
\item A shear is not an isometry (distances and angles change).
\end{itemize}
\end{propriete}

\begin{popfigure}
\begin{tikzpicture}[scale=1.5, >=stealth]
% Original unit square
\draw[popdark, thick, fill=popblueL] (0,0) rectangle (1,1);
\node[popdark] at (0.5,0.5) {$S_0$};
% Rotation by 60 degrees
\begin{scope}[shift={(3,0)}]
\draw[->] (-0.5,0) -- (1.5,0) node[right] {$x$};
\draw[->] (0,-0.5) -- (0,1.5) node[above] {$y$};
\draw[poporange, thick, fill=poporange!20]
(0,0) -- (0.5,0.866) -- (-0.366,1.366) -- (-0.866,0.5) -- cycle;
\node[poporange] at (0.1,0.8) {$R_{60^\circ}(S_0)$};
\node[align=center] at (0.5,-0.8) {Rotation $60^\circ$\\ $\begin{pmatrix} 1/2 & -\sqrt3/2 \\ \sqrt3/2 & 1/2 \end{pmatrix}$};
\end{scope}
% Shear k=1 parallel to x-axis
\begin{scope}[shift={(0,-2.5)}]
\draw[->] (-0.5,0) -- (2,0) node[right] {$x$};
\draw[->] (0,-0.5) -- (0,1.5) node[above] {$y$};
\draw[popgreen, thick, fill=popgreen!20]
(0,0) -- (1,0) -- (2,1) -- (1,1) -- cycle;
\node[popgreen] at (1,0.5) {$H_x(1)(S_0)$};
\node[align=center] at (1,-0.8) {Shear $k=1$ (x-axis)\\ $\begin{pmatrix} 1 & 1 \\ 0 & 1 \end{pmatrix}$};
\end{scope}
% Stretch diag(2,1)
\begin{scope}[shift={(3,-2.5)}]
\draw[->] (-0.5,0) -- (2.5,0) node[right] {$x$};
\draw[->] (0,-0.5) -- (0,1.5) node[above] {$y$};
\draw[poppurple, thick, fill=poppurple!20]
(0,0) -- (2,0) -- (2,1) -- (0,1) -- cycle;
\node[poppurple] at (1,0.5) {$S(S_0)$};
\node[align=center] at (1,-0.8) {Stretch $\begin{pmatrix} 2 & 0 \\ 0 & 1 \end{pmatrix}$};
\end{scope}
\end{tikzpicture}
\legende{Images of the unit square under a rotation of $60^\circ$, a shear with $k=1$ (parallel to $x$-axis), and a stretch $\operatorname{diag}(2,1)$.}
\end{popfigure}

\courssub{Composite Transformations}

If a transformation $T_1$ with matrix $A$ is followed by $T_2$ with matrix $B$, the combined transformation $T = T_2 \circ T_1$ has matrix $BA$ (note the \emph{reverse order}: the matrix of the first transformation appears on the right).

\begin{methode}[Matrix of a composite transformation]
To find the matrix of the transformation ``apply $T_1$ then $T_2$'':
\begin{enumerate}
\item Write the matrices $M_1$ (for $T_1$) and $M_2$ (for $T_2$).
\item Multiply in the order $M = M_2 M_1$ (rightmost applied first).
\item The resulting matrix $M$ represents the composite transformation.
\end{enumerate}
\end{methode}

\begin{attention}[Order of multiplication]
A very common mistake is to multiply in the order $M_1 M_2$. Remember: if you apply $T_1$ then $T_2$, the matrix is $M_2 M_1$. In function notation $(T_2 \circ T_1)(\mathbf{x}) = T_2(T_1(\mathbf{x})) = M_2(M_1\mathbf{x}) = (M_2 M_1)\mathbf{x}$.
\end{attention}

\begin{exemplebox}[Rotation then shear]
Rotate by $90^\circ$ ($R = \begin{pmatrix} 0 & -1 \\ 1 & 0 \end{pmatrix}$), then apply shear $H_x(1) = \begin{pmatrix} 1 & 1 \\ 0 & 1 \end{pmatrix}$.
Composite matrix:
\[
M = H_x(1) \, R = \begin{pmatrix} 1 & 1 \\ 0 & 1 \end{pmatrix}
\begin{pmatrix} 0 & -1 \\ 1 & 0 \end{pmatrix}
= \begin{pmatrix} 1 & -1 \\ 1 & 0 \end{pmatrix}.
\]
If we reversed the order (shear then rotation) we would get
$R\,H_x(1) = \begin{pmatrix} 0 & -1 \\ 1 & 1 \end{pmatrix} \neq M$, showing non-commutativity.
\end{exemplebox}

\begin{exemplebox}[Reflection then rotation]
Reflect in the line $y=x$ ($M_{y=x} = \begin{pmatrix} 0 & 1 \\ 1 & 0 \end{pmatrix}$), then rotate by $90^\circ$ ($R_{90} = \begin{pmatrix} 0 & -1 \\ 1 & 0 \end{pmatrix}$).
Composite matrix:
\[
M = R_{90} M_{y=x} = \begin{pmatrix} 0 & -1 \\ 1 & 0 \end{pmatrix}
\begin{pmatrix} 0 & 1 \\ 1 & 0 \end{pmatrix}
= \begin{pmatrix} -1 & 0 \\ 0 & 1 \end{pmatrix} = M_y.
\]
The composite is a reflection in the $y$-axis. This illustrates that the product of a reflection (det $-1$) and a rotation (det $1$) is a reflection (det $-1$).
\end{exemplebox}

\courssub{Determinant and Area Scale Factor}

\begin{propriete}[Determinant as signed area scale factor]
Let $T$ be a linear transformation with matrix $M$. For any region $R$ in the plane with area $A$, the area of the image $T(R)$ is $|\det M| \cdot A$. The sign of $\det M$ indicates whether orientation is preserved ($+$) or reversed ($-$).
\end{propriete}

\begin{aretenir}[Key facts about determinants]
\begin{itemize}
\item $\det(kI) = k^2$ (enlargement by $k$ multiplies area by $k^2$).
\item $\det\begin{pmatrix} a & 0 \\ 0 & b \end{pmatrix} = ab$ (stretch).
\item $\det H_x(k) = \det H_y(k) = 1$ (shear preserves area).
\item $\det R(\theta) = 1$ (rotation preserves area and orientation).
\item $\det M_{\text{reflection}} = -1$ (reflection preserves area but reverses orientation).
\item If $\det M = 0$, the transformation is \cle{singular}: it collapses the plane onto a line or a point (e.g., projection onto an axis). Such a transformation has no inverse.
\end{itemize}
\end{aretenir}

\courssub{Inverse Transformations and Identifying a Transformation from its Matrix}

If $\det M \neq 0$, the transformation is invertible and its inverse has matrix $M^{-1} = \frac{1}{\det M}\begin{pmatrix} d & -b \\ -c & a \end{pmatrix}$ for $M = \begin{pmatrix} a & b \\ c & d \end{pmatrix}$.

\begin{methode}[Identifying a transformation from its matrix (exam technique)]
Given a $2\times2$ matrix $M = \begin{pmatrix} a & b \\ c & d \end{pmatrix}$:
\begin{enumerate}
\item Compute $\det M$.
\item If $\det M = 0$: singular transformation (projection or collapse).
\item If $\det M > 0$ and $M^T M = (\det M)I$ (i.e. columns are orthogonal and have same length $\sqrt{\det M}$): it is a \textbf{similitude} (rotation + enlargement). The scale factor is $\sqrt{\det M}$, and the rotation angle $\theta$ satisfies $\cos\theta = \frac{a}{\sqrt{\det M}}$, $\sin\theta = \frac{c}{\sqrt{\det M}}$.
\item If $\det M = 1$ and $M^T M = I$: it is a \textbf{rotation} (find $\theta$ from $\cos\theta = a$, $\sin\theta = c$).
\item If $\det M = -1$ and $M^T M = I$: it is a \textbf{reflection} (find the line of fixed points: solve $M\mathbf{x} = \mathbf{x}$).
\item If $M$ is diagonal $\begin{pmatrix} a & 0 \\ 0 & b \end{pmatrix}$: it is a \textbf{stretch} (or enlargement if $a=b$).
\item If $M = \begin{pmatrix} 1 & k \\ 0 & 1 \end{pmatrix}$ or $\begin{pmatrix} 1 & 0 \\ k & 1 \end{pmatrix}$: it is a \textbf{shear}.
\item If $M$ is a product of the above, decompose it (e.g., by recognising standard forms or using polar decomposition) — but in exams you usually recognise standard forms directly.
\end{enumerate}
\end{methode}

\begin{exoresolu}[Identifying a transformation]
The matrix $M = \begin{pmatrix} 0 & -1 \\ -1 & 0 \end{pmatrix}$ represents a single geometric transformation. Identify it fully.
\par
\tcblower
\textbf{Solution:}
\begin{enumerate}
\item $\det M = (0)(0) - (-1)(-1) = -1$. Orientation is reversed.
\item $M^T M = \begin{pmatrix} 0 & -1 \\ -1 & 0 \end{pmatrix}\begin{pmatrix} 0 & -1 \\ -1 & 0 \end{pmatrix} = \begin{pmatrix} 1 & 0 \\ 0 & 1 \end{pmatrix} = I$. So $M$ is orthogonal.
\item An orthogonal matrix with determinant $-1$ is a reflection.
\item The line of reflection is the set of fixed points: $M\begin{pmatrix} x \\ y \end{pmatrix} = \begin{pmatrix} x \\ y \end{pmatrix} \Rightarrow \begin{pmatrix} -y \\ -x \end{pmatrix} = \begin{pmatrix} x \\ y \end{pmatrix} \Rightarrow x = -y$. Hence the line is $y = -x$ (or $\theta = -45^\circ$).
\end{enumerate}
\textbf{Answer:} Reflection in the line $y = -x$.
\end{exoresolu}

\begin{exoresolu}[Identifying a similitude]
Identify the transformation represented by $M = \begin{pmatrix} 2 & -2 \\ 2 & 2 \end{pmatrix}$.
\par
\tcblower
\textbf{Solution:}
\begin{enumerate}
\item $\det M = 4 - (-4) = 8 > 0$.
\item $M^T M = \begin{pmatrix} 2 & 2 \\ -2 & 2 \end{pmatrix}\begin{pmatrix} 2 & -2 \\ 2 & 2 \end{pmatrix} = \begin{pmatrix} 8 & 0 \\ 0 & 8 \end{pmatrix} = 8I = (\det M)I$.
\item Columns are orthogonal and have length $\sqrt{8} = 2\sqrt{2}$.
\item Scale factor $k = \sqrt{\det M} = 2\sqrt{2}$.
\item Rotation angle: $\cos\theta = \frac{2}{2\sqrt{2}} = \frac{1}{\sqrt{2}}$, $\sin\theta = \frac{2}{2\sqrt{2}} = \frac{1}{\sqrt{2}} \Rightarrow \theta = 45^\circ$.
\end{enumerate}
\textbf{Answer:} Enlargement scale factor $2\sqrt{2}$ followed by rotation $45^\circ$ anticlockwise (or equivalently, multiplication by $2\sqrt{2}e^{i\pi/4} = 2+2i$ in the complex plane).
\end{exoresolu}

\begin{exercice}[Practice]
\begin{enumerate}
\item Find the matrix of the transformation that reflects in the line $y = \sqrt3\,x$ (i.e. $\theta = 60^\circ$). Verify that $M^2 = I$ and $\det M = -1$.
\item A triangle has vertices $A(1,2)$, $B(3,1)$, $C(2,4)$. It undergoes a shear $H_y(2)$ followed by a rotation of $-90^\circ$. Find the coordinates of the image vertices.
\item The matrix $M = \begin{pmatrix} 3 & 4 \\ 4 & -3 \end{pmatrix}$ represents a reflection combined with an enlargement. Determine the scale factor of the enlargement and the line of reflection.
\item Show that the transformation with matrix $\begin{pmatrix} 2 & 1 \\ 4 & 2 \end{pmatrix}$ is singular. Describe its geometric effect.
\item Find the matrix of the composite transformation: stretch $S = \begin{pmatrix} 2 & 0 \\ 0 & 3 \end{pmatrix}$ followed by rotation $R(90^\circ)$. Is the result a similitude? Justify.
\end{enumerate}
\end{exercice}

\begin{corrige}[Exercise 1]
The line makes angle $\theta = 60^\circ$, so $2\theta = 120^\circ$. $\cos120^\circ = -\frac12$, $\sin120^\circ = \frac{\sqrt3}{2}$.
$M = \begin{pmatrix} -\frac12 & \frac{\sqrt3}{2} \\ \frac{\sqrt3}{2} & \frac12 \end{pmatrix}$.
$M^2 = \begin{pmatrix} \frac14+\frac34 & -\frac{\sqrt3}{4}+\frac{\sqrt3}{4} \\ -\frac{\sqrt3}{4}+\frac{\sqrt3}{4} & \frac34+\frac14 \end{pmatrix} = I$, $\det M = -\frac14 - \frac34 = -1$.
\end{corrige}

\begin{corrige}[Exercise 2]
$H_y(2) = \begin{pmatrix} 1 & 0 \\ 2 & 1 \end{pmatrix}$, $R(-90^\circ) = \begin{pmatrix} 0 & 1 \\ -1 & 0 \end{pmatrix}$.
Composite $M = R(-90^\circ) H_y(2) = \begin{pmatrix} 0 & 1 \\ -1 & 0 \end{pmatrix}\begin{pmatrix} 1 & 0 \\ 2 & 1 \end{pmatrix} = \begin{pmatrix} 2 & 1 \\ -1 & 0 \end{pmatrix}$.
Apply to each vertex:
$A' = M\binom{1}{2} = \binom{2\cdot1+2}{-1} = \binom{4}{-1}$,
$B' = M\binom{3}{1} = \binom{6+1}{-3} = \binom{7}{-3}$,
$C' = M\binom{2}{4} = \binom{4+4}{-2} = \binom{8}{-2}$.
\end{corrige}

\begin{corrige}[Exercise 3]
$\det M = -9 - 16 = -25 = -5^2$. The matrix is $5$ times an orthogonal matrix:
$M = 5 \begin{pmatrix} 3/5 & 4/5 \\ 4/5 & -3/5 \end{pmatrix}$.
The orthogonal part has determinant $-1$, so it is a reflection.
The line of reflection satisfies $\begin{pmatrix} 3/5 & 4/5 \\ 4/5 & -3/5 \end{pmatrix}\binom{x}{y} = \binom{x}{y} \Rightarrow \frac{3}{5}x + \frac{4}{5}y = x \Rightarrow 4y = 2x \Rightarrow y = \frac12 x$.
Enlargement scale factor $= 5$.
\end{corrige}

\begin{corrige}[Exercise 4]
$\det = 4 - 4 = 0$, singular. The matrix has rank 1; its columns are $\binom{2}{4}$ and $\binom{1}{2}$, which are linearly dependent (second is half the first). The transformation maps the whole plane onto the line spanned by $\binom{1}{2}$ (equation $y=2x$). It is a projection onto that line combined with a scaling (the image of $(x,y)$ is $(2x+y)\binom{1}{2}$).
\end{corrige}

\begin{corrige}[Exercise 5]
$R(90^\circ) = \begin{pmatrix} 0 & -1 \\ 1 & 0 \end{pmatrix}$.
Composite $M = R(90^\circ) S = \begin{pmatrix} 0 & -1 \\ 1 & 0 \end{pmatrix}\begin{pmatrix} 2 & 0 \\ 0 & 3 \end{pmatrix} = \begin{pmatrix} 0 & -3 \\ 2 & 0 \end{pmatrix}$.
$\det M = 0 - (-6) = 6 > 0$.
$M^T M = \begin{pmatrix} 0 & 2 \\ -3 & 0 \end{pmatrix}\begin{pmatrix} 0 & -3 \\ 2 & 0 \end{pmatrix} = \begin{pmatrix} 4 & 0 \\ 0 & 9 \end{pmatrix} \neq 6I$.
The columns are orthogonal but have different lengths ($2$ and $3$). Hence $M^T M \neq (\det M)I$.
\textbf{Answer:} The result is \emph{not} a similitude; it is a rotation combined with a stretch (different scale factors in perpendicular directions), which distorts angles.
\end{corrige}

\coursec{Similitudes: Complex Form, Matrix Interpretation and Applications}

In the previous section we studied plane transformations through matrices: rotations, reflections, enlargements, stretches and shears, all described by $2\times 2$ matrices acting on position vectors. We now single out a particularly beautiful family of transformations, the \cle{similitudes}, which combine a rotation, an enlargement and possibly a translation or reflection into a single geometric object. The extraordinary fact is that complex numbers describe these transformations with a single formula: $z' = az + b$ (or $z' = a\bar{z} + b$). This is one of the most elegant meetings of algebra and geometry in the whole Further Mathematics syllabus, and it is a favourite of GCE Advanced Level examiners.

\courssub{Definition and First Examples}

\begin{experience}[Growing a map without distorting it]
Imagine a photocopier set to $150\%$. The map of Cameroon it produces is bigger, but every angle is unchanged: the angle between two roads is the same on the copy as on the original, and the ratio of two distances is preserved. A transformation with this behaviour --- all distances multiplied by the same positive constant --- is exactly what we call a similitude.
\end{experience}

\begin{definition}[Similitude and ratio]
A \cle{similitude} of the plane is a transformation $s$ for which there exists a fixed real number $k > 0$, called the \cle{ratio} of the similitude, such that for all points $M, N$ with images $M' = s(M)$, $N' = s(N)$,
\[ M'N' = k \cdot MN. \]
If $k = 1$ the similitude is an \cle{isometry} (it preserves distances). If $k > 1$ it enlarges; if $0 < k < 1$ it contracts.
\end{definition}

Rotations, translations and reflections (ratio $1$), and homotheties (enlargements) of ratio $k$, are all similitudes. The new phenomenon is that we may \emph{combine} a rotation and an enlargement about the same centre: this is still a similitude, and it is the typical case.

\courssub{Complex Form of a Similitude}

\begin{propriete}[Complex form of similitudes]
Every similitude of the plane has one of the following two complex forms:
\begin{itemize}
\item \textbf{Direct similitude} (preserves orientation): $\;z' = az + b$, with $a, b \in \mathbb{C}$, $a \neq 0$;
\item \textbf{Indirect (opposite) similitude} (reverses orientation): $\;z' = a\bar{z} + b$, with $a, b \in \mathbb{C}$, $a \neq 0$.
\end{itemize}
Conversely, every map of one of these forms is a similitude of ratio $k = |a|$. (The derivation below is examinable in spirit: you must be able to justify that $z' = az+b$ multiplies distances by $|a|$.)
\end{propriete}

\textbf{Justification.} Let $z_1, z_2$ be the affixes of $M, N$ and $z_1', z_2'$ those of their images.
\begin{itemize}
\item For a direct similitude $z' = az + b$:
\[ |z_1' - z_2'| = |(az_1 + b) - (az_2 + b)| = |a(z_1 - z_2)| = |a|\,|z_1 - z_2|, \]
so $M'N' = |a|\cdot MN$.
\item For an indirect similitude $z' = a\bar{z} + b$:
\[ |z_1' - z_2'| = |a\bar{z}_1 + b - (a\bar{z}_2 + b)| = |a(\bar{z}_1 - \bar{z}_2)| = |a|\,|\bar{z}_1 - \bar{z}_2| = |a|\,|z_1 - z_2|, \]
since conjugation preserves moduli ($|\bar{z}| = |z|$).
\end{itemize}
In both cases all distances are multiplied by $k = |a|$. $\blacksquare$

\begin{aretenir}[Reading a direct similitude]
For the direct similitude $z' = az + b$, write $a = re^{i\theta}$ ($r>0$). Then:
\begin{itemize}
\item the \cle{ratio} is $k = |a| = r$;
\item the \cle{angle} of the similitude is $\theta = \arg(a)$ (defined modulo $2\pi$);
\item geometrically, $s$ is the composite of the homothety of ratio $r$ and the rotation of angle $\theta$, \emph{about the same centre}.
\end{itemize}
\end{aretenir}

\courssub{Classification of Direct Similitudes}

The nature of $z' = az + b$ depends entirely on $a$:

\begin{center}
\begin{tabularx}{\linewidth}{|l|X|X|}
\hline
\rowcolor{popblueL} \textbf{Condition on $a$} & \textbf{Transformation} & \textbf{Fixed points} \\
\hline
$a = 1$, $b \neq 0$ & Translation by vector of affix $b$ & none \\
\hline
$a = 1$, $b = 0$ & Identity & all points \\
\hline
$|a| = 1$, $a \neq 1$ & Rotation of angle $\arg a$ & one (the centre) \\
\hline
$a \in \mathbb{R}$, $a > 0$, $a \neq 1$ & Homothety (enlargement) of ratio $a$ & one (the centre) \\
\hline
general $a \neq 1$ & Rotation of angle $\arg a$ \textbf{and} homothety of ratio $|a|$ about the same centre & one (the centre) \\
\hline
\end{tabularx}
\end{center}

\begin{methode}[Centre, ratio and angle of $z' = az + b$ ($a \neq 1$)]
\begin{enumerate}
\item \textbf{Ratio:} compute $k = |a|$.
\item \textbf{Angle:} compute $\theta = \arg a$ (modulo $2\pi$).
\item \textbf{Centre:} the centre $\Omega$ is the unique \cle{fixed point}; solve
\[ z = az + b \quad\Longleftrightarrow\quad z_0 = \frac{b}{1 - a}. \]
\item \textbf{Conclusion:} state ``$s$ is the similitude of centre $\Omega(z_0)$, ratio $k$ and angle $\theta$'', i.e. the composite of the homothety $h(\Omega, k)$ and the rotation $r(\Omega, \theta)$ (these two commute).
\end{enumerate}
\end{methode}

\begin{attention}[The centre is not the origin!]
The most frequent error is to say that $z' = az + b$ is ``a rotation about the origin''. The rotation part acts \emph{about the centre $\Omega$}, the fixed point obtained by solving $z = az + b$. Only when $b = 0$ is the centre the origin. Always compute $z_0 = \dfrac{b}{1-a}$ first.
\end{attention}

\begin{exoresolu}[Identifying a similitude]
Characterise the transformation $s:\ z' = (1 + i)z + (2 - i)$.
\tcblower
\textbf{Solution.} Here $a = 1 + i$ and $b = 2 - i$.

\textbf{Ratio:} $k = |a| = \sqrt{1^2 + 1^2} = \sqrt{2}$.

\textbf{Angle:} $\theta = \arg(1+i) = \dfrac{\pi}{4}$.

\textbf{Centre:} solve $z = (1+i)z + 2 - i$:
\[ z - (1+i)z = 2 - i \;\Longleftrightarrow\; -iz = 2 - i \;\Longleftrightarrow\; z = \frac{2-i}{-i} = \frac{(2-i)\,i}{-i \cdot i} = \frac{2i + 1}{1} = 1 + 2i. \]
\textbf{Conclusion:} $s$ is the direct similitude of centre $\Omega(1 + 2i)$, ratio $\sqrt{2}$ and angle $\dfrac{\pi}{4}$: the composite of the rotation of centre $\Omega$ and angle $\pi/4$ with the homothety of centre $\Omega$ and ratio $\sqrt{2}$.
\end{exoresolu}

\begin{popfigure}
\begin{center}
\begin{tikzpicture}[scale=1.15, >=stealth]
\coordinate (O) at (0,0);
\coordinate (A) at (1.6,0.4);
\coordinate (B) at (2.6,0.9);
\coordinate (C) at (1.9,1.5);
% images: rotate by 60 deg and scale by 2 about O
\coordinate (A2) at ($(O)!2!60:(A)$);
\coordinate (B2) at ($(O)!2!60:(B)$);
\coordinate (C2) at ($(O)!2!60:(C)$);
\fill[popink] (O) circle (1.6pt) node[below left] {$\Omega$};
\draw[popblue, thick] (A)--(B)--(C)--cycle;
\foreach \p/\n in {A/A, B/B, C/C}{\fill[popblue] (\p) circle (1.3pt) node[below] {$\n$};}
\draw[poporange, thick] (A2)--(B2)--(C2)--cycle;
\foreach \p/\n in {A2/A', B2/B', C2/C'}{\fill[poporange] (\p) circle (1.3pt) node[above left] {$\n$};}
\draw[popmuted, dashed] (O)--(A2);
\draw[popmuted, dashed] (O)--(B2);
\draw[popmuted, dashed] (O)--(C2);
\draw[popgreen, ->] (1.15,0.29) arc (14:74:1.19) node[midway, right] {$\frac{\pi}{3}$};
\node[popdark] at (3.4,2.6) {$\Omega A' = 2\,\Omega A$};
\end{tikzpicture}
\end{center}
\legende{Direct similitude of centre $\Omega$, ratio $2$ and angle $\pi/3$: triangle $ABC$ is mapped to the similar triangle $A'B'C'$, twice as large and rotated by $60^\circ$ about $\Omega$.}
\end{popfigure}

\courssub{Indirect Similitudes}

An \cle{indirect similitude} has the form $z' = a\bar{z} + b$. The conjugation $z \mapsto \bar{z}$ is the reflection in the real axis; composing it with a direct similitude reverses orientation.

\begin{exemplebox}[Special cases of indirect similitudes]
\begin{itemize}
\item $z' = \bar{z}$: reflection in the real axis.
\item $z' = e^{i\theta}\bar{z}$: reflection in the line through $O$ making angle $\theta/2$ with the real axis.
\item $z' = \bar{z} + b$ with $b$ real and non-zero: \cle{glide reflection} --- reflection in the real axis followed by a translation parallel to it.
\item $z' = a\bar{z} + b$ with $|a| \neq 1$: a reflection (or glide reflection) combined with an enlargement of ratio $|a|$.
\end{itemize}
\end{exemplebox}

\begin{attention}[Direct or indirect?]
Do not confuse the two families. \textbf{Direct} similitudes $z' = az + b$ \emph{preserve} orientation: a triangle labelled anticlockwise keeps its anticlockwise labelling. \textbf{Indirect} similitudes $z' = a\bar{z} + b$ \emph{reverse} orientation, because of the conjugation. In an examination, the presence of $\bar{z}$ in the formula is the unmistakable signature of an indirect similitude.
\end{attention}

\courssub{Matrix Interpretation}

Let $a = p + iq = re^{i\theta}$ and write $z = x + iy$, $z' = x' + iy'$. Expanding $z' = az + b$ with $b = b_1 + ib_2$:
\[ x' + iy' = (p + iq)(x + iy) + b_1 + ib_2 = (px - qy + b_1) + i(qx + py + b_2). \]
Hence, in matrix form,
\[ \begin{pmatrix} x' \\ y' \end{pmatrix} = \begin{pmatrix} p & -q \\ q & p \end{pmatrix} \begin{pmatrix} x \\ y \end{pmatrix} + \begin{pmatrix} b_1 \\ b_2 \end{pmatrix}. \]

\begin{propriete}[Similarity matrices]
A matrix of the form
\[ M = \begin{pmatrix} p & -q \\ q & p \end{pmatrix}, \qquad (p,q) \neq (0,0), \]
is called a \cle{similarity matrix}. Writing $r = \sqrt{p^2+q^2}$ and $\theta = \arg(p+iq)$, it factorises as
\[ M = r\begin{pmatrix} \cos\theta & -\sin\theta \\ \sin\theta & \cos\theta \end{pmatrix} = rR(\theta), \]
i.e. a rotation matrix scaled by $r$. Thus $z' = az + b$ corresponds exactly to the matrix map $\vec{x} \mapsto rR(\theta)\vec{x} + \vec{b}$: a rotation, an enlargement of ratio $r$, then a translation. Note $\det M = p^2 + q^2 = r^2 > 0$, confirming that direct similitudes preserve orientation.

For an indirect similitude $z' = a\bar{z} + b$ with $a = p+iq$, we have
\[ x' + iy' = (p+iq)(x-iy) + b_1 + ib_2 = (px + qy + b_1) + i(qx - py + b_2), \]
giving the matrix form
\[ \begin{pmatrix} x' \\ y' \end{pmatrix} = \begin{pmatrix} p & q \\ q & -p \end{pmatrix} \begin{pmatrix} x \\ y \end{pmatrix} + \begin{pmatrix} b_1 \\ b_2 \end{pmatrix}. \]
The determinant is $-p^2 - q^2 = -r^2 < 0$, confirming orientation reversal.
\end{propriete}

\begin{exemplebox}[From complex form to matrix]
For $s:\ z' = (1+i)z + (2-i)$, we have $p = q = 1$, so
\[ \begin{pmatrix} x' \\ y' \end{pmatrix} = \begin{pmatrix} 1 & -1 \\ 1 & 1 \end{pmatrix}\begin{pmatrix} x \\ y \end{pmatrix} + \begin{pmatrix} 2 \\ -1 \end{pmatrix} = \sqrt{2}\,R\!\left(\tfrac{\pi}{4}\right)\begin{pmatrix} x \\ y \end{pmatrix} + \begin{pmatrix} 2 \\ -1 \end{pmatrix}, \]
which matches the ratio $\sqrt{2}$ and angle $\pi/4$ found earlier.
\end{exemplebox}

\courssub{Geometric Properties of Similitudes}

\begin{propriete}[Conservation properties]
Every similitude of ratio $k$:
\begin{enumerate}
\item multiplies all distances by $k$ and all areas by $k^2$;
\item \textbf{preserves angles} (in size; direct similitudes also preserve their orientation);
\item preserves \textbf{ratios of distances} and division of a segment (in particular it maps the midpoint of $[MN]$ to the midpoint of $[M'N']$);
\item maps \textbf{lines to lines} and \textbf{circles to circles} (a line may be viewed as a circle of infinite radius);
\item maps any figure to a \textbf{similar} figure: triangles are mapped to similar triangles.
\end{enumerate}
(Proofs of (1)--(3) follow directly from $|z_1' - z_2'| = k|z_1 - z_2|$ and are examinable.)
\end{propriete}

\textbf{Proof that midpoints are preserved.} If $M, N$ have affixes $z_1, z_2$, the midpoint $I$ has affix $\frac{z_1+z_2}{2}$. Its image under $z' = az+b$ has affix
\[ a\cdot\frac{z_1+z_2}{2} + b = \frac{(az_1+b) + (az_2+b)}{2} = \frac{z_1' + z_2'}{2}, \]
which is the affix of the midpoint of $[M'N']$. $\blacksquare$

\textbf{Proof that circles map to circles.} A circle with centre $z_c$ and radius $R$ has equation $|z - z_c| = R$. Under $z' = az + b$ ($a \neq 0$), the image satisfies $|z' - (az_c+b)| = |a|\,|z - z_c| = |a|R$, which is a circle of centre $az_c+b$ and radius $|a|R$. For $z' = a\bar{z}+b$, the equation becomes $|z' - (a\bar{z}_c+b)| = |a|R$, also a circle. $\blacksquare$

\begin{popfigure}
\begin{center}
\begin{tikzpicture}[scale=2.0, >=stealth]
\coordinate (P0) at (1,0);
\coordinate (P1) at (0.7071,0.7071);
\coordinate (P2) at (0,1);
\coordinate (P3) at (-0.7071,0.7071);
\coordinate (P4) at (-1,0);
\coordinate (P5) at (-0.7071,-0.7071);
\draw[popblue, thick] (P0)--(P1);
\draw[poporange, thick] (P1)--(P2);
\draw[popgreen, thick] (P2)--(P3);
\draw[poppurple, thick] (P3)--(P4);
\draw[poppink, thick] (P4)--(P5);
\foreach \p in {P0,P1,P2,P3,P4,P5}{\fill[popdark] (\p) circle (0.9pt);}
\fill[popink] (0,0) circle (1.1pt) node[below right] {$O$};
\node[popmuted] at (1.12,-0.12) {$M_0$};
\node[popmuted] at (0.83,0.83) {$M_1$};
\node[popmuted] at (0.05,1.12) {$M_2$};
\node[popmuted] at (-0.86,0.82) {$M_3$};
\node[popmuted] at (-1.14,-0.1) {$M_4$};
\end{tikzpicture}
\end{center}
\legende{Spiral similarity: successive images of the segment $[OM_0]$ under $z' = \dfrac{1+i}{\sqrt{2}}\,z$ (ratio $1$, angle $\pi/4$) trace out a rotational chain; with ratio $k \neq 1$ the points would spiral inwards ($k<1$) or outwards ($k>1$) --- hence the name \emph{spiral similarity}.}
\end{popfigure}

\courssub{Composition of Similitudes and Applications}

\begin{propriete}[Composition of similitudes]
The composite of two similitudes of ratios $k_1$ and $k_2$ is a similitude of ratio $k_1 k_2$. For direct similitudes $s_1: z' = a_1 z + b_1$ and $s_2: z' = a_2 z + b_2$,
\[ s_2 \circ s_1:\ z' = a_2(a_1 z + b_1) + b_2 = a_1 a_2\, z + (a_2 b_1 + b_2), \]
so the angles add: $\arg(a_1 a_2) = \arg a_1 + \arg a_2$. In particular, the composite of two rotations is a rotation (angles add) or a translation (angles sum to $0$ modulo $2\pi$), and the composite of two homotheties is a homothety or a translation.
\end{propriete}

\begin{exoresolu}[Composing two similitudes]
Let $s_1: z' = 2iz + 1$ and $s_2: z' = (1 - i)z - 2$. Determine the nature and characteristic elements of $s = s_2 \circ s_1$.
\tcblower
\textbf{Solution.} Compute the composite:
\[ s(z) = (1-i)(2iz + 1) - 2 = 2i(1-i)z + (1 - i) - 2 = (2i + 2)z - 1 - i. \]
So $s:\ z' = (2 + 2i)z - 1 - i$, a direct similitude with $a = 2 + 2i$.

\textbf{Ratio:} $k = |2+2i| = 2\sqrt{2}$ (check: $|2i|\cdot|1-i| = 2\sqrt{2}$ \checkmark).

\textbf{Angle:} $\theta = \arg(2+2i) = \dfrac{\pi}{4}$ (check: $\arg(2i) + \arg(1-i) = \tfrac{\pi}{2} - \tfrac{\pi}{4} = \tfrac{\pi}{4}$ \checkmark).

\textbf{Centre:} $z_0 = \dfrac{b}{1-a} = \dfrac{-1-i}{1 - 2 - 2i} = \dfrac{-1-i}{-1-2i} = \dfrac{(-1-i)(-1+2i)}{(-1)^2 + 2^2} = \dfrac{1 - 2i + i + 2}{5} = \dfrac{3 - i}{5}$.

\textbf{Conclusion:} $s$ is the direct similitude of centre $\Omega\left(\dfrac{3-i}{5}\right)$, ratio $2\sqrt{2}$ and angle $\dfrac{\pi}{4}$.
\end{exoresolu}

\begin{exoresolu}[Writing the complex expression from geometric data]
Find the complex form of the direct similitude $s$ of centre $\Omega(2+i)$, ratio $\dfrac{1}{2}$ and angle $-\dfrac{\pi}{2}$, and determine the image of $A(3+3i)$.
\tcblower
\textbf{Solution.} We need $a$ with $|a| = \tfrac12$ and $\arg a = -\tfrac{\pi}{2}$: $a = \tfrac12 e^{-i\pi/2} = -\tfrac{i}{2}$.

The centre is the fixed point: $z_0 = \dfrac{b}{1-a}$, so $b = z_0(1 - a) = (2+i)\left(1 + \tfrac{i}{2}\right) = 2 + i + i + \tfrac{i^2}{2} = \dfrac{3}{2} + 2i$.

Hence $s:\ z' = -\dfrac{i}{2}z + \dfrac{3}{2} + 2i$.

Image of $A$: $z_A' = -\dfrac{i}{2}(3+3i) + \dfrac{3}{2} + 2i = -\dfrac{3i}{2} + \dfrac{3}{2} + \dfrac{3}{2} + 2i = 3 + \dfrac{i}{2}$. So $A'\left(3 + \dfrac{i}{2}\right)$.
\end{exoresolu}

\begin{exemplebox}[Proving triangles similar]
To prove that triangles $ABC$ and $DEF$ are directly similar, exhibit a direct similitude mapping one to the other: find $a, b$ with $z_D = az_A + b$ and $z_E = az_B + b$ (two equations, two unknowns), then check that $z_F = az_C + b$. Since similitudes preserve angles and ratios of distances, the triangles are similar with ratio $|a|$.
\end{exemplebox}

\begin{attention}[Examination traps]
\begin{itemize}
\item Solving $z = az + b$ when $a = 1$: there is no fixed point (unless $b = 0$) --- the map is a translation, not a similitude with centre.
\item The ratio is $|a|$, never $a$ itself: for $a = -3$ the ratio is $3$ (and the angle is $\pi$).
\item When composing, the order matters for the centre (though not for the ratio and angle): $s_2 \circ s_1$ means apply $s_1$ \emph{first}.
\item For indirect similitudes $z' = a\bar{z} + b$, fixed points satisfy $z = a\bar{z} + b$, which generally gives a whole \emph{line} of fixed points (reflection) or none (glide reflection) when $|a| = 1$.
\end{itemize}
\end{attention}

\begin{savaistu}[Why ``spiral similarity''?]
A direct similitude with $k \neq 1$ and $\theta \neq 0$ sends any point $M$ to $M'$ such that triangle $\Omega M M'$ always has the same shape: the angle at $\Omega$ is $\theta$ and $\dfrac{\Omega M'}{\Omega M} = k$. Iterating the map makes the images of a point spiral around $\Omega$ along a \emph{logarithmic spiral} --- the curve seen in snail shells and sunflower heads. Similitudes are also the mathematical engine behind fractals: each piece of a fractal is the image of the whole under a collection of contracting similitudes.
\end{savaistu}

\begin{exercice}[Practice]
\begin{enumerate}
\item Characterise (centre, ratio, angle) each similitude: \textbf{(a)} $z' = 3z - 4 + 2i$; \textbf{(b)} $z' = iz + 1 + i$; \textbf{(c)} $z' = (-1 + i\sqrt{3})z$.
\item Write the complex form of the similitude of centre $\Omega(1 - i)$, ratio $2$, angle $\dfrac{\pi}{3}$, and find the image of $B(2 + i)$.
\item Let $s_1: z' = (1+i)z$ and $s_2: z' = \dfrac{1-i}{2}z + 3$. Determine the nature of $s_2 \circ s_1$.
\item Show that $t: z' = 2\bar{z} + 3i$ is an indirect similitude and give its ratio.
\end{enumerate}
\end{exercice}

\begin{corrige}[Exercise 1]
\textbf{(a)} $a = 3$ real positive: homothety of ratio $3$, centre $z_0 = \dfrac{-4+2i}{1-3} = 2 - i$.
\textbf{(b)} $a = i$: $|a| = 1$, $\arg a = \pi/2$: rotation of angle $\pi/2$, centre $z_0 = \dfrac{1+i}{1-i} = \dfrac{(1+i)^2}{2} = i$.
\textbf{(c)} $a = -1 + i\sqrt{3} = 2e^{2i\pi/3}$, $b = 0$: centre $O$, ratio $2$, angle $\dfrac{2\pi}{3}$.
\end{corrige}

\begin{corrige}[Exercise 2]
$a = 2e^{i\pi/3} = 1 + i\sqrt{3}$; $b = z_0(1-a) = (1-i)(-i\sqrt{3}) = -\sqrt{3} - i\sqrt{3}$. So $z' = (1+i\sqrt{3})z - \sqrt{3}(1 + i)$. Image of $B$: $z_B' = (1+i\sqrt3)(2+i) - \sqrt3 - i\sqrt3 = 2 + i + 2i\sqrt3 - \sqrt3 - \sqrt3 - i\sqrt3 = (2 - 2\sqrt3) + i(1 + \sqrt3)$.
\end{corrige}

\begin{corrige}[Exercise 3]
$s_2 \circ s_1: z' = \dfrac{1-i}{2}(1+i)z + 3 = \dfrac{2}{2}z + 3 = z + 3$. Since $a = 1$, this is the \textbf{translation} by the vector of affix $3$ (the rotation angles $\pi/4$ and $-\pi/4$ cancel, and the ratios $\sqrt2$ and $1/\sqrt2$ multiply to $1$).
\end{corrige}

\begin{corrige}[Exercise 4]
The formula contains $\bar{z}$, so $t$ is an indirect similitude; its ratio is $k = |a| = |2| = 2$. It is the composite of a reflection-type map with an enlargement of ratio $2$.
\end{corrige}

\begin{aretenir}[Key points of the section]
\begin{itemize}
\item A similitude multiplies all distances by a fixed ratio $k > 0$.
\item Direct: $z' = az + b$; indirect: $z' = a\bar{z} + b$; in both cases $k = |a|$.
\item For $a \neq 1$, the centre is the fixed point $z_0 = \dfrac{b}{1-a}$; the angle is $\arg a$.
\item Matrix form: $\begin{pmatrix} p & -q \\ q & p \end{pmatrix} = rR(\theta)$ plus a translation for direct; $\begin{pmatrix} p & q \\ q & -p \end{pmatrix}$ for indirect.
\item Similitudes preserve angles, ratios, midpoints; they send lines to lines and circles to circles.
\item Ratios multiply and angles add under composition.
\end{itemize}
\end{aretenir}

\coursec{Modulus Inequalities and Rigid Motions}

In the previous section we studied similitudes, which combine a rotation, a scaling and a translation.  If we remove the scaling factor (i.e. set the ratio to 1) we obtain the \emph{rigid motions} or \emph{isometries} of the plane — transformations that preserve distances.  Before classifying them we first need a solid tool for estimating and comparing moduli: the triangle inequality and its consequences.  These inequalities are also the key to describing regions in the Argand diagram defined by conditions such as $|z-a|<r$ or $|z-a|\ge |z-b|$.

\courssub{The Triangle Inequality and its Consequences}

\begin{popfigure}
\begin{tikzpicture}[scale=1.2,>=stealth]
\coordinate (O) at (0,0);
\coordinate (A) at (2.5,0.8);
\coordinate (B) at (4.5,1.8);
\draw[->,thick,popblue] (O) -- (A) node[midway,above left] {$z_1$};
\draw[->,thick,poporange] (A) -- (B) node[midway,above right] {$z_2$};
\draw[->,thick,popgreen] (O) -- (B) node[midway,below] {$z_1+z_2$};
\draw[dashed,popmuted] (O) -- (2.5,0) -- (A);
\draw[dashed,popmuted] (A) -- (4.5,0.8) -- (B);
\node at (O) [below left] {$0$};
\end{tikzpicture}
\legende{Geometric proof of the triangle inequality: the sum of two sides of a triangle is at least the third side.}
\end{popfigure}

\begin{propriete}[Triangle Inequality]
For any two complex numbers $z_1$ and $z_2$,
\[
|z_1+z_2|\le |z_1|+|z_2|.
\]
Equality holds \textbf{if and only if} $z_1$ and $z_2$ have the same argument (i.e. they point in the same direction), or one of them is zero.
\end{propriete}

\textbf{Proof (examinable).}  Write $z_1=r_1e^{i\theta_1}$, $z_2=r_2e^{i\theta_2}$ with $r_1,r_2\ge0$.  Then
\[
|z_1+z_2|^2 = (z_1+z_2)(\overline{z_1}+\overline{z_2})
= |z_1|^2+|z_2|^2+2\Re(z_1\overline{z_2})
= r_1^2+r_2^2+2r_1r_2\cos(\theta_1-\theta_2).
\]
Since $\cos(\theta_1-\theta_2)\le 1$, we obtain $|z_1+z_2|^2\le (r_1+r_2)^2$, hence the inequality.  Equality requires $\cos(\theta_1-\theta_2)=1$, i.e. $\theta_1=\theta_2\pmod{2\pi}$.

\begin{propriete}[Reverse Triangle Inequality]
For any $z_1,z_2\in\mathbb{C}$,
\[
\bigl||z_1|-|z_2|\bigr|\le |z_1-z_2|.
\]
\end{propriete}
\textit{Proof.}  Apply the triangle inequality to $z_1 = (z_1-z_2)+z_2$ and to $z_2 = (z_2-z_1)+z_1$:
\[
|z_1| = |(z_1-z_2)+z_2| \le |z_1-z_2|+|z_2| \implies |z_1|-|z_2| \le |z_1-z_2|.
\]
Swapping $z_1$ and $z_2$ gives $|z_2|-|z_1| \le |z_1-z_2|$.  Combining both yields the result.

\begin{aretenir}[Key Uses of the Triangle Inequalities]
\begin{itemize}
\item Bounding a sum: $|z_1+z_2|\le |z_1|+|z_2|$ (upper bound).
\item Bounding a difference from below: $|z_1-z_2|\ge \bigl||z_1|-|z_2|\bigr|$ (lower bound).
\item Estimating moduli of polynomials: if $|z|=R$ then $|a_nz^n+\dots+a_0|\le |a_n|R^n+\dots+|a_0|$ (useful for locating roots — a GCE‑useful extra).
\item Error regions: a GPS receiver in Douala gives a position $z_0$ with accuracy $r$; the true position lies in the disc $|z-z_0|\le r$.
\end{itemize}
\end{aretenir}

\begin{attention}[Common Mistake]
$|z_1+z_2| = |z_1|+|z_2|$ is \textbf{not} true in general.  It holds only when $z_1$ and $z_2$ have exactly the same argument (or one is zero).  Always check the arguments before claiming equality.
\end{attention}

\courssub{Solving Modulus Inequalities in the Argand Diagram}

The inequalities $|z-a|<r$, $|z-a|\ge r$, $|z-a|\le |z-b|$, and double inequalities $r_1<|z-a|\le r_2$ define discs, half‑planes, annuli and their combinations.  The boundary is always a circle or a line (the perpendicular bisector of $a$ and $b$ for $|z-a|=|z-b|$).

\begin{methode}[Sketching the Solution Region of a Modulus Inequality]
\begin{enumerate}
\item Identify the geometric locus of the corresponding equality (circle, line, or ray).
\item Choose a test point not on the boundary (often the origin, if it is not on the boundary) and substitute it into the inequality.
\item Shade the side containing the test point if it satisfies the inequality; otherwise shade the opposite side.
\item For a system of inequalities, the solution region is the intersection of the individual shaded regions.
\item Indicate whether boundaries are included (solid line) or excluded (dashed line).
\end{enumerate}
\end{methode}

\begin{exoresolu}[System of Modulus Inequalities]
Sketch the region defined by
\[
|z-1|\le 2 \quad\text{and}\quad |z-1|\ge |z+1|.
\]
\end{exoresolu}
\begin{center}
\textbf{Solution.}
\end{center}
\begin{itemize}
\item $|z-1|\le 2$ is the closed disc centred at $1$ with radius $2$.  Boundary: solid circle.
\item $|z-1|\ge |z+1|$ means the distance to $1$ is at least the distance to $-1$.  The boundary $|z-1|=|z+1|$ is the perpendicular bisector of $1$ and $-1$, i.e. the imaginary axis $\Re(z)=0$.  Test $z=2$: $|2-1|=1$, $|2+1|=3$, so $1\ge3$ is false.  Hence the region is the left half‑plane $\Re(z)\le0$ (including the axis).
\item The intersection is the part of the disc $|z-1|\le2$ that lies on or to the left of the imaginary axis.
\end{itemize}

\begin{popfigure}
\begin{tikzpicture}[scale=1.2,>=stealth]
\draw[popline] (-3.5,0) -- (3.5,0) node[right] {$\Re(z)$};
\draw[popline] (0,-3) -- (0,3) node[above] {$\Im(z)$};
% disc centered at (1,0) radius 2
\draw[thick,popblue,fill=popblueL,opacity=0.3] (1,0) circle (2);
% imaginary axis
\draw[thick,poporange,dashed] (0,-3) -- (0,3);
% shade intersection: clip to left half-plane and disc
\begin{scope}
\clip (-3.5,-3) rectangle (0,3);
\fill[popgreenL,opacity=0.4] (1,0) circle (2);
\end{scope}
% points
\fill (1,0) circle (1.5pt) node[below right] {$1$};
\fill (-1,0) circle (1.5pt) node[below left] {$-1$};
\fill (0,0) circle (1.5pt) node[above left] {$0$};
\end{tikzpicture}
\legende{Solution region of $|z-1|\le2$ and $|z-1|\ge|z+1|$ (shaded).}
\end{popfigure}

\courssub{Rigid Motions (Isometries) of the Plane}

\begin{definition}[Rigid Motion / Isometry]
A transformation $T:\mathbb{C}\to\mathbb{C}$ is a \emph{rigid motion} (or \emph{isometry}) if it preserves distances:
\[
|T(z_1)-T(z_2)| = |z_1-z_2|\qquad\forall\,z_1,z_2\in\mathbb{C}.
\]
\end{definition}

Every isometry of the plane can be written in one of two complex forms:

\begin{propriete}[Complex Form of Plane Isometries]
Let $T$ be an isometry of $\mathbb{C}$.  Then there exist $\theta\in\mathbb{R}$ and $b\in\mathbb{C}$ such that either
\[
\text{(direct)}\quad T(z)=e^{i\theta}z+b
\qquad\text{or}\qquad
\text{(indirect)}\quad T(z)=e^{i\theta}\overline{z}+b.
\]
\end{propriete}
\textit{Proof sketch.}  An isometry fixes the origin after a translation, so we may assume $T(0)=0$.  Then $|T(z)|=|z|$ and $|T(z)-T(w)|=|z-w|$.  Using the polarization identity one shows $T$ is $\mathbb{R}$-linear and preserves the inner product, hence its matrix is orthogonal.  In complex notation orthogonal matrices correspond to multiplication by $e^{i\theta}$ (determinant $+1$, orientation preserving) or multiplication by $e^{i\theta}$ followed by conjugation (determinant $-1$, orientation reversing).  Adding back the translation gives the stated forms.

\begin{definition}[Orientation]
A direct isometry (determinant $+1$) \emph{preserves} orientation; an indirect isometry (determinant $-1$) \emph{reverses} orientation.  Equivalently, a direct isometry maps a positively oriented triangle to a positively oriented triangle.
\end{definition}

\courssub{Classification of Plane Isometries}

\begin{propriete}[Classification of Plane Isometries]
Every isometry of the plane is exactly one of the following four types:
\begin{enumerate}
\item \textbf{Translation:} $z\mapsto z+b$ ($b\neq0$).  Direct, no fixed points.
\item \textbf{Rotation:} $z\mapsto e^{i\theta}z+b$ with $\theta\not\equiv0\pmod{2\pi}$.  Direct, exactly one fixed point (the centre).
\item \textbf{Reflection:} $z\mapsto e^{i\theta}\overline{z}+b$ where the translation component $b$ is orthogonal to the direction $e^{i\theta/2}$ (i.e. $b$ has no component parallel to the mirror axis).  Indirect, fixed points form a line (the mirror axis).
\item \textbf{Glide Reflection:} $z\mapsto e^{i\theta}\overline{z}+b$ where the translation component $b$ has a non‑zero component parallel to the mirror axis.  Indirect, no fixed points.
\end{enumerate}
\end{propriete}

\begin{propriete}[Characterisation of Direct Isometries]
A direct isometry with a fixed point is a rotation.  Indeed, if $T(z)=e^{i\theta}z+b$ and $T(z_0)=z_0$, then $b=(1-e^{i\theta})z_0$, so $T(z)=e^{i\theta}(z-z_0)+z_0$, a rotation about $z_0$ by angle $\theta$.
\end{propriete}

\begin{propriete}[Characterisation of Indirect Isometries]
An indirect isometry $T(z)=e^{i\theta}\overline{z}+b$:
\begin{itemize}
\item If $b$ is orthogonal to the direction $e^{i\theta/2}$ (i.e. $\Re(b e^{-i\theta/2})=0$), then $T$ is a reflection in a line.
\item Otherwise $T$ is a glide reflection: a reflection in a line followed by a translation parallel to that line.
\end{itemize}
\end{propriete}

\begin{popfigure}
\begin{tikzpicture}[scale=1.2,>=stealth]
% axis line at 45 degrees
\draw[thick,poppurple] (-3,-3) -- (3,3) node[above right] {axis $L$};
% translation vector parallel to axis
\draw[->,thick,poporange] (0,0) -- (1.5,1.5) node[midway,above left] {$\vec{v}$};
% original flag (triangle)
\fill[popblueL,opacity=0.6] (0.5,1.5) -- (1.2,1.5) -- (0.5,2.2) -- cycle;
\node at (0.7,1.7) {F};
% reflected flag across axis
\fill[poppink,opacity=0.6] (-1.5,-0.5) -- (-2.2,-0.5) -- (-1.5,-1.2) -- cycle;
\node at (-1.7,-0.7) {F'};
% glide: reflect then translate along axis
\draw[->,dashed,popgreen] (0.5,1.5) .. controls (1,2) and (2,1) .. (1.5,1.5);
\node at (1.8,1.8) {glide};
\end{tikzpicture}
\legende{Glide reflection: reflection in axis $L$ followed by translation $\vec{v}$ parallel to $L$.  The flag changes orientation (F $\to$ F').}
\end{popfigure}

\begin{aretenir}[Composition Rules for Isometries]
\begin{itemize}
\item Direct $\circ$ Direct = Direct (translation or rotation).
\item Direct $\circ$ Indirect = Indirect (reflection or glide reflection).
\item Indirect $\circ$ Direct = Indirect.
\item Indirect $\circ$ Indirect = Direct.
\end{itemize}
\end{aretenir}

\courssub{Worked Example: Identifying an Isometry}

\begin{exoresolu}[Identify the Isometry]
Determine the type of the transformation $T(z)=i\overline{z}+2-i$.
\end{exoresolu}
\begin{center}
\textbf{Solution.}
\end{center}
The map has the form $e^{i\theta}\overline{z}+b$ with $e^{i\theta}=i$ ($\theta=\pi/2$) and $b=2-i$.  It is indirect.  The pure reflection part is $z\mapsto i\overline{z}$, which is reflection in the line through the origin making an angle $\theta/2 = \pi/4$ with the real axis (the line $y=x$).  The direction vector of this axis is $e^{i\pi/4} = \frac{1}{\sqrt{2}}(1+i)$.

Decompose $b=2-i$ into components parallel and perpendicular to the axis:
\[
b_{\parallel} = \Re(b e^{-i\pi/4}) e^{i\pi/4},\qquad
b_{\perp} = b - b_{\parallel}.
\]
Compute $b e^{-i\pi/4} = (2-i)\frac{1}{\sqrt{2}}(1-i) = \frac{1}{\sqrt{2}}(2-2i -i +i^2) = \frac{1}{\sqrt{2}}(1-3i)$.
Hence $\Re(b e^{-i\pi/4}) = 1/\sqrt{2} \neq 0$, so there is a non‑zero parallel component.  Therefore $T$ is a \textbf{glide reflection}.

The glide vector is $b_{\parallel} = \frac{1}{\sqrt{2}} e^{i\pi/4} = \frac{1}{2}+\frac{1}{2}i$.
The axis of the glide reflection is the line parallel to $e^{i\pi/4}$ that is mapped onto itself.  Solving $T(z) \in L$ for $z\in L$ with $L: y=x+c$ gives $c=-3/2$.  Thus the axis is the line $y=x-\frac{3}{2}$ (or $\Im(z)=\Re(z)-\frac{3}{2}$).

\courssub{Common Mistakes and Examination Tips}

\begin{attention}[Pitfalls in Modulus Inequalities]
\begin{itemize}
\item Forgetting to test a point when the boundary is a line (the origin may lie on the line).
\item Confusing $|z-a|<r$ (open disc) with $|z-a|\le r$ (closed disc) — use dashed vs solid circles.
\item Misapplying the triangle inequality: $|z_1+z_2|\le|z_1|+|z_2|$ gives an \emph{upper} bound; the reverse inequality gives a \emph{lower} bound for $|z_1-z_2|$.
\end{itemize}
\end{attention}

\begin{attention}[Pitfalls in Rigid Motions]
\begin{itemize}
\item Writing a rotation as $z\mapsto e^{i\theta}z$ without the translation part when the centre is not the origin.
\item Confusing the condition for a glide reflection: the translation vector must have a component \emph{parallel} to the reflection axis (direction $e^{i\theta/2}$), not just any non‑zero translation.
\item Forgetting that the composition of two reflections in parallel lines is a translation, while in intersecting lines it is a rotation.
\end{itemize}
\end{attention}

\begin{savaistu}[Did You Know?]
The classification of plane isometries into four types (translation, rotation, reflection, glide reflection) is a classic result of Euclidean geometry.  In crystallography, the symmetry groups of wallpaper patterns are built from these four isometries combined with discrete translations.  The famous 17 wallpaper groups are exactly the distinct ways to tile the plane with a repeating motif using only rigid motions.
\end{savaistu}

\begin{exercice}[Practice]
\begin{enumerate}
\item Solve and sketch: $|z-2i|<3$ and $|z-2i|\ge|z+2|$.
\item Show that the transformation $T(z)=-\overline{z}+3+4i$ is a glide reflection.  Find its axis and glide vector.
\item Prove that the composition of a rotation by $\alpha$ about $a$ and a rotation by $\beta$ about $b$ is a rotation by $\alpha+\beta$ about some point $c$ (unless $\alpha+\beta\equiv0\pmod{2\pi}$, in which case it is a translation).
\item A GPS receiver in Yaoundé gives a position $z_0=3+4i$ (in km) with an error radius of $0.5$ km.  Write the inequality describing the possible true positions and sketch the region.
\end{enumerate}
\end{exercice}

\begin{corrige}[Exercise 1]
\begin{enumerate}
\item $|z-2i|<3$ is the open disc centre $2i$ (point $(0,2)$), radius $3$.  Boundary: dashed circle.
\item $|z-2i|\ge|z+2|$: the boundary is the perpendicular bisector of $2i$ and $-2$.  Points: $2i=(0,2)$, $-2=(-2,0)$.  Midpoint $=(-1,1)$.  Slope of segment $=1$, so perpendicular slope $=-1$.  Equation: $y-1 = -(x+1) \implies y = -x$.  This is the line $\Re(z)=-\Im(z)$ (or $\arg z = -\pi/4, 3\pi/4$).  Test $z=0$: $|0-2i|=2$, $|0+2|=2$, so $0$ satisfies the inequality (equality).  The region is the closed half‑plane containing the origin, i.e. the side of the line $y=-x$ where $y \ge -x$ (or $\Re(z)+\Im(z)\ge 0$).
\item The intersection is the part of the open disc $|z-2i|<3$ that lies on or above the line $y=-x$.
\end{enumerate}
\end{corrige}

\begin{corrige}[Exercise 2]
$T(z)=-\overline{z}+3+4i = e^{i\pi}\overline{z}+3+4i$.  Here $\theta=\pi$, so the pure reflection is in the line through the origin with direction $e^{i\pi/2}=i$ (the imaginary axis, $x=0$).  The translation vector is $b=3+4i$.

Decompose $b$ relative to the axis direction $i$ (vertical):
Parallel component (vertical): $4i$.
Perpendicular component (horizontal): $3$.

Since the parallel component $4i \neq 0$, $T$ is a \textbf{glide reflection}.

The axis is a vertical line $x=c$ mapped onto itself.  For $z=c+iy$, $T(z) = -c+3 + i(y+4)$.  This lies on $x=c$ iff $-c+3=c \implies c=1.5$.  So the axis is the line $x=1.5$ (or $\Re(z)=1.5$).

The glide vector is the parallel component $4i$ (translation of $4$ units upward along the axis).
\end{corrige}

\begin{corrige}[Exercise 3]
Let $R_a(z)=e^{i\alpha}(z-a)+a$ and $R_b(z)=e^{i\beta}(z-b)+b$.
Then $R_b(R_a(z)) = e^{i\beta}\bigl(e^{i\alpha}(z-a)+a - b\bigr) + b = e^{i(\alpha+\beta)}z + \bigl[ e^{i\beta}(a-b) - e^{i\alpha}a + b \bigr]$.
The linear part is $e^{i(\alpha+\beta)}$.
\begin{itemize}
\item If $\alpha+\beta \not\equiv 0 \pmod{2\pi}$, then $e^{i(\alpha+\beta)}\neq 1$.  The transformation has a unique fixed point $c$ satisfying $c = e^{i(\alpha+\beta)}c + K$, where $K$ is the constant bracket.  Solving gives $c = \frac{K}{1-e^{i(\alpha+\beta)}}$.  Hence it is a rotation by $\alpha+\beta$ about $c$.
\item If $\alpha+\beta \equiv 0 \pmod{2\pi}$, then $e^{i(\alpha+\beta)}=1$ and the transformation reduces to $z \mapsto z + K$, a translation by vector $K$.
\end{itemize}
\end{corrige}

\begin{corrige}[Exercise 4]
The true position $z$ satisfies $|z - (3+4i)| \le 0.5$.  This is a closed disc centred at $(3,4)$ with radius $0.5$ km.
\end{corrige}

\newpage

\coursec{Integration activity}

\courssub{Integration Activity: Camtel-Sud Relay Network}

\textbf{Aims of the activity}\par
This integration activity brings together the whole chapter in a single modelling task. By the end of it, you should be able to:
\begin{itemize}
\item convert complex numbers between \cle{Cartesian}, \cle{polar} and \cle{exponential} forms and plot affixes in an Argand diagram;
\item use \cle{De Moivre's theorem} to find all the $n$-th roots of a complex number and interpret them geometrically;
\item describe and shade \cle{loci} defined by modulus conditions (circles, annuli, perpendicular bisectors);
\item recognise a \cle{direct similitude} from its complex form $z' = az+b$, identify its ratio, angle and centre, and write its \cle{matrix};
\item apply \cle{modulus inequalities} (triangle inequality) to justify bounds in a concrete situation;
\item classify a \cle{rigid motion} as orientation-preserving or orientation-reversing;
\item communicate a complete mathematical argument, as required at the GCE Advanced Level.
\end{itemize}

\begin{aretenir}[Key skills assessed]
\begin{itemize}
\item De Moivre's theorem: $(\cos\theta + i\sin\theta)^n = \cos n\theta + i\sin n\theta$ and $n$-th roots.
\item Loci: $|z-a|=r$ (circle), $|z-a|\le r$ (disc), $|z-a|=|z-b|$ (perpendicular bisector).
\item Direct similitude $z' = az+b$: ratio $|a|$, angle $\arg a$, centre $\frac{b}{1-a}$ (if $a\neq1$).
\item Matrix of $z\mapsto \alpha z + \beta \bar{z}$; for $z\mapsto az$ (direct) it is $\begin{pmatrix} \Re a & -\Im a \\ \Im a & \Re a \end{pmatrix}$.
\item Rigid motion: isometry $\Leftrightarrow |a|=1$; orientation-preserving $\Leftrightarrow \det>0$ (direct), orientation-reversing $\Leftrightarrow \det<0$ (opposite).
\end{itemize}
\end{aretenir}

\textbf{Situation}\par
A telecommunications company, \emph{Camtel-Sud}, installs a relay mast at a point $\Omega$ on a hill in Yaoundé. The surrounding region is modelled by the complex plane (the \cle{Argand diagram}), with $\Omega$ at the origin; distances are measured in kilometres, the positive real axis pointing East and the positive imaginary axis pointing North.

Three relay stations $A$, $B$ and $C$ are placed so that their affixes $z_A$, $z_B$, $z_C$ are exactly the three \cle{cube roots of $8$}, i.e.\ the solutions of $z^3 = 8$.

The mast carries a \textbf{rotating directional antenna}. Each time the control system sends a pulse, the beam rotates and slightly contracts: a signal point at affix $z$ is moved to the affix
\[ z' = \frac{1}{\sqrt{2}}\,e^{i\pi/6}\, z. \]
The engineers apply this transformation repeatedly to track the successive positions $z_0, z_1, z_2, \dots$ of a test signal initially sent to $z_0 = 2\sqrt{2}$ (a point $2\sqrt{2}$ km East of the mast).

The \textbf{coverage zone} $\mathcal{D}$ guaranteed by the company is the set of points $M(z)$ satisfying simultaneously
\[ 1 \le |z| \le 3 \qquad \text{and} \qquad |z - 1| \ge |z + 1|. \]
A client $K$ lives at the point of affix $z_K = 2e^{i\pi/4}$.

The company wants a complete mathematical report: positions of the relays, behaviour of the rotating beam, shape of the coverage zone, and a decision on whether the client $K$ is covered.

\begin{experience}[Guided investigation]
Work through the following tasks step by step. Use a separate sheet for your Argand diagram (scale: $1$ cm for $1$ km).
\end{experience}

\textbf{Tasks}\par
\begin{enumerate}
\item \textbf{Relays via De Moivre.} Write $8$ in exponential form. Using De Moivre's theorem, determine the three cube roots of $8$ in exponential form, then in Cartesian form. Plot $\Omega$, $A$, $B$, $C$ on an Argand diagram (scale: $1$ cm for $1$ km). What is the nature of triangle $ABC$? Justify.

\item \textbf{The rotating beam.} Write $a = \dfrac{1}{\sqrt{2}}e^{i\pi/6}$ in the form $a = r(\cos\theta + i\sin\theta)$.
\begin{enumerate}
\item Identify the \cle{ratio} and the \cle{angle} of the similitude $S:\ z' = az$. State its centre and its type.
\item Compute $z_1 = S(z_0)$, $z_2 = S(z_1)$ and $z_3 = S(z_2)$ where $z_0 = 2\sqrt{2}$. Give each in exponential form.
\item Show that $z_n = 2\sqrt{2}\left(\dfrac{1}{\sqrt{2}}\right)^{\!n} e^{in\pi/6}$, and find the smallest $n$ for which the signal point enters the disc $|z| < 1$.
\item Describe geometrically the path of the successive signal points.
\end{enumerate}

\item \textbf{The coverage zone.}
\begin{enumerate}
\item Show that the condition $1 \le |z| \le 3$ describes a closed annulus; give its boundary circles, centres and radii.
\item Show that $|z-1| \ge |z+1|$ is equivalent to $\operatorname{Re}(z) \le 0$. Which line is the boundary, and how is it related to the points of affixes $1$ and $-1$?
\item Shade the coverage zone $\mathcal{D}$ on your Argand diagram.
\end{enumerate}

\item \textbf{Modulus inequalities.} Using the triangle inequality, prove that every point of $\mathcal{D}$ satisfies $|z - 1| \le 4$. Is the bound $4$ attained? If yes, at which point(s) of $\mathcal{D}$?

\item \textbf{The client.} Compute $|z_K|$ and $\operatorname{Re}(z_K)$ for $z_K = 2e^{i\pi/4}$. Decide, with justification, whether the client $K$ is covered.

\item \textbf{Matrix and rigid motion.}
\begin{enumerate}
\item Write the matrix of the linear map $z \mapsto az$ (rotation part combined with the contraction) in the canonical basis.
\item Consider now the pure rotation $R:\ z' = e^{i\pi/6}z$. Write its matrix, compute its determinant, and classify $R$ as a rigid motion: is it orientation-preserving or orientation-reversing? Justify.
\end{enumerate}

\item \textbf{Synthesis.} Prepare a group poster: the annotated Argand diagram (relays, beam positions, shaded coverage zone, client $K$), and a five-line conclusion for the company.
\end{enumerate}

\begin{attention}[Common pitfalls]
\begin{itemize}
\item When using De Moivre for roots, remember to add $2k\pi$ \emph{before} dividing by $n$: $z_k = \sqrt[n]{r}\, e^{i(\theta+2k\pi)/n}$.
\item The condition $|z-1|\ge|z+1|$ means ``distance to $1$ is at least distance to $-1$''; the boundary is the perpendicular bisector of $[-1,1]$, i.e. the imaginary axis. Do not confuse with $|z-1|\le|z+1|$.
\item Triangle inequality: $|z_1+z_2|\le|z_1|+|z_2|$. Equality holds iff $z_1$ and $z_2$ have the same argument (or one is zero).
\item A direct similitude $z'=az+b$ with $a\neq1$ has centre $\frac{b}{1-a}$. If $b=0$, the centre is the origin.
\item For a rigid motion $z'=e^{i\theta}z$, the matrix is a rotation matrix with determinant $+1$ (orientation-preserving). A reflection would have determinant $-1$.
\end{itemize}
\end{attention}

\textbf{Model Answer}\par

\begin{exoresolu}[Complete solution]
\textbf{1. Cube roots of $8$.} 
$8 = 8e^{i\cdot 0}$. By De Moivre's theorem, the cube roots are
\[ w_k = 8^{1/3} e^{i(0 + 2k\pi)/3} = 2e^{2ik\pi/3}, \qquad k = 0, 1, 2. \]
Explicitly:
\[ z_A = 2e^{0} = 2, \quad z_B = 2e^{2i\pi/3} = -1 + i\sqrt{3}, \quad z_C = 2e^{4i\pi/3} = -1 - i\sqrt{3}. \]
The three points lie on the circle of centre $\Omega$ and radius $2$, separated by angles $\dfrac{2\pi}{3}$. They form an \textbf{equilateral triangle} inscribed in that circle (side length $|z_A - z_B| = |3 - i\sqrt{3}| = \sqrt{9+3} = 2\sqrt{3}$).

\textbf{2. The rotating beam.} 
(a) $a = \dfrac{1}{\sqrt{2}}\left(\cos\dfrac{\pi}{6} + i\sin\dfrac{\pi}{6}\right)$. The similitude $S$ has \cle{ratio} $r = |a| = \dfrac{1}{\sqrt{2}}$, \cle{angle} $\theta = \arg a = \dfrac{\pi}{6}$, and \cle{centre} $\Omega$ (since $b = 0$). As $r \ne 1$, $S$ is a \textbf{direct similitude that is not an isometry}: the composite of the rotation of angle $\pi/6$ about $\Omega$ and the enlargement of ratio $1/\sqrt{2}$ centred at $\Omega$.

(b) With $z_0 = 2\sqrt{2} = 2\sqrt{2}\,e^{i\cdot 0}$:
\[ z_1 = \tfrac{1}{\sqrt{2}}e^{i\pi/6}\cdot 2\sqrt{2} = 2e^{i\pi/6}, \quad z_2 = \sqrt{2}\,e^{i\pi/3}, \quad z_3 = e^{i\pi/2} = i. \]

(c) By induction, $z_n = a^n z_0$, and $a^n = \left(\tfrac{1}{\sqrt{2}}\right)^{\!n} e^{in\pi/6}$ by De Moivre, hence
\[ z_n = 2\sqrt{2}\left(\tfrac{1}{\sqrt{2}}\right)^{\!n} e^{in\pi/6}. \]
Then $|z_n| = 2\sqrt{2}\left(\tfrac{1}{\sqrt{2}}\right)^{n} = 2^{(3-n)/2}$, and $|z_n| < 1 \iff \dfrac{3-n}{2} < 0 \iff n > 3$. The smallest such integer is $\boxed{n = 4}$ (indeed $|z_4| = \tfrac{1}{\sqrt{2}} < 1$).

(d) The points spiral in towards $\Omega$: each step rotates by $30^{\circ}$ and multiplies the distance to $\Omega$ by $\dfrac{1}{\sqrt{2}}$; they lie on an equiangular (logarithmic-type) spiral sampled at equal angles.

\textbf{3. Coverage zone.} 
(a) $|z| = 1$ and $|z| = 3$ are the circles of centre $\Omega$ and radii $1$ and $3$; $1 \le |z| \le 3$ is the closed annulus between them (boundaries included).

(b) Writing $z = x + iy$:
\[ |z-1|^2 \ge |z+1|^2 \iff (x-1)^2 + y^2 \ge (x+1)^2 + y^2 \iff -4x \ge 0 \iff x \le 0. \]
The boundary $x = 0$ (the imaginary axis) is the \textbf{perpendicular bisector} of the segment joining the points of affixes $1$ and $-1$; the condition selects the half-plane closer to $-1$ than to $1$.

(c) $\mathcal{D}$ is the \textbf{left half-annulus}: the part of the annulus $1 \le |z| \le 3$ lying in the half-plane $\operatorname{Re}(z) \le 0$, boundaries included.

\begin{popfigure}
\begin{tikzpicture}[scale=0.85]
\fill[popblueL, even odd rule] (0,0) circle (3) (0,0) circle (1);
\fill[white] (0,-3) rectangle (3,3);
\draw[popline] (0,0) circle (1);
\draw[popline] (0,0) circle (3);
\draw[->, gray] (-3.6,0) -- (3.6,0) node[right]{\scriptsize Re};
\draw[->, gray] (0,-3.6) -- (0,3.6) node[above]{\scriptsize Im};
\draw[popdark, thick] (0,-3.4) -- (0,3.4);
\fill (0,0) circle (1.5pt) node[below right]{\scriptsize $\Omega$};
\fill[poporange] (2,0) circle (2pt) node[below right]{\scriptsize $A$};
\fill[poporange] (-1,1.732) circle (2pt) node[above left]{\scriptsize $B$};
\fill[poporange] (-1,-1.732) circle (2pt) node[below left]{\scriptsize $C$};
\fill[popgreen] (1.414,1.414) circle (2pt) node[right]{\scriptsize $K$};
\fill[poppurple] (1.732,1) circle (1.5pt) node[right]{\scriptsize $z_1$};
\fill[poppurple] (0.707,1.225) circle (1.5pt) node[above right]{\scriptsize $z_2$};
\fill[poppurple] (0,1) circle (1.5pt) node[above left]{\scriptsize $z_3$};
\end{tikzpicture}
\legende{Coverage zone $\mathcal{D}$ (shaded), relays $A,B,C$, signal points $z_1,z_2,z_3$, client $K$.}
\end{popfigure}

\textbf{4. Modulus inequality.} For $z \in \mathcal{D}$, $|z| \le 3$, so by the triangle inequality
\[ |z - 1| \le |z| + |-1| = |z| + 1 \le 3 + 1 = 4. \]
Equality in $|z-1| \le |z|+1$ holds iff $z$ and $-1$ have the same direction, i.e.\ $z$ is a negative real number. Combined with $|z|=3$ and $\operatorname{Re}(z)\le0$, this gives $z = -3$. Since $-3 \in \mathcal{D}$ (as $|-3| = 3 \in [1,3]$ and $\operatorname{Re}(-3) = -3 \le 0$), the bound $4$ \textbf{is attained} at $z = -3$.

\textbf{5. The client.} $z_K = 2e^{i\pi/4}$, so $|z_K| = 2 \in [1,3]$: the first condition holds. But
\[ \operatorname{Re}(z_K) = 2\cos\frac{\pi}{4} = \sqrt{2} > 0, \]
so the second condition fails. \textbf{Conclusion: the client $K$ is NOT covered} by the guaranteed zone (he is in the Eastern half-plane).

\textbf{6. Matrix and rigid motion.} 
(a) $a = \dfrac{1}{\sqrt{2}}\left(\dfrac{\sqrt{3}}{2} + i\dfrac{1}{2}\right) = \dfrac{\sqrt{6}}{4} + i\dfrac{\sqrt{2}}{4}$. The matrix of $z \mapsto az$ in the canonical basis is
\[ M = \begin{pmatrix} \dfrac{\sqrt{6}}{4} & -\dfrac{\sqrt{2}}{4} \\[2mm] \dfrac{\sqrt{2}}{4} & \dfrac{\sqrt{6}}{4} \end{pmatrix} = \frac{1}{\sqrt{2}}\begin{pmatrix} \cos\frac{\pi}{6} & -\sin\frac{\pi}{6} \\[1mm] \sin\frac{\pi}{6} & \cos\frac{\pi}{6} \end{pmatrix}. \]

(b) For $R:\ z' = e^{i\pi/6}z$, the matrix is the rotation matrix
\[ M_R = \begin{pmatrix} \cos\frac{\pi}{6} & -\sin\frac{\pi}{6} \\[1mm] \sin\frac{\pi}{6} & \cos\frac{\pi}{6} \end{pmatrix}, \qquad \det M_R = \cos^2\tfrac{\pi}{6} + \sin^2\tfrac{\pi}{6} = 1. \]
Since $|e^{i\pi/6}| = 1$, $R$ preserves distances: it is a \textbf{rigid motion} (an isometry). Since $\det M_R = +1 > 0$, $R$ is \textbf{orientation-preserving}: it is a \emph{positive} (direct) rigid motion, namely the rotation of centre $\Omega$ and angle $\dfrac{\pi}{6}$.

\textbf{7. Synthesis (sample conclusion).} The three relays form an equilateral triangle of side $2\sqrt{3}$ km around the mast. The beam describes a contracting spiral (angle $30^{\circ}$, ratio $\tfrac{1}{\sqrt{2}}$) and enters the 1 km dead-zone after 4 pulses. The guaranteed coverage is the half-annulus $\{1 \le |z| \le 3,\ \operatorname{Re}(z) \le 0\}$; the client $K$ lies outside it and requires an additional relay.
\end{exoresolu}

\begin{savaistu}[Did you know?]
The spiral described by the signal points is a discrete version of the \emph{logarithmic spiral} $r = r_0 e^{k\theta}$, which appears in nature (nautilus shells, galaxies) and in engineering (antenna patterns). The transformation $z\mapsto az$ with $|a|<1$ and $\arg a\neq0$ is a \emph{spiral similarity}; its matrix is a scaled rotation matrix. In the GCE exam, you may be asked to identify the centre of a similitude given by $z'=az+b$ ($a\neq1$): it is the unique fixed point $z = \frac{b}{1-a}$.
\end{savaistu}

\newpage

\coursec{Methods and worked examples}

\courssub{Methods and Worked Examples}

\begin{methode}[Converting Forms and Applying De Moivre's Theorem for Powers and Roots]
\textbf{Goal:} Switch fluently between Cartesian ($z=x+iy$), polar ($z=r(\cos\theta+i\sin\theta)$) and exponential ($z=re^{i\theta}$) forms, and compute integer powers or $n$-th roots using De Moivre's Theorem.

\textbf{Steps:}
\begin{enumerate}
    \item \textbf{Identify the given form.} If Cartesian, compute the modulus $r=\sqrt{x^2+y^2}$ and the argument $\theta=\arg(z)$ (fix the quadrant using $\tan\theta=y/x$ and the signs of $x$ and $y$). Principal argument: $-\pi < \theta \le \pi$.
    \item \textbf{Write in polar/exponential form:} $z=r(\cos\theta+i\sin\theta)=re^{i\theta}$.
    \item \textbf{For powers $z^n$ ($n\in\mathbb{Z}$):} apply De Moivre: $z^n = r^n(\cos n\theta + i\sin n\theta)$. Convert back to Cartesian form if required.
    \item \textbf{For $n$-th roots $z^{1/n}$:} the $n$ distinct roots are
    \[
    w_k = r^{1/n}\left[\cos\left(\frac{\theta+2k\pi}{n}\right) + i\sin\left(\frac{\theta+2k\pi}{n}\right)\right],\quad k=0,1,\dots,n-1.
    \]
    \item \textbf{Geometric interpretation:} the roots lie on the circle of radius $r^{1/n}$ centred at the origin, at the vertices of a regular $n$-gon. One root (the principal root) has argument $\theta/n$.
\end{enumerate}

\textbf{Key reflexes:}
\begin{itemize}
    \item Always reduce the argument of the final answer to the principal range $(-\pi,\pi]$.
    \item For negative powers: $z^{-n} = \dfrac{1}{z^n} = r^{-n}(\cos n\theta - i\sin n\theta)$.
    \item Check: the sum of the $n$-th roots of a non-zero complex number is $0$; their product is $(-1)^{n-1}z$ (Vieta's formulas applied to $w^n - z = 0$).
\end{itemize}
\end{methode}

\begin{exoresolu}[Powers, Roots and Geometric Properties]
\textbf{Statement:} Let $z = -\sqrt{3} + i$.
\begin{enumerate}
    \item Express $z$ in exponential form $re^{i\theta}$ with $r>0$ and $-\pi < \theta \le \pi$.
    \item Find $z^6$ in Cartesian form $a+ib$.
    \item Determine the three cube roots of $z$ in exponential form. Plot them on an Argand diagram.
    \item Show that the sum of these three cube roots is zero and find their product.
\end{enumerate}

\tcblower
\textbf{Detailed Solution:}

\textbf{1. Exponential form of $z$.}
Modulus: $r = |z| = \sqrt{(-\sqrt{3})^2 + 1^2} = \sqrt{3+1} = 2$.
Argument: $\tan\theta = \dfrac{1}{-\sqrt{3}} = -\dfrac{1}{\sqrt{3}}$. The point $(-\sqrt{3}, 1)$ is in the \textbf{second quadrant}.
Reference angle: $\dfrac{\pi}{6}$. Principal argument: $\theta = \pi - \dfrac{\pi}{6} = \dfrac{5\pi}{6}$.
\[
z = 2e^{i\frac{5\pi}{6}} \quad \left(\text{or } 2\left(\cos\frac{5\pi}{6}+i\sin\frac{5\pi}{6}\right)\right).
\]

\textbf{2. Compute $z^6$.}
By De Moivre's Theorem:
\[
z^6 = 2^6 e^{i\left(6\times\frac{5\pi}{6}\right)} = 64 e^{i5\pi}.
\]
Reduce the argument: $5\pi \equiv 5\pi - 2(2\pi) = \pi \pmod{2\pi}$.
\[
z^6 = 64 e^{i\pi} = 64(\cos\pi + i\sin\pi) = 64(-1 + 0i) = \boxed{-64}.
\]

\textbf{3. Cube roots of $z$.}
We seek $w_k$ such that $w_k^3 = z = 2e^{i\frac{5\pi}{6}}$.
Modulus of the roots: $\sqrt[3]{2} = 2^{1/3}$.
Arguments: $\dfrac{\frac{5\pi}{6} + 2k\pi}{3} = \dfrac{5\pi}{18} + \dfrac{2k\pi}{3},\quad k=0,1,2$.
\begin{align*}
k=0:\quad & w_0 = 2^{1/3} e^{i\frac{5\pi}{18}} \\
k=1:\quad & w_1 = 2^{1/3} e^{i\left(\frac{5\pi}{18}+\frac{2\pi}{3}\right)} = 2^{1/3} e^{i\frac{17\pi}{18}} \\
k=2:\quad & w_2 = 2^{1/3} e^{i\left(\frac{5\pi}{18}+\frac{4\pi}{3}\right)} = 2^{1/3} e^{i\frac{29\pi}{18}}.
\end{align*}
Reduce the argument of $w_2$ to the principal range: $\dfrac{29\pi}{18} - 2\pi = \dfrac{29\pi}{18} - \dfrac{36\pi}{18} = -\dfrac{7\pi}{18}$.
So $w_2 = 2^{1/3} e^{-i\frac{7\pi}{18}}$.

\textbf{Argand diagram} (note $\frac{5\pi}{18}=50^\circ$, $\frac{17\pi}{18}=170^\circ$, $-\frac{7\pi}{18}=-70^\circ$, and $2^{1/3}\approx 1.26$):
\begin{popfigure}
\begin{tikzpicture}[scale=1.5]
\draw[->] (-1.6,0) -- (1.6,0) node[right] {$\Re$};
\draw[->] (0,-1.6) -- (0,1.6) node[above] {$\Im$};
\draw[popline, dashed] (0,0) circle (1.26);
\draw[popblue, thick] (0,0) -- (50:1.26);
\draw[popblue, thick] (0,0) -- (170:1.26);
\draw[popblue, thick] (0,0) -- (-70:1.26);
\fill[popblue] (50:1.26) circle (1.5pt) node[above right, font=\small] {$w_0$};
\fill[popblue] (170:1.26) circle (1.5pt) node[above left, font=\small] {$w_1$};
\fill[popblue] (-70:1.26) circle (1.5pt) node[below right, font=\small] {$w_2$};
\draw[popmuted] (0.55,0) arc (0:50:0.55) node[midway, right, font=\tiny] {$\frac{5\pi}{18}$};
\end{tikzpicture}
\legende{Cube roots of $z$: vertices of an equilateral triangle on the circle of radius $\sqrt[3]{2}$.}
\end{popfigure}

\textbf{4. Sum and product of the roots.}
The roots $w_0, w_1, w_2$ are the solutions of $w^3 - z = 0$, i.e. $w^3 + 0w^2 + 0w - z = 0$.
By Vieta's formulas:
Sum $= w_0+w_1+w_2 = -\dfrac{\text{coefficient of }w^2}{\text{coefficient of }w^3} = 0$.
Product $= w_0 w_1 w_2 = -\dfrac{\text{constant term}}{\text{leading coefficient}} = -(-z) = z = \boxed{-\sqrt{3}+i}$.
\emph{Check:} sum of arguments: $\frac{5\pi}{18}+\frac{17\pi}{18}-\frac{7\pi}{18} = \frac{15\pi}{18} = \frac{5\pi}{6}$; product of moduli: $(2^{1/3})^3 = 2$. This matches $z = 2e^{i5\pi/6}$.
\end{exoresolu}

\begin{methode}[Deriving Trigonometric Identities via De Moivre's Theorem]
\textbf{Goal:} express $\cos n\theta$, $\sin n\theta$ in powers of $\cos\theta$, $\sin\theta$; or express $\cos^n\theta$, $\sin^n\theta$ in terms of cosines/sines of multiple angles.

\textbf{Steps (Powers $\to$ Multiple angles):}
\begin{enumerate}
    \item Write $2\cos\theta = z + z^{-1}$ and $2i\sin\theta = z - z^{-1}$, where $z=e^{i\theta}$.
    \item Raise to the power $n$: $(2\cos\theta)^n = (z+z^{-1})^n$, $(2i\sin\theta)^n = (z-z^{-1})^n$.
    \item Expand using the Binomial Theorem: $(z+z^{-1})^n = \displaystyle\sum_{k=0}^n \binom{n}{k} z^{n-2k}$.
    \item Group the symmetric terms using $z^m + z^{-m} = 2\cos m\theta$ and $z^m - z^{-m} = 2i\sin m\theta$.
    \item Simplify to obtain a linear combination of $\cos m\theta$ (or $\sin m\theta$).
\end{enumerate}

\textbf{Steps (Multiple angles $\to$ Powers):}
\begin{enumerate}
    \item Use De Moivre: $(\cos\theta + i\sin\theta)^n = \cos n\theta + i\sin n\theta$.
    \item Expand the left-hand side using the Binomial Theorem.
    \item Equate real parts $\to \cos n\theta$; equate imaginary parts $\to \sin n\theta$.
    \item Use $\sin^2\theta = 1-\cos^2\theta$ (or vice versa) to express the result purely in powers of one function.
\end{enumerate}

\textbf{Key reflexes:}
\[
\cos n\theta = \Re\!\left(e^{in\theta}\right),\qquad \sin n\theta = \Im\!\left(e^{in\theta}\right),\qquad z^m+z^{-m}=2\cos m\theta,\qquad z^m-z^{-m}=2i\sin m\theta .
\]
\end{methode}

\begin{exoresolu}[Trigonometric Expansions and Integration Application]
\textbf{Statement:}
\begin{enumerate}
    \item Use De Moivre's theorem to show that $\cos 5\theta = 16\cos^5\theta - 20\cos^3\theta + 5\cos\theta$.
    \item Hence express $\cos^5\theta$ in terms of cosines of multiple angles.
    \item Evaluate the exact value of $\displaystyle\int_0^{\pi/2} \cos^5\theta\, d\theta$.
\end{enumerate}

\tcblower
\textbf{Detailed Solution:}

\textbf{1. Expansion of $\cos 5\theta$.}
Let $z = \cos\theta + i\sin\theta = e^{i\theta}$.
By De Moivre: $z^5 = \cos 5\theta + i\sin 5\theta$.
Binomial expansion of $(\cos\theta + i\sin\theta)^5$:
\begin{align*}
z^5 &= \cos^5\theta + 5i\cos^4\theta\sin\theta + 10i^2\cos^3\theta\sin^2\theta + 10i^3\cos^2\theta\sin^3\theta + 5i^4\cos\theta\sin^4\theta + i^5\sin^5\theta \\
&= \cos^5\theta + 5i\cos^4\theta\sin\theta - 10\cos^3\theta\sin^2\theta - 10i\cos^2\theta\sin^3\theta + 5\cos\theta\sin^4\theta + i\sin^5\theta.
\end{align*}
Equate real parts:
\[
\cos 5\theta = \cos^5\theta - 10\cos^3\theta\sin^2\theta + 5\cos\theta\sin^4\theta.
\]
Replace $\sin^2\theta = 1-\cos^2\theta$:
\begin{align*}
\cos 5\theta &= \cos^5\theta - 10\cos^3\theta(1-\cos^2\theta) + 5\cos\theta(1-\cos^2\theta)^2 \\
&= \cos^5\theta - 10\cos^3\theta + 10\cos^5\theta + 5\cos\theta(1 - 2\cos^2\theta + \cos^4\theta) \\
&= 11\cos^5\theta - 10\cos^3\theta + 5\cos\theta - 10\cos^3\theta + 5\cos^5\theta \\
&= \boxed{16\cos^5\theta - 20\cos^3\theta + 5\cos\theta}.
\end{align*}

\textbf{2. Express $\cos^5\theta$ in multiple angles.}
Rearrange the identity:
\[
16\cos^5\theta = \cos 5\theta + 20\cos^3\theta - 5\cos\theta.
\]
We also need $\cos^3\theta$. The same method (or the standard identity) gives $\cos 3\theta = 4\cos^3\theta - 3\cos\theta$, hence $\cos^3\theta = \frac{1}{4}(\cos 3\theta + 3\cos\theta)$.
Substitute:
\begin{align*}
16\cos^5\theta &= \cos 5\theta + 20\left[\frac{1}{4}(\cos 3\theta + 3\cos\theta)\right] - 5\cos\theta \\
&= \cos 5\theta + 5\cos 3\theta + 15\cos\theta - 5\cos\theta \\
&= \cos 5\theta + 5\cos 3\theta + 10\cos\theta.
\end{align*}
\[
\boxed{\cos^5\theta = \frac{1}{16}(\cos 5\theta + 5\cos 3\theta + 10\cos\theta)}.
\]

\textbf{3. Evaluate the integral.}
\[
I = \int_0^{\pi/2} \cos^5\theta\, d\theta = \frac{1}{16} \int_0^{\pi/2} (\cos 5\theta + 5\cos 3\theta + 10\cos\theta)\, d\theta.
\]
Integrate term by term using $\displaystyle\int \cos k\theta\, d\theta = \frac{1}{k}\sin k\theta$:
\[
I = \frac{1}{16} \left[ \frac{1}{5}\sin 5\theta + \frac{5}{3}\sin 3\theta + 10\sin\theta \right]_0^{\pi/2}.
\]
At $\theta = \pi/2$: $\sin(5\pi/2) = 1$; $\sin(3\pi/2) = -1$; $\sin(\pi/2) = 1$. At $\theta=0$: all terms vanish.
\[
I = \frac{1}{16} \left( \frac{1}{5} - \frac{5}{3} + 10 \right) = \frac{1}{16}\times\frac{3 - 25 + 150}{15} = \frac{1}{16} \times \frac{128}{15} = \boxed{\frac{8}{15}}.
\]
\end{exoresolu}

\begin{methode}[Sketching Loci in the Argand Diagram]
\textbf{Goal:} interpret and sketch loci defined by equations or inequalities in $z$ (involving $|z|$, $\arg z$, $\Re(z)$, $\Im(z)$).

\textbf{Standard forms and shapes:}
\begin{center}
\begin{tabularx}{\linewidth}{|l|X|l|}\hline
\rowcolor{popblueL}
\textbf{Equation / Inequality} & \textbf{Geometric Description} & \textbf{Boundary} \\ \hline
$|z - a| = r$ & Circle, centre $a$, radius $r$ & Solid circle \\ \hline
$|z - a| < r$ / $> r$ & Interior / exterior of the circle & Dashed circle \\ \hline
$\arg(z - a) = \alpha$ & Half-line from $a$ ($a$ excluded) at angle $\alpha$ & Solid ray, open dot at $a$ \\ \hline
$\alpha < \arg(z - a) < \beta$ & Angular sector, vertex $a$ excluded & Dashed rays \\ \hline
$|z - a| = |z - b|$ & Perpendicular bisector of segment $AB$ & Solid line \\ \hline
$|z - a| = k|z - b|$ ($k>0$, $k\ne 1$) & Circle of Apollonius & Solid circle \\ \hline
$\Re(z)=c$ or $\Im(z)=d$ & Vertical / horizontal line & Solid line \\ \hline
\end{tabularx}
\end{center}

\textbf{Steps for sketching:}
\begin{enumerate}
    \item \textbf{Identify the type} (circle, line, ray, region).
    \item \textbf{Find key points:} centre and radius for circles; intercepts for lines; vertex for rays.
    \item \textbf{Determine the boundary style:} solid for $\le, \ge, =$; dashed for $<, >$.
    \item \textbf{Test a point} (usually the origin $z=0$, if not on the boundary) to shade the correct region for inequalities.
    \item \textbf{Label clearly:} axes, scale, centre coordinates, radius, angles.
\end{enumerate}

\textbf{Traps:}
\begin{itemize}
    \item $\arg(z-a)=\alpha$ excludes the point $a$ itself (open dot), because $\arg$ is undefined there.
    \item $|z-a|=k|z-b|$ with $k=1$ is a line (perpendicular bisector), not a circle.
    \item For $\arg(z) = \alpha$, the half-line starts at the origin, which is excluded.
\end{itemize}
\end{methode}

\begin{exoresolu}[Compound Locus and Extremal Moduli]
\textbf{Statement:} On a single Argand diagram, sketch the region $R$ defined by the simultaneous inequalities
\[
|z - 3i| \le 3 \qquad \text{and} \qquad \frac{\pi}{4} \le \arg(z) \le \frac{\pi}{2}.
\]
Determine the greatest value of $|z|$ for $z\in R$, and discuss whether $|z|$ attains a least value on $R$.

\tcblower
\textbf{Detailed Solution:}

\textbf{Analysis of the inequalities:}
\begin{enumerate}
    \item $|z - 3i| \le 3$: closed disc, centre $C(0,3)$ (affix $3i$), radius $3$. Boundary: solid circle, interior included. Note that $|0-3i|=3$, so the circle \emph{passes through the origin}.
    \item $\frac{\pi}{4} \le \arg(z) \le \frac{\pi}{2}$: closed sector between the ray $\arg(z)=\pi/4$ (the line $y=x$, $x\ge 0$) and the ray $\arg(z)=\pi/2$ (the positive imaginary axis). Boundaries: solid rays. Since $\arg(0)$ is undefined, the origin is \emph{not} in $R$.
\end{enumerate}

\textbf{Key boundary points.} The ray $\arg z=\pi/2$ meets the circle where $z=iy$, $|iy-3i|=|y-3|=3$, giving $y=0$ (origin, excluded) or $y=6$: point $A=6i$.
The ray $\arg z=\pi/4$ meets the circle where $z=x+ix$:
\[
|x+i(x-3)|^2 = x^2+(x-3)^2 = 9 \implies 2x^2-6x=0 \implies x=0 \text{ or } x=3,
\]
giving the origin (excluded) and the point $B=3+3i$.

\textbf{Sketch:}
\begin{popfigure}
\begin{tikzpicture}[scale=1.0]
\draw[->] (-1,0) -- (4.5,0) node[right] {$\Re$};
\draw[->] (0,-0.5) -- (0,6.8) node[above] {$\Im$};
% Shaded region: intersection of disc and sector
\begin{scope}
\clip (0,0) -- (4.5,4.5) -- (0,6.5) -- cycle;
\clip (0,3) circle (3);
\fill[popblueL] (-1,-1) rectangle (5,7);
\end{scope}
% Boundaries
\draw[popblue, thick] (0,3) circle (3);
\draw[popblue, thick] (0,0) -- (4.2,4.2) node[pos=1.05, right, font=\small] {$\arg z=\frac{\pi}{4}$};
\draw[popblue, thick] (0,0) -- (0,6.4) node[above left, font=\small] {$\arg z=\frac{\pi}{2}$};
% Points
\draw[poporange, fill=white, thick] (0,0) circle (2.5pt) node[below left] {$O$ (excluded)};
\fill[poporange] (0,6) circle (2pt) node[above left] {$A=6i$};
\fill[poporange] (3,3) circle (2pt) node[right] {$B=3+3i$};
\fill[popdark] (0,3) circle (2pt) node[left] {$C=3i$};
\draw[popmuted, dashed] (0,3) -- (3,3) node[midway, above] {$3$};
\end{tikzpicture}
\legende{Region $R$: intersection of the disc $|z-3i|\le 3$ with the sector $\pi/4 \le \arg z \le \pi/2$. The origin is excluded.}
\end{popfigure}

\textbf{Greatest modulus.}
$|z|$ is the distance from $O$. The farthest point of the closed disc from $O$ lies on the line through $O$ and the centre $C$, i.e. on the imaginary axis: this is $A=6i$, which lies in $R$ (it is on the ray $\arg z=\pi/2$). Hence
\[
|z|_{\max} = OA = \boxed{6}, \quad \text{attained at } z=6i.
\]

\textbf{Least modulus — discussion.}
Points of the ray $\arg z=\pi/4$ satisfy $z=x+ix$, and
\[
|z-3i|^2 = 2x^2-6x+9 \le 9 \iff 2x(x-3)\le 0 \iff 0\le x\le 3.
\]
So the whole segment from $O$ to $B$ (excluding $O$) lies in $R$, and $|z|=x\sqrt{2}$ can be made \emph{arbitrarily close to $0$}. Since $O\notin R$ (argument undefined), $|z|$ has \textbf{no least value} on $R$; its greatest lower bound (infimum) is $0$, not attained.

\begin{attention}[A classic trap]
One might be tempted to answer $|z|_{\min}=OB=3\sqrt{2}$ at $B=3+3i$. This is wrong: $B$ is merely the closest point to $O$ \emph{on the circular arc}, but the region also contains the segment $OB$ of the ray $\arg z=\pi/4$, along which $|z|$ decreases towards $0$. Always check whether the region accumulates at the origin when the boundary circle passes through $O$.
\end{attention}
\end{exoresolu}

\begin{methode}[Transformations from the $z$-plane to the $w$-plane]
\textbf{Goal:} find the image (locus of $w$) of a given locus of $z$ under a transformation $w = f(z)$, or find a pre-image.

\textbf{General algebraic strategy:}
\begin{enumerate}
    \item Write $w = u+iv$ and $z = x+iy$.
    \item Express $z$ in terms of $w$ (invert the mapping): $z = f^{-1}(w)$.
    \item Substitute $z = f^{-1}(w)$ into the equation of the $z$-locus (e.g. $|z-a|=r$, $\arg(z-a)=\alpha$, or a Cartesian equation).
    \item Simplify to obtain a Cartesian equation in $u,v$ (or a condition on $|w|$, $\arg w$).
    \item Identify the locus in the $w$-plane (line, circle, ray, etc.) and mark any excluded points.
\end{enumerate}

\textbf{Key standard transformations:}
\begin{itemize}
    \item \textbf{Translation:} $w = z + b$: shift by the vector affix $b$.
    \item \textbf{Rotation--enlargement:} $w = az$ ($a \in \mathbb{C}^*$): rotation by $\arg a$ about $O$, enlargement factor $|a|$.
    \item \textbf{Reciprocal:} $w = \dfrac{1}{z}$:
        \begin{itemize}
            \item circle through the origin $\to$ line not through the origin;
            \item line not through the origin $\to$ circle through the origin;
            \item line through the origin $\to$ line through the origin;
            \item circle not through the origin $\to$ circle not through the origin.
        \end{itemize}
    \item \textbf{Möbius (bilinear) transformation:} $w = \dfrac{az+b}{cz+d}$ ($ad-bc \ne 0$): maps generalized circles (lines and circles) to generalized circles.
\end{itemize}

\textbf{Useful formula for $w=1/z$.} If the $z$-locus is the circle $|z-a|=r$ with $|a|\ne r$ (not through the origin), its image is the circle of centre $w_c = \dfrac{\bar{a}}{|a|^2 - r^2}$ and radius $R = \dfrac{r}{\left||a|^2 - r^2\right|}$.
\end{methode}

\begin{exoresolu}[Image of a Line under a Möbius Transformation]
\textbf{Statement:} The transformation $T$ from the $z$-plane to the $w$-plane is defined by
\[
w = \frac{z - i}{z + i}, \quad z \ne -i.
\]
The line $L$ in the $z$-plane has equation $\Im(z) = 1$ (i.e. $y=1$).
\begin{enumerate}
    \item Find the Cartesian equation of the image of $L$ under $T$ in the $w$-plane.
    \item Sketch the image locus, indicating clearly any excluded points.
    \item Find the image of the point $z = 1+i$.
\end{enumerate}

\tcblower
\textbf{Detailed Solution:}

\textbf{1. Equation of the image.}
Invert the mapping:
\[
w(z+i) = z-i \implies wz + wi = z - i \implies z(w-1) = -i(1+w) \implies z = i\,\frac{1+w}{1-w}.
\]
Let $w = u+iv$. Compute $z$ in terms of $u,v$:
\[
z = i\,\frac{(1+u)+iv}{(1-u)-iv}
   = i\,\frac{[(1+u)+iv][(1-u)+iv]}{(1-u)^2 + v^2}
   = i\,\frac{(1-u^2-v^2) + 2iv}{(1-u)^2+v^2}.
\]
Hence
\[
\Re(z) = \frac{-2v}{(1-u)^2+v^2}, \qquad \Im(z) = \frac{1-u^2-v^2}{(1-u)^2+v^2}.
\]
The condition $\Im(z) = 1$ gives
\[
1-u^2-v^2 = (1-u)^2+v^2 = 1 - 2u + u^2 + v^2
\implies 2u^2 + 2v^2 - 2u = 0
\implies u^2 + v^2 - u = 0.
\]
Completing the square:
\[
\boxed{\left(u - \frac{1}{2}\right)^2 + v^2 = \frac{1}{4}},
\]
a \textbf{circle} of centre $\left(\frac{1}{2}, 0\right)$ and radius $\frac{1}{2}$.

\textbf{2. Sketch and excluded points.}
The pole $z=-i$ does not lie on $L$ (since $\Im(-i)=-1\ne 1$), so no point of $L$ is sent to infinity: the image is a genuine circle. However, as $z\to\infty$ along $L$, $w = \frac{z-i}{z+i} \to 1$; since no finite point of $L$ maps to $w=1$, the point $w=1$ is \textbf{excluded} (open dot). Also $z=i\in L$ maps to $w=0$, so the circle passes through the origin.

\begin{popfigure}
\begin{tikzpicture}[scale=2.2]
\draw[->] (-0.5,0) -- (1.5,0) node[right] {$u$};
\draw[->] (0,-0.8) -- (0,0.8) node[above] {$v$};
\draw[popblue, thick] (0.5,0) circle (0.5);
\fill[poporange] (0,0) circle (1.5pt) node[below left] {$0$};
\draw[poporange, fill=white, thick] (1,0) circle (2.5pt) node[below right] {$1$ excluded};
\fill[popdark] (0.5,0) circle (1pt) node[below] {$\left(\frac12,0\right)$};
\draw[popmuted, dashed] (0.5,0) -- (0.5,0.5) node[midway, right, font=\small] {$\frac12$};
\end{tikzpicture}
\legende{Image of $\Im(z)=1$: circle of centre $(1/2,0)$, radius $1/2$; the point $w=1$ is excluded.}
\end{popfigure}

\textbf{3. Image of $z=1+i$.}
\[
w = \frac{(1+i) - i}{(1+i) + i} = \frac{1}{1+2i} = \frac{1-2i}{(1+2i)(1-2i)} = \frac{1-2i}{5} = \boxed{\frac{1}{5} - \frac{2}{5}i}.
\]
Check on the circle: $u=0.2$, $v=-0.4$: $(0.2-0.5)^2 + (-0.4)^2 = 0.09+0.16=0.25 = (0.5)^2$. \checkmark
\end{exoresolu}

\begin{methode}[Matrix Representation of Plane Transformations]
\textbf{Goal:} represent geometric transformations (rotation, reflection, enlargement, stretch, shear, and their composites) as $2\times 2$ matrices acting on column vectors $\begin{pmatrix} x \\ y \end{pmatrix}$.

\textbf{Standard matrices (with respect to the standard basis):}
\begin{center}
\begin{tabularx}{\linewidth}{|l|X|c|}\hline
\rowcolor{popblueL}
\textbf{Transformation} & \textbf{Matrix $\mathbf{M}$} & \textbf{$\det\mathbf{M}$} \\ \hline
Rotation by $\theta$ about $O$ (anticlockwise) & $\begin{pmatrix} \cos\theta & -\sin\theta \\ \sin\theta & \cos\theta \end{pmatrix}$ & $+1$ \\ \hline
Reflection in the line $y = (\tan\theta)x$ & $\begin{pmatrix} \cos 2\theta & \sin 2\theta \\ \sin 2\theta & -\cos 2\theta \end{pmatrix}$ & $-1$ \\ \hline
Enlargement, factor $k$, centre $O$ & $\begin{pmatrix} k & 0 \\ 0 & k \end{pmatrix}$ & $k^2$ \\ \hline
Stretch parallel to the $x$-axis, factor $a$ & $\begin{pmatrix} a & 0 \\ 0 & 1 \end{pmatrix}$ & $a$ \\ \hline
Stretch parallel to the $y$-axis, factor $b$ & $\begin{pmatrix} 1 & 0 \\ 0 & b \end{pmatrix}$ & $b$ \\ \hline
Shear parallel to the $x$-axis, factor $k$ & $\begin{pmatrix} 1 & k \\ 0 & 1 \end{pmatrix}$ & $1$ \\ \hline
Shear parallel to the $y$-axis, factor $k$ & $\begin{pmatrix} 1 & 0 \\ k & 1 \end{pmatrix}$ & $1$ \\ \hline
Reflection in the $x$-axis & $\begin{pmatrix} 1 & 0 \\ 0 & -1 \end{pmatrix}$ & $-1$ \\ \hline
Reflection in the $y$-axis & $\begin{pmatrix} -1 & 0 \\ 0 & 1 \end{pmatrix}$ & $-1$ \\ \hline
Reflection in the line $y=x$ & $\begin{pmatrix} 0 & 1 \\ 1 & 0 \end{pmatrix}$ & $-1$ \\ \hline
Reflection in the line $y=-x$ & $\begin{pmatrix} 0 & -1 \\ -1 & 0 \end{pmatrix}$ & $-1$ \\ \hline
\end{tabularx}
\end{center}

\textbf{Composition:} ``$T_1$ then $T_2$'' has matrix $\mathbf{M}_2 \mathbf{M}_1$ (right-to-left order).
\textbf{Inverse:} $\mathbf{M}^{-1}$ exists iff $\det\mathbf{M} \ne 0$.
\textbf{Invariant points:} solve $\mathbf{M}\mathbf{x} = \mathbf{x}$. \textbf{Invariant lines through $O$:} solve $\mathbf{M}\mathbf{x} = \lambda \mathbf{x}$ (eigenvectors).
\end{methode}

\begin{exoresolu}[Composite Transformation: Rotation followed by Reflection]
\textbf{Statement:} $R$ is the rotation of $60^\circ$ anticlockwise about the origin; $S$ is the reflection in the line $y = \sqrt{3}\,x$. Let $T = S \circ R$ (apply $R$ first, then $S$).
\begin{enumerate}
    \item Write down the matrices $\mathbf{R}$ and $\mathbf{S}$.
    \item Find the matrix $\mathbf{T}$ representing $T$.
    \item Identify the single geometric transformation represented by $\mathbf{T}$, describing it fully.
    \item Find the equation of the line of invariant points of $T$.
\end{enumerate}

\tcblower
\textbf{Detailed Solution:}

\textbf{1. Matrices $\mathbf{R}$ and $\mathbf{S}$.}
$\theta = 60^\circ$: $\cos 60^\circ = \tfrac12$, $\sin 60^\circ = \tfrac{\sqrt3}{2}$.
\[
\mathbf{R} = \begin{pmatrix} \frac{1}{2} & -\frac{\sqrt{3}}{2} \\[2pt] \frac{\sqrt{3}}{2} & \frac{1}{2} \end{pmatrix}.
\]
The mirror line $y = \sqrt{3}x$ has $\tan\theta = \sqrt{3}$, so $\theta = 60^\circ$ and $2\theta = 120^\circ$: $\cos 120^\circ = -\tfrac12$, $\sin 120^\circ = \tfrac{\sqrt3}{2}$.
\[
\mathbf{S} = \begin{pmatrix} -\frac{1}{2} & \frac{\sqrt{3}}{2} \\[2pt] \frac{\sqrt{3}}{2} & \frac{1}{2} \end{pmatrix}.
\]

\textbf{2. Matrix $\mathbf{T} = \mathbf{S}\mathbf{R}$.}
\[
\mathbf{T} = \begin{pmatrix} -\frac{1}{2} & \frac{\sqrt{3}}{2} \\[2pt] \frac{\sqrt{3}}{2} & \frac{1}{2} \end{pmatrix} \begin{pmatrix} \frac{1}{2} & -\frac{\sqrt{3}}{2} \\[2pt] \frac{\sqrt{3}}{2} & \frac{1}{2} \end{pmatrix}.
\]
Entries:
$T_{11} = -\frac{1}{4} + \frac{3}{4} = \frac{1}{2}$;\quad
$T_{12} = \frac{\sqrt{3}}{4} + \frac{\sqrt{3}}{4} = \frac{\sqrt{3}}{2}$;\quad
$T_{21} = \frac{\sqrt{3}}{4} + \frac{\sqrt{3}}{4} = \frac{\sqrt{3}}{2}$;\quad
$T_{22} = -\frac{3}{4} + \frac{1}{4} = -\frac{1}{2}$.
\[
\mathbf{T} = \begin{pmatrix} \frac{1}{2} & \frac{\sqrt{3}}{2} \\[2pt] \frac{\sqrt{3}}{2} & -\frac{1}{2} \end{pmatrix}.
\]

\textbf{3. Identification.}
$\det \mathbf{T} = \left(\frac{1}{2}\right)\left(-\frac{1}{2}\right) - \left(\frac{\sqrt{3}}{2}\right)^{\!2} = -\frac{1}{4} - \frac{3}{4} = -1$, so $T$ reverses orientation. The matrix matches the reflection pattern $\begin{pmatrix} \cos 2\alpha & \sin 2\alpha \\ \sin 2\alpha & -\cos 2\alpha \end{pmatrix}$ with
\[
\cos 2\alpha = \tfrac12,\quad \sin 2\alpha = \tfrac{\sqrt3}{2} \implies 2\alpha = 60^\circ \implies \alpha = 30^\circ.
\]
\textbf{Conclusion:} $T$ is the \textbf{reflection in the line $y = (\tan 30^\circ)\,x = \dfrac{x}{\sqrt{3}}$}.

\textbf{4. Line of invariant points.}
Solve $\mathbf{T}\mathbf{x} = \mathbf{x}$:
\[
\tfrac{1}{2}x + \tfrac{\sqrt{3}}{2}y = x \implies x = \sqrt{3}\,y ; \qquad
\tfrac{\sqrt{3}}{2}x - \tfrac{1}{2}y = y \implies x = \sqrt{3}\,y \quad (\text{consistent}).
\]
Equation of the invariant line: $\boxed{y = \dfrac{x}{\sqrt{3}}}$ \; (equivalently $x - \sqrt{3}\,y = 0$) — exactly the mirror line found above, as expected for a reflection.
\end{exoresolu}

\begin{methode}[Similitudes: Complex Form, Matrix Form and Properties]
\textbf{Goal:} recognise, analyse and apply similitudes (direct and opposite) in both complex and matrix forms.

\textbf{Definitions.} A \textbf{similitude} is a transformation that preserves angles (hence shapes) and multiplies all distances by a fixed ratio $k>0$.
\begin{itemize}
    \item \textbf{Direct similitude (orientation-preserving):} $z' = az + b$, $a,b \in \mathbb{C}$, $a \ne 0$. Ratio $k = |a|$, angle of rotation $\theta = \arg a$.
    Matrix (affine) form:
    \[
    \begin{pmatrix} x' \\ y' \end{pmatrix} = \begin{pmatrix} \Re(a) & -\Im(a) \\ \Im(a) & \Re(a) \end{pmatrix} \begin{pmatrix} x \\ y \end{pmatrix} + \begin{pmatrix} \Re(b) \\ \Im(b) \end{pmatrix}.
    \]
    \item \textbf{Opposite similitude (orientation-reversing):} $z' = a\bar{z} + b$, $a \ne 0$. Ratio $k = |a|$.
    Matrix (affine) form:
    \[
    \begin{pmatrix} x' \\ y' \end{pmatrix} = \begin{pmatrix} \Re(a) & \Im(a) \\ \Im(a) & -\Re(a) \end{pmatrix} \begin{pmatrix} x \\ y \end{pmatrix} + \begin{pmatrix} \Re(b) \\ \Im(b) \end{pmatrix}.
    \]
\end{itemize}

\textbf{Key properties:}
\begin{enumerate}
    \item \textbf{Fixed point (centre):} a direct similitude $z' = az+b$ with $a\ne 1$ has a unique fixed point $z_0 = \dfrac{b}{1-a}$: it is the centre of the rotation--enlargement.
    \item \textbf{Composition:} direct $\circ$ direct $=$ direct (ratios multiply, angles add); opposite $\circ$ opposite $=$ direct; direct $\circ$ opposite $=$ opposite.
    \item \textbf{Inverse:} the inverse of a direct similitude is direct; the inverse of an opposite similitude is opposite.
\end{enumerate}

\textbf{Determining a direct similitude from two pairs $(z_1, z_1')$, $(z_2, z_2')$:}
\[
a = \frac{z_1' - z_2'}{z_1 - z_2}, \qquad b = z_1' - a z_1, \qquad z_0 = \frac{b}{1-a}.
\]
\end{methode}

\begin{exoresolu}[Analysing a Direct Similitude from Point Pairs]
\textbf{Statement:} A direct similitude $S$ maps $z_1 = 1+i$ to $z_1' = 4+2i$ and $z_2 = 2-i$ to $z_2' = 5-4i$.
\begin{enumerate}
    \item Determine the complex form $z' = az + b$ of $S$.
    \item Find the ratio of similitude $k$ and the angle of rotation $\theta$.
    \item Find the centre (fixed point) $z_0$ of $S$.
    \item Write the matrix (affine) representation of $S$.
    \item Determine the Cartesian equation of the image of the line $\Re(z) = 1$.
\end{enumerate}

\tcblower
\textbf{Detailed Solution:}

\textbf{1. Find $a$ and $b$.}
\[
a = \frac{z_1' - z_2'}{z_1 - z_2} = \frac{(4+2i) - (5-4i)}{(1+i) - (2-i)} = \frac{-1 + 6i}{-1 + 2i}
= \frac{(-1+6i)(-1-2i)}{(-1)^2+2^2} = \frac{13 - 4i}{5} = 2.6 - 0.8i.
\]
Then, using $z_1$:
\[
a z_1 = (2.6 - 0.8i)(1+i) = 2.6 + 2.6i - 0.8i - 0.8i^2 = 3.4 + 1.8i,
\]
\[
b = z_1' - a z_1 = (4+2i) - (3.4 + 1.8i) = 0.6 + 0.2i.
\]
\[
\boxed{z' = (2.6 - 0.8i)\,z + (0.6 + 0.2i)} \qquad \left(\text{equivalently } z' = \tfrac{1}{5}\big[(13-4i)z + (3+i)\big]\right).
\]
\emph{Check with $z_2$:} $az_2 = (2.6-0.8i)(2-i) = 5.2 - 2.6i - 1.6i + 0.8i^2 = 4.4 - 4.2i$; then $az_2 + b = 5 - 4i = z_2'$. \checkmark

\textbf{2. Ratio $k$ and angle $\theta$.}
\[
k = |a| = \frac{|13-4i|}{5} = \frac{\sqrt{169+16}}{5} = \boxed{\frac{\sqrt{185}}{5}} \approx 2.72.
\]
$\tan\theta = \dfrac{-4}{13}$ with $\Re(a)>0$, $\Im(a)<0$ (fourth quadrant):
\[
\theta = -\arctan\!\frac{4}{13} \approx -17.1^\circ .
\]

\textbf{3. Centre $z_0$.}
\[
1-a = \frac{5-(13-4i)}{5} = \frac{-8+4i}{5}, \qquad
z_0 = \frac{b}{1-a} = \frac{3+i}{-8+4i} = \frac{(3+i)(-8-4i)}{64+16} = \frac{-20-20i}{80} = \boxed{-\frac14 - \frac14 i}.
\]
\emph{Check:} $az_0 = (2.6-0.8i)(-0.25-0.25i) = -0.85 - 0.45i$; $az_0 + b = -0.25 - 0.25i = z_0$. \checkmark

\textbf{4. Matrix representation.}
With $\Re(a)=2.6$, $\Im(a)=-0.8$:
\[
\boxed{\begin{pmatrix} x' \\ y' \end{pmatrix} = \begin{pmatrix} 2.6 & 0.8 \\ -0.8 & 2.6 \end{pmatrix} \begin{pmatrix} x \\ y \end{pmatrix} + \begin{pmatrix} 0.6 \\ 0.2 \end{pmatrix}}.
\]

\textbf{5. Image of the line $x=1$.}
Write $a=\alpha+i\beta$ with $\alpha=2.6$, $\beta=-0.8$, and $b = u_0+iv_0$ with $u_0=0.6$, $v_0=0.2$. For $z=x+iy$:
\[
u = \alpha x - \beta y + u_0, \qquad v = \beta x + \alpha y + v_0.
\]
Setting $x=1$:
\[
u = 2.6 + 0.8y + 0.6 = 3.2 + 0.8y, \qquad v = -0.8 + 2.6y + 0.2 = -0.6 + 2.6y.
\]
Eliminate $y$: from the first equation $y = \dfrac{u-3.2}{0.8}$; substituting,
\[
v = -0.6 + 2.6\cdot\frac{u-3.2}{0.8} = -0.6 + 3.25(u-3.2) = \boxed{3.25\,u - 11}.
\]
The image is a straight line (as expected: a similitude maps lines to lines). \emph{Consistency check:} the original line is vertical (direction angle $90^\circ$); after rotation by $\theta\approx-17.1^\circ$ the direction angle is $\approx 72.9^\circ$, and $\tan 72.9^\circ \approx 3.25$ — matching the slope found.
\end{exoresolu}

\begin{methode}[Solving Modulus Inequalities and Classifying Rigid Motions]
\textbf{Goal:} solve inequalities involving $|z-a| \lesseqgtr |z-b|$ and interpret them geometrically; classify rigid motions (isometries) of the plane.

\textbf{Modulus inequalities — algebraic method.} For $|z-a| \lesseqgtr |z-b|$, write $z=x+iy$, $a=x_a+iy_a$, $b=x_b+iy_b$ and square both sides (valid since both sides are $\ge 0$):
\[
(x-x_a)^2 + (y-y_a)^2 \lesseqgtr (x-x_b)^2 + (y-y_b)^2.
\]
The terms $x^2, y^2$ cancel, leaving the linear inequality
\[
2(x_b - x_a)\,x + 2(y_b - y_a)\,y \lesseqgtr |b|^2 - |a|^2,
\]
a \textbf{half-plane} bounded by the perpendicular bisector of the segment $AB$.

\textbf{Rigid motions (isometries):} transformations preserving distances, $|z_1' - z_2'| = |z_1 - z_2|$.
\begin{itemize}
    \item \textbf{Direct (orientation-preserving, $\det=+1$):}
        \begin{itemize}
            \item Translation $z' = z + b$ ($b \ne 0$): no fixed point.
            \item Rotation $z' - c = e^{i\theta}(z - c)$ ($\theta \ne 2k\pi$): unique fixed point, the centre $c$.
        \end{itemize}
    \item \textbf{Opposite (orientation-reversing, $\det=-1$):}
        \begin{itemize}
            \item Reflection: a line of fixed points (the mirror).
            \item Glide reflection (reflection composed with a translation parallel to the mirror): no fixed point.
        \end{itemize}
\end{itemize}

\textbf{Classification algorithm} for $\mathbf{x}' = \mathbf{M}\mathbf{x} + \mathbf{t}$ with $\mathbf{M}$ orthogonal ($\mathbf{M}^{T}\mathbf{M}=\mathbf{I}$):
\begin{enumerate}
    \item Compute $\det \mathbf{M}$.
    \item If $\det \mathbf{M} = +1$ (direct): if $\mathbf{M} = \mathbf{I}$, it is a \textbf{translation} by $\mathbf{t}$; otherwise it is a \textbf{rotation}, with centre $\mathbf{c} = (\mathbf{I}-\mathbf{M})^{-1}\mathbf{t}$ and angle read from $\mathbf{M}$.
    \item If $\det \mathbf{M} = -1$ (opposite): solve $(\mathbf{I}-\mathbf{M})\mathbf{x} = \mathbf{t}$. If there is a line of solutions, it is a \textbf{reflection} (the solution line is the mirror); if there is no solution, it is a \textbf{glide reflection}.
\end{enumerate}
\end{methode}

\begin{exoresolu}[Modulus Inequality Region and Rigid Motion Classification]
\textbf{Statement:}
\begin{enumerate}
    \item Solve the inequality $|z - 3| < |z - i|$ and shade the region on an Argand diagram.
    \item A transformation $T$ is defined by
    \[
    \begin{pmatrix} x' \\ y' \end{pmatrix} = \begin{pmatrix} 0 & -1 \\ -1 & 0 \end{pmatrix} \begin{pmatrix} x \\ y \end{pmatrix} + \begin{pmatrix} 2 \\ 2 \end{pmatrix}.
    \]
    \begin{enumerate}
        \item Show that $T$ is an isometry (rigid motion).
        \item Classify $T$ fully (type and mirror line).
        \item Find the image of the circle $|z| = 2$ under $T$.
    \end{enumerate}
\end{enumerate}

\tcblower
\textbf{Detailed Solution:}

\textbf{1. Inequality $|z-3| < |z-i|$.}
\textbf{Geometric reading:} points closer to $A(3,0)$ than to $B(0,1)$; the boundary is the perpendicular bisector of $AB$.
Midpoint of $AB$: $M(1.5,\,0.5)$. Gradient of $AB$: $\dfrac{1-0}{0-3} = -\dfrac13$, so the bisector has gradient $3$:
\[
y - 0.5 = 3(x - 1.5) \implies y = 3x - 4.
\]
\textbf{Algebraic check} with $z=x+iy$:
\[
(x-3)^2 + y^2 < x^2 + (y-1)^2
\implies -6x + 9 < -2y + 1
\implies \boxed{y < 3x - 4}.
\]
Test the origin: $|0-3|=3 > |0-i|=1$, so $O$ is \emph{not} in the region; the region is the half-plane below the line (the side containing $A(3,0)$). Boundary dashed (strict inequality).

\begin{popfigure}
\begin{tikzpicture}[scale=0.75]
\draw[->] (-2,0) -- (5,0) node[right] {$\Re$};
\draw[->] (0,-3) -- (0,5) node[above] {$\Im$};
% Shaded half-plane y < 3x - 4, clipped to the viewing window
\begin{scope}
\clip (-2,-3) rectangle (5,5);
\fill[popblueL] (1.333,-3) -- (3,5) -- (5,5) -- (5,-3) -- cycle;
\draw[popblue, dashed, thick] (0.333,-3) -- (3,5);
\end{scope}
\fill[poporange] (3,0) circle (2pt) node[below right] {$A=3$};
\fill[poporange] (0,1) circle (2pt) node[above left] {$B=i$};
\fill[popdark] (1.5,0.5) circle (2pt) node[above right, font=\small] {$M$};
\node[popblue, font=\small] at (3.6,1.2) {$y=3x-4$};
\end{tikzpicture}
\legende{Region $|z-3| < |z-i|$: open half-plane below the bisector $y=3x-4$ (dashed).}
\end{popfigure}

\textbf{2(a). Isometry check.}
\[
\mathbf{M}^{T}\mathbf{M} = \begin{pmatrix} 0 & -1 \\ -1 & 0 \end{pmatrix} \begin{pmatrix} 0 & -1 \\ -1 & 0 \end{pmatrix} = \begin{pmatrix} 1 & 0 \\ 0 & 1 \end{pmatrix} = \mathbf{I},
\]
so $\mathbf{M}$ is orthogonal: $T$ preserves distances and is a rigid motion.

\textbf{2(b). Classification.}
$\det \mathbf{M} = (0)(0) - (-1)(-1) = -1$: $T$ is an \textbf{opposite} isometry.
Fixed points: solve $(\mathbf{I}-\mathbf{M})\mathbf{x} = \mathbf{t}$:
\[
\begin{pmatrix} 1 & 1 \\ 1 & 1 \end{pmatrix} \begin{pmatrix} x \\ y \end{pmatrix} = \begin{pmatrix} 2 \\ 2 \end{pmatrix} \implies x + y = 2,
\]
a whole line of fixed points. Hence $T$ is a \textbf{reflection} in the mirror line
\[
\boxed{x + y = 2}.
\]
\emph{Consistency:} $\mathbf{M}$ alone is the reflection in $y=-x$; the translation $(2,2)$ is perpendicular to that line, and composing gives the reflection in the parallel line through $(1,1)$, i.e. $x+y=2$. \checkmark

\textbf{2(c). Image of the circle $|z|=2$.}
An isometry maps circles to circles of the same radius. The centre $O(0,0)$ maps to
\[
\mathbf{M}\begin{pmatrix}0\\0\end{pmatrix} + \begin{pmatrix}2\\2\end{pmatrix} = \begin{pmatrix}2\\2\end{pmatrix}.
\]
Hence the image is the circle of centre $(2,2)$ and radius $2$:
\[
\boxed{|z - (2+2i)| = 2} \qquad \big(\text{i.e. } (x-2)^2 + (y-2)^2 = 4\big).
\]
\end{exoresolu}

\newpage

\coursec{Exercises}

\courssub{I check my knowledge}
\begin{exercice}[MCQ: Forms and De Moivre]
Which of the following is the exponential form of $z = 1 + i\sqrt{3}$?
\begin{enumerate}
\item $2e^{i\pi/3}$
\item $2e^{i\pi/6}$
\item $4e^{i\pi/3}$
\item $2e^{-i\pi/3}$
\end{enumerate}
\end{exercice}

\begin{exercice}[True/False with justification]
State whether each statement is true or false. Justify your answer.
\begin{enumerate}
\item $\arg(z_1 z_2) = \arg(z_1) + \arg(z_2)$ for all non-zero $z_1, z_2 \in \mathbb{C}$.
\item The locus $|z - 2| = |z + 2i|$ is a straight line.
\item The transformation $w = \frac{1}{z}$ maps the line $\operatorname{Re}(z) = 2$ to a circle.
\item A similitude $z' = az + b$ with $|a| = 1$ and $a \neq 1$ is a rotation.
\end{enumerate}
\end{exercice}

\begin{exercice}[Definitions and direct calculations]
\begin{enumerate}
\item Write the complex number $z = -2 + 2i$ in trigonometric form $r(\cos\theta + i\sin\theta)$ and exponential form $re^{i\theta}$, with $r>0$ and $\theta \in (-\pi,\pi]$.
\item Use De Moivre's theorem to compute $(1 - i)^8$ exactly.
\item Find all the cube roots of $-8i$, giving them in exponential form.
\end{enumerate}
\end{exercice}

\begin{exercice}[Quick computation]
Given $z_1 = 3e^{i\pi/4}$ and $z_2 = 2e^{-i\pi/6}$, compute exactly:
\begin{enumerate}
\item $z_1 z_2$ and $\frac{z_1}{z_2}$ in exponential form.
\item $\operatorname{Re}(z_1 \overline{z_2})$ and $\operatorname{Im}(z_1 \overline{z_2})$.
\end{enumerate}
\end{exercice}

\courssub{I practise}
\begin{exercice}[Trigonometric and exponential forms]
Let $z = 1 + i\sqrt{3}$.
\begin{enumerate}
\item Write $z$ in trigonometric form and in exponential form.
\item Compute $z^{12}$ exactly.
\item Find all the fourth roots of $-4z$. Plot them on an Argand diagram.
\end{enumerate}
\end{exercice}

\begin{exercice}[De Moivre and $\cos 5\theta$]
\begin{enumerate}
\item Use De Moivre's theorem to show that $\cos 5\theta = 16\cos^5\theta - 20\cos^3\theta + 5\cos\theta$.
\item Deduce the exact value of $\cos(\pi/10)$ as a surd.
\end{enumerate}
\end{exercice}

\begin{exercice}[Loci in the Argand diagram]
On separate Argand diagrams, sketch the following loci:
\begin{enumerate}
\item $|z - 2 + i| = 2$.
\item $|z - 2| = |z + 2i|$.
\item $\arg(z - 1) = \pi/4$.
\end{enumerate}
On a single diagram, shade the region satisfying all three conditions simultaneously.
\end{exercice}

\begin{exercice}[Transformation $w = 1/z$]
Find the image of the line $\operatorname{Re}(z) = 2$ under the transformation $w = 1/z$. Give the centre and radius of the resulting circle.
\end{exercice}

\courssub{I prepare for the GCE}
\begin{exercice}[{Composite transformation and matrix] [12 marks}]
A transformation $T_1$ is represented by the matrix $M = \begin{pmatrix} 0 & -1 \\ 1 & 0 \end{pmatrix}$ and a transformation $T_2$ by $N = \begin{pmatrix} 2 & 0 \\ 0 & 2 \end{pmatrix}$.
\begin{enumerate}
\item Find the single matrix representing the composite transformation $T_2 \circ T_1$.
\item State the geometric effect of this composite transformation.
\item Determine its area scale factor.
\item Find the inverse transformation and its matrix.
\end{enumerate}
\end{exercice}

\begin{exercice}[{Similitude analysis] [13 marks}]
A similitude is given by $z' = (1 - i)z + 3i$.
\begin{enumerate}
\item Determine its ratio $k$ and its angle of rotation $\theta$.
\item Find its centre $c$.
\item Classify the similitude (rotation, homothety, translation, etc.).
\item Find the image of the triangle with vertices $0$, $1$, $i$.
\end{enumerate}
\end{exercice}

\begin{exercice}[{Modulus inequality and rigid motions] [15 marks}]
\begin{enumerate}
\item Prove the triangle inequality $|z_1 + z_2| \le |z_1| + |z_2|$ for all $z_1, z_2 \in \mathbb{C}$.
\item Solve in $\mathbb{C}$ the inequality $|z - 3i| \le |z + 3|$. Shade the solution set on an Argand diagram.
\item Classify the rigid motions:
\begin{enumerate}
\item $z' = iz + 2 - i$,
\item $z' = i\overline{z} + 1$.
\end{enumerate}
For each, state whether it preserves or reverses orientation and give its characteristic elements (centre and angle, or axis and vector).
\end{enumerate}
\end{exercice}

\newpage

\coursec{Solutions to the exercises}

\coursec{Worked Exercises and Examination Practice}

\begin{corrige}[Exercise 1 — MCQ: Forms and De Moivre]
The correct answer is \textbf{a)} $2e^{i\pi/3}$.

\textbf{Justification.} For $z = 1 + i\sqrt{3}$:
\begin{itemize}
\item Modulus: $|z| = \sqrt{1^2 + (\sqrt{3})^2} = \sqrt{1+3} = 2$.
\item Argument: $\cos\theta = \dfrac{1}{2}$ and $\sin\theta = \dfrac{\sqrt{3}}{2}$, so $\theta = \dfrac{\pi}{3}$ (principal value).
\end{itemize}
Hence $z = 2\left(\cos\dfrac{\pi}{3} + i\sin\dfrac{\pi}{3}\right) = 2e^{i\pi/3}$.

Options b) and d) have the wrong argument, and c) has the wrong modulus ($4 = |z|^2$, a classic trap).
\end{corrige}

\begin{corrige}[Exercise 2 — True/False with justification]
\textbf{a) TRUE (modulo $2\pi$).} If $z_1 = r_1 e^{i\theta_1}$ and $z_2 = r_2 e^{i\theta_2}$, then $z_1 z_2 = r_1 r_2 e^{i(\theta_1+\theta_2)}$, so $\arg(z_1 z_2) = \arg z_1 + \arg z_2 \pmod{2\pi}$. The statement is true provided we accept equality up to a multiple of $2\pi$; if principal values in $(-\pi,\pi]$ are required, the sum may need adjusting by $\pm 2\pi$.

\textbf{b) TRUE.} $|z-2| = |z+2i| = |z-(-2i)|$ says that $P(z)$ is equidistant from $A(2,0)$ and $B(0,-2)$. The locus of points equidistant from two fixed points is the \cle{perpendicular bisector} of $[AB]$, which is a straight line. Algebraically, with $z = x+iy$: $(x-2)^2 + y^2 = x^2 + (y+2)^2$ gives $-4x = 4y$, i.e. $y = -x$.

\textbf{c) TRUE.} The line $\operatorname{Re}(z) = 2$ does not pass through the origin, and inversion $w = 1/z$ maps any line not through $O$ to a circle through $O$. Explicitly (see Exercise 8), the image is the circle $\left|w - \tfrac14\right| = \tfrac14$.

\textbf{d) TRUE.} $z' = az + b$ with $|a| = 1$, $a \neq 1$, has a unique fixed point $c = \dfrac{b}{1-a}$, and writing $z' - c = a(z - c)$ shows it is the rotation of centre $c$ and angle $\arg(a)$. (If $a = 1$ it would be a translation.)
\end{corrige}

\begin{corrige}[Exercise 3 — Definitions and direct calculations]
\textbf{a)} $z = -2 + 2i$: $|z| = \sqrt{4+4} = 2\sqrt{2}$. Then $z = 2\sqrt{2}\left(-\dfrac{1}{\sqrt{2}} + i\dfrac{1}{\sqrt{2}}\right)$, so $\cos\theta = -\frac{\sqrt2}{2}$, $\sin\theta = \frac{\sqrt2}{2}$, giving $\theta = \dfrac{3\pi}{4}$ (in $(-\pi,\pi]$).
\[ z = 2\sqrt{2}\left(\cos\tfrac{3\pi}{4} + i\sin\tfrac{3\pi}{4}\right) = 2\sqrt{2}\,e^{3i\pi/4}. \]

\textbf{b)} $1 - i = \sqrt{2}\,e^{-i\pi/4}$. By De Moivre:
\[ (1-i)^8 = (\sqrt{2})^8 e^{-8i\pi/4} = 16\,e^{-2i\pi} = 16. \]

\textbf{c)} $-8i = 8e^{-i\pi/2} = 8e^{i( -\pi/2 + 2k\pi)}$. The cube roots are
\[ z_k = 2\,e^{i(-\pi/6 + 2k\pi/3)}, \quad k = 0, 1, 2, \]
i.e. $z_0 = 2e^{-i\pi/6}$, $z_1 = 2e^{i\pi/2} = 2i$, $z_2 = 2e^{7i\pi/6} = 2e^{-5i\pi/6}$.
\end{corrige}

\begin{corrige}[Exercise 4 — Quick computation]
\textbf{a)} $z_1 z_2 = 3 \cdot 2 \, e^{i(\pi/4 - \pi/6)} = 6e^{i\pi/12}$. \quad $\dfrac{z_1}{z_2} = \dfrac{3}{2}e^{i(\pi/4 + \pi/6)} = \dfrac{3}{2}e^{5i\pi/12}$.

\textbf{b)} $z_1\overline{z_2} = 3e^{i\pi/4}\cdot 2e^{i\pi/6} = 6e^{5i\pi/12}$. Now $\dfrac{5\pi}{12} = \dfrac{\pi}{4} + \dfrac{\pi}{6}$, so
\[ \cos\tfrac{5\pi}{12} = \cos\tfrac{\pi}{4}\cos\tfrac{\pi}{6} - \sin\tfrac{\pi}{4}\sin\tfrac{\pi}{6} = \tfrac{\sqrt2}{2}\cdot\tfrac{\sqrt3}{2} - \tfrac{\sqrt2}{2}\cdot\tfrac12 = \tfrac{\sqrt6 - \sqrt2}{4}, \]
\[ \sin\tfrac{5\pi}{12} = \sin\tfrac{\pi}{4}\cos\tfrac{\pi}{6} + \cos\tfrac{\pi}{4}\sin\tfrac{\pi}{6} = \tfrac{\sqrt6 + \sqrt2}{4}. \]
Therefore
\[ \operatorname{Re}(z_1\overline{z_2}) = 6\cdot\tfrac{\sqrt6-\sqrt2}{4} = \tfrac{3(\sqrt6-\sqrt2)}{2}, \qquad \operatorname{Im}(z_1\overline{z_2}) = \tfrac{3(\sqrt6+\sqrt2)}{2}. \]
\end{corrige}

\begin{corrige}[Exercise 5 — Trigonometric and exponential forms]
\textbf{a)} From Exercise 1: $z = 1 + i\sqrt3 = 2\left(\cos\dfrac{\pi}{3} + i\sin\dfrac{\pi}{3}\right) = 2e^{i\pi/3}$.

\textbf{b)} By De Moivre: $z^{12} = 2^{12}e^{12i\pi/3} = 4096\,e^{4i\pi} = 4096$.

\textbf{c)} $-4z = -4(1+i\sqrt3)$. Modulus: $|-4z| = 4|z| = 8$. Argument: $-4z$ is $z$ rotated by $\pi$ (multiplication by the negative real $-4$), so $\arg(-4z) = \dfrac{\pi}{3} + \pi = \dfrac{4\pi}{3} \equiv -\dfrac{2\pi}{3}$. Hence $-4z = 8e^{-2i\pi/3}$.

The fourth roots are $w_k = 8^{1/4} e^{i(-2\pi/3 + 2k\pi)/4} = 2^{3/4} e^{i(-\pi/6 + k\pi/2)}$, $k = 0,1,2,3$:
\[ w_0 = 2^{3/4}e^{-i\pi/6},\quad w_1 = 2^{3/4}e^{i\pi/3},\quad w_2 = 2^{3/4}e^{5i\pi/6},\quad w_3 = 2^{3/4}e^{4i\pi/3} = 2^{3/4}e^{-2i\pi/3}. \]

\begin{popfigure}
\begin{center}
\begin{tikzpicture}[scale=1.1]
\draw[->,popline] (-2.4,0) -- (2.4,0) node[right]{\small Re};
\draw[->,popline] (0,-2.4) -- (0,2.4) node[above]{\small Im};
\draw[popmuted,dashed] (0,0) circle (1.682);
\foreach \a/\n in {-30/w_0, 60/w_1, 150/w_2, -120/w_3}{
  \fill[popblue] (\a:1.682) circle (1.6pt);
  \draw[->,poporange] (0,0) -- (\a:1.682);
  \node[popdark] at (\a:2.05) {\small $\n$};
}
\end{tikzpicture}
\end{center}
\legende{The four fourth roots of $-4z$: vertices of a square on the circle of radius $2^{3/4}$.}
\end{popfigure}

The roots form a square inscribed in the circle of radius $2^{3/4} \approx 1.68$ centred at $O$, successive roots being separated by an angle $\dfrac{\pi}{2}$.
\end{corrige}

\begin{corrige}[Exercise 6 — De Moivre and $\cos 5\theta$]
\textbf{a)} By De Moivre, $\cos 5\theta + i\sin 5\theta = (\cos\theta + i\sin\theta)^5$. Writing $c = \cos\theta$, $s = \sin\theta$ and expanding by the binomial theorem:
\[ (c+is)^5 = c^5 + 5ic^4s - 10c^3s^2 - 10ic^2s^3 + 5cs^4 + is^5. \]
Equating real parts: $\cos 5\theta = c^5 - 10c^3s^2 + 5cs^4$. Using $s^2 = 1 - c^2$:
\begin{align*}
\cos 5\theta &= c^5 - 10c^3(1-c^2) + 5c(1-c^2)^2 \\
&= c^5 - 10c^3 + 10c^5 + 5c - 10c^3 + 5c^5 \\
&= 16c^5 - 20c^3 + 5c.
\end{align*}
Hence $\cos 5\theta = 16\cos^5\theta - 20\cos^3\theta + 5\cos\theta$. $\blacksquare$

\textbf{b)} Put $\theta = \dfrac{\pi}{10}$: then $5\theta = \dfrac{\pi}{2}$, so $\cos 5\theta = 0$. With $c = \cos\dfrac{\pi}{10} \neq 0$:
\[ 16c^5 - 20c^3 + 5c = 0 \;\Longrightarrow\; 16c^4 - 20c^2 + 5 = 0. \]
This is quadratic in $c^2$: $c^2 = \dfrac{20 \pm \sqrt{400 - 320}}{32} = \dfrac{20 \pm \sqrt{80}}{32} = \dfrac{5 \pm \sqrt5}{8}$.

Since $\dfrac{\pi}{10} < \dfrac{\pi}{4}$, we have $\cos\dfrac{\pi}{10} > \cos\dfrac{\pi}{4} = \dfrac{1}{\sqrt2}$, so $c^2 > \dfrac12$. Check: $\dfrac{5-\sqrt5}{8} \approx 0.345 < \tfrac12$ and $\dfrac{5+\sqrt5}{8} \approx 0.905 > \tfrac12$. Hence $c^2 = \dfrac{5+\sqrt5}{8}$, and since $c > 0$:
\[ \cos\frac{\pi}{10} = \sqrt{\frac{5+\sqrt5}{8}} = \frac{\sqrt{10 + 2\sqrt5}}{4}. \]
\end{corrige}

\begin{corrige}[Exercise 7 — Loci in the Argand diagram]
\textbf{a)} $|z - (2 - i)| = 2$: circle of centre $C(2,-1)$ and radius $2$.

\textbf{b)} $|z - 2| = |z + 2i|$: perpendicular bisector of $A(2,0)$ and $B(0,-2)$, i.e. the line $y = -x$.

\textbf{c)} $\arg(z-1) = \dfrac{\pi}{4}$: half-line (ray) from the point $(1,0)$ — excluded — making an angle $\dfrac{\pi}{4}$ with the positive real axis, i.e. the ray $y = x - 1$, $x > 1$.

\begin{popfigure}
\begin{center}
\begin{tikzpicture}[scale=0.85]
\draw[->,popline] (-3,0) -- (4.5,0) node[right]{\small Re};
\draw[->,popline] (0,-3.5) -- (0,3.5) node[above]{\small Im};
% circle
\draw[popblue,thick] (2,-1) circle (2);
\fill[popblue] (2,-1) circle (1.5pt) node[below right]{\small $C(2,-1)$};
% bisector y=-x
\draw[poporange,thick] (-2.5,2.5) -- (3,-3) node[below left]{\small $y=-x$};
\fill[popdark] (2,0) circle (1.5pt) node[above right]{\small $A$};
\fill[popdark] (0,-2) circle (1.5pt) node[below right]{\small $B$};
% ray
\draw[popgreen,thick] (1,0) -- (4,3);
\draw[popgreen,fill=white] (1,0) circle (1.8pt);
\node[popgreen] at (3.9,2.5) {\small $\arg(z-1)=\frac{\pi}{4}$};
\end{tikzpicture}
\end{center}
\legende{The three loci: circle, perpendicular bisector, and half-line (open point at $(1,0)$).}
\end{popfigure}

\textbf{Simultaneous region.} The three conditions are equalities, so the solution set is the set of points lying on all three loci at once. The ray $y = x - 1$ meets the line $y = -x$ where $x - 1 = -x$, i.e. $x = \tfrac12$, which does not satisfy $x > 1$: the ray and the bisector do not intersect. Hence \textbf{no point satisfies all three conditions simultaneously}; the region is empty. (If the intended question was the region defined by the corresponding inequalities — inside the circle, the half-plane containing $O$ bounded by $y=-x$, and the sector $\arg(z-1) \in [\frac{\pi}{4}, \frac{\pi}{2}]$ say — one shades the intersection of the disc, the half-plane and the angular sector.)
\end{corrige}

\begin{corrige}[Exercise 8 — Transformation $w = 1/z$]
Let $z = x + iy$ on the line $\operatorname{Re}(z) = 2$, so $z = 2 + iy$, and write $w = u + iv = \dfrac{1}{z}$, i.e. $z = \dfrac{1}{w}$.

Then $\operatorname{Re}(z) = \operatorname{Re}\left(\dfrac{1}{w}\right) = \operatorname{Re}\left(\dfrac{\overline{w}}{|w|^2}\right) = \dfrac{u}{u^2+v^2} = 2$.

Hence $u = 2(u^2 + v^2)$, i.e. $u^2 + v^2 - \dfrac{u}{2} = 0$. Completing the square:
\[ \left(u - \frac14\right)^2 + v^2 = \frac{1}{16}, \qquad\text{i.e.}\qquad \left|w - \frac14\right| = \frac14. \]

The image is the \textbf{circle of centre $\left(\tfrac14, 0\right)$ and radius $\tfrac14$}, with the origin excluded (the point $w=0$ is never reached since $|w| = 1/|z|$ is finite). The circle passes through $O$, as expected: a line not through $O$ is mapped by inversion to a circle through $O$.
\end{corrige}

\begin{corrige}[Exercise 9 — Composite transformation and matrix (12 marks)]
\textbf{a)} The composite $T_2 \circ T_1$ (apply $T_1$ first, then $T_2$) is represented by
\[ NM = \begin{pmatrix} 2 & 0 \\ 0 & 2 \end{pmatrix}\begin{pmatrix} 0 & -1 \\ 1 & 0 \end{pmatrix} = \begin{pmatrix} 0 & -2 \\ 2 & 0 \end{pmatrix}. \]
\hfill [M1 A1]

\textbf{b)} $M$ is the rotation about $O$ through $+\dfrac{\pi}{2}$ (since $\begin{pmatrix}\cos\theta & -\sin\theta \\ \sin\theta & \cos\theta\end{pmatrix}$ with $\theta = \frac{\pi}{2}$), and $N$ is the enlargement (homothety) of centre $O$ and scale factor $2$. The composite is a \textbf{rotation of $\dfrac{\pi}{2}$ about $O$ combined with an enlargement of scale factor $2$} — a spiral similarity of ratio $2$ and angle $\dfrac{\pi}{2}$. \hfill [B1 B1]

\textbf{c)} The area scale factor is $|\det(NM)| = |0\cdot 0 - (-2)(2)| = 4$. (Equivalently, enlargement factor $2$ squares to $4$; rotation preserves area.) \hfill [M1 A1]

\textbf{d)} The inverse undoes the enlargement then the rotation: scale by $\tfrac12$ and rotate by $-\dfrac{\pi}{2}$. Its matrix:
\[ (NM)^{-1} = \frac{1}{4}\begin{pmatrix} 0 & 2 \\ -2 & 0 \end{pmatrix} = \begin{pmatrix} 0 & \tfrac12 \\[2pt] -\tfrac12 & 0 \end{pmatrix}. \]
Check: $\begin{pmatrix} 0 & -2 \\ 2 & 0 \end{pmatrix}\begin{pmatrix} 0 & \tfrac12 \\ -\tfrac12 & 0 \end{pmatrix} = \begin{pmatrix} 1 & 0 \\ 0 & 1 \end{pmatrix}$. \hfill [M1 A1 A1]

\textbf{Examiner expectations:} correct order of multiplication ($NM$, not $MN$) for the first mark; geometric description must mention \emph{both} components; area factor must be the \emph{determinant}, not the linear scale factor; the inverse matrix must be verified or derived from $\frac{1}{\det}\operatorname{adj}$. Common error: giving area factor $2$ instead of $4$.
\end{corrige}

\begin{corrige}[Exercise 10 — Similitude analysis (13 marks)]
Given $z' = (1-i)z + 3i$, a direct similitude $z' = az + b$ with $a = 1 - i$, $b = 3i$.

\textbf{a)} Ratio: $k = |a| = |1-i| = \sqrt{2}$. Angle: $\theta = \arg(1-i) = -\dfrac{\pi}{4}$. \hfill [B1 B1]

\textbf{b)} The centre is the unique fixed point (since $a \neq 1$): solve $c = (1-i)c + 3i$:
\[ c - (1-i)c = 3i \;\Longrightarrow\; ic = 3i \;\Longrightarrow\; c = 3. \]
Centre: the point $\Omega(3, 0)$. Check: $(1-i)\cdot 3 + 3i = 3 - 3i + 3i = 3$. \hfill [M1 A1]

\textbf{c)} Since $k = \sqrt2 \neq 1$ and $a \notin \mathbb{R}$, this is a \textbf{spiral similarity (similitude with rotation and enlargement)}: centre $\Omega(3,0)$, ratio $\sqrt2$, angle $-\dfrac{\pi}{4}$. It is neither a pure rotation ($k \neq 1$), nor a homothety ($a$ not real), nor a translation ($a \neq 1$). \hfill [B1 B1]

\textbf{d)} Images of the vertices:
\begin{align*}
0 &\mapsto 3i, \\
1 &\mapsto (1-i) + 3i = 1 + 2i, \\
i &\mapsto (1-i)i + 3i = i - i^2 + 3i = 1 + 4i.
\end{align*}
The image triangle has vertices $A'(0,3)$, $B'(1,2)$, $C'(1,4)$. \hfill [M1 A1 A1]

Check of similarity: original side from $0$ to $1$ has length $1$; image side from $3i$ to $1+2i$ has length $|1 - i| = \sqrt2 = k \cdot 1$. $\checkmark$ \hfill [B1]

\textbf{Examiner expectations:} quote $k = |a|$ and $\theta = \arg(a)$ explicitly; the fixed-point equation is the standard route to the centre; classification must be justified using the values of $k$ and $\theta$; images computed with correct complex arithmetic. A frequent error is solving $c = ac + b$ as $c = \frac{b}{1-a}$ and mishandling the division — here $\frac{3i}{i} = 3$ directly.
\end{corrige}

\begin{corrige}[Exercise 11 — Modulus inequality and rigid motions (15 marks)]
\textbf{a) Proof of the triangle inequality.} For any $z_1, z_2 \in \mathbb{C}$, using $|w|^2 = w\overline{w}$:
\begin{align*}
|z_1 + z_2|^2 &= (z_1+z_2)(\overline{z_1}+\overline{z_2}) = |z_1|^2 + z_1\overline{z_2} + \overline{z_1}z_2 + |z_2|^2 \\
&= |z_1|^2 + 2\operatorname{Re}(z_1\overline{z_2}) + |z_2|^2.
\end{align*}
Now $\operatorname{Re}(z_1\overline{z_2}) \le |z_1\overline{z_2}| = |z_1||z_2|$, so
\[ |z_1+z_2|^2 \le |z_1|^2 + 2|z_1||z_2| + |z_2|^2 = (|z_1|+|z_2|)^2. \]
Both sides are non-negative, so taking square roots: $|z_1+z_2| \le |z_1| + |z_2|$. $\blacksquare$

Equality holds iff $\operatorname{Re}(z_1\overline{z_2}) = |z_1\overline{z_2}|$, i.e. iff $z_1\overline{z_2}$ is a non-negative real, i.e. $z_1$ and $z_2$ have the same direction (one is a non-negative real multiple of the other). \hfill [M1 A1 A1 B1]

\textbf{b)} $|z - 3i| \le |z + 3|$: the set of points $P(z)$ at least as close to $A(0,3)$ as to $B(-3,0)$. The boundary is the perpendicular bisector of $[AB]$. With $z = x+iy$:
\[ x^2 + (y-3)^2 \le (x+3)^2 + y^2 \;\Longrightarrow\; -6y + 9 \le 6x + 9 \;\Longrightarrow\; y \ge -x. \]
Solution set: the closed half-plane $\{(x,y) : y \ge -x\}$, bounded by the line $y = -x$ and containing $A(0,3)$. \hfill [M1 A1 A1]

\begin{popfigure}
\begin{center}
\begin{tikzpicture}[scale=0.7]
\fill[popblueL] (-3,3) -- (3,-3) -- (3,3) -- cycle;
\draw[->,popline] (-3.4,0) -- (3.6,0) node[right]{\small Re};
\draw[->,popline] (0,-3.4) -- (0,3.6) node[above]{\small Im};
\draw[popblue,thick] (-3,3) -- (3,-3) node[below left]{\small $y=-x$};
\fill[popdark] (0,3) circle (1.8pt) node[above right]{\small $A(0,3)$};
\fill[popdark] (-3,0) circle (1.8pt) node[below left]{\small $B(-3,0)$};
\node[popdark] at (1.6,1.9) {\small solution};
\end{tikzpicture}
\end{center}
\legende{The half-plane $|z-3i| \le |z+3|$, i.e. $y \ge -x$ (boundary included).}
\end{popfigure}

\textbf{c)} \textbf{(i)} $z' = iz + 2 - i$: direct isometry ($a = i$, $|a| = 1$, $a \neq 1$), hence a \textbf{rotation}. It \textbf{preserves orientation}. Fixed point: $c = ic + 2 - i \Rightarrow c(1-i) = 2 - i \Rightarrow c = \dfrac{2-i}{1-i} = \dfrac{(2-i)(1+i)}{2} = \dfrac{3+i}{2}$. Angle: $\arg(i) = \dfrac{\pi}{2}$. So: rotation of centre $\left(\tfrac32, \tfrac12\right)$ and angle $+\dfrac{\pi}{2}$. \hfill [B1 M1 A1]

\textbf{(ii)} $z' = i\overline{z} + 1$: indirect isometry (conjugate form $z' = a\overline{z} + b$, $|a| = 1$), hence an \textbf{orientation-reversing} rigid motion. Fixed points: $z = i\overline{z} + 1$. With $z = x+iy$: $x + iy = i(x - iy) + 1 = y + 1 + ix$, giving $x = y + 1$ and $y = x$ — impossible ($x = x+1$). No fixed point, so it is a \textbf{glide reflection} (not a pure reflection).

\emph{Analysis of the glide reflection.} The transformation is $z' = i\overline{z} + 1$.
\begin{itemize}
\item The reflection part $z \mapsto i\overline{z}$ is reflection in the line through the origin with angle $\frac12\arg(i) = \frac{\pi}{4}$, i.e. the line $y = x$.
\item The translation vector is $b = 1$ (real axis). Decompose $b$ into components parallel and perpendicular to the axis $y=x$ (direction vector $(1,1)$):
\[ b_\parallel = \frac{1}{2}(1+i) \quad\text{(vector $(\tfrac12,\tfrac12)$)},\qquad b_\perp = \frac{1}{2}(1-i) \quad\text{(vector $(\tfrac12,-\tfrac12)$)}. \]
\item The \textbf{glide reflection axis} is the line $y=x$ shifted by $\frac12 b_\perp = (\tfrac14,-\tfrac14)$, giving the line $y = x - \tfrac12$.
\item The \textbf{glide vector} is $b_\parallel = (\tfrac12,\tfrac12)$ (parallel to the axis, length $\frac{1}{\sqrt2}$).
\end{itemize}
Conclusion: $z' = i\overline{z} + 1$ is the \textbf{glide reflection} with axis $y = x - \tfrac12$ and translation vector $(\tfrac12,\tfrac12)$; it \textbf{reverses orientation}. \hfill [B1 M1 A1 A1]

\textbf{Examiner expectations:} in (a) the proof must use $|w|^2 = w\bar w$ and the bound $\operatorname{Re}(w) \le |w|$; in (b) the boundary must be identified as a perpendicular bisector and the correct half-plane selected (test a point, e.g. $A$); in (c) the form $az+b$ vs $a\bar z + b$ decides orientation, and the existence or non-existence of fixed points distinguishes rotation from translation, and reflection from glide reflection.
\end{corrige}

\newpage

\coursec{Summary sheet}

\begin{popfigure}
\begin{tikzpicture}[mindmap, grow cyclic, every node/.style=concept, concept color=popblue,
    level 1/.append style={level distance=4.5cm, sibling angle=360/7},
    level 2/.append style={level distance=3cm, sibling angle=360/7}]
\node[concept] {Further Complex Numbers}
    child[concept color=poporange] {node[concept] {Forms \& De Moivre}}
    child[concept color=popgreen] {node[concept] {Applications of De Moivre}}
    child[concept color=poppurple] {node[concept] {Loci in Argand Diagram}}
    child[concept color=poppink] {node[concept] {Transformations $z\to w$}}
    child[concept color=popgold] {node[concept] {Plane Transformations \& Matrices}}
    child[concept color=popblue] {node[concept] {Similitudes}}
    child[concept color=popgreen] {node[concept] {Modulus Inequalities \& Rigid Motions}};
\end{tikzpicture}
\legende{Mind map of Chapter FMA02: Further Complex Numbers}
\end{popfigure}

\begin{bilan}
\textbf{Key Definitions}
\begin{itemize}
\item \cle{Complex number}: $z = x+iy = r(\cos\theta + i\sin\theta) = re^{i\theta}$.
\item \cle{Modulus}: $|z| = r = \sqrt{x^2+y^2}$.
\item \cle{Argument}: $\arg z = \theta$ (principal value in $(-\pi,\pi]$).
\item \cle{De Moivre's theorem}: $(\cos\theta + i\sin\theta)^n = \cos n\theta + i\sin n\theta$ for $n\in\mathbb{Z}$.
\item \cle{Locus}: set of points satisfying a condition on $z$.
\item \cle{Transformation $w=f(z)$}: mapping from $z$-plane to $w$-plane.
\item \cle{Similitude}: transformation $w = az + b$ or $w = a\bar{z}+b$ with $a\neq0$.
\item \cle{Rigid motion}: isometry preserving distances; orientation preserving (rotation, translation) or reversing (reflection, glide reflection).
\end{itemize}

\textbf{Key Formulas \& Theorems}
\begin{itemize}
\item Euler: $e^{i\theta}=\cos\theta+i\sin\theta$.
\item $z^n = r^n e^{in\theta}$; $z^{1/n} = r^{1/n} e^{i(\theta+2k\pi)/n}$, $k=0,\dots,n-1$.
\item $\cos n\theta$, $\sin n\theta$ expressed via binomial expansion of $(\cos\theta+i\sin\theta)^n$.
\item Loci: $|z-z_0|=r$ (circle), $|z-z_1|=|z-z_2|$ (perpendicular bisector), $\arg(z-z_0)=\alpha$ (half-line), $|z-z_1|+|z-z_2|=k$ (ellipse), $||z-z_1|-|z-z_2||=k$ (hyperbola).
\item Linear transformation: $w = \frac{az+b}{cz+d}$ (Möbius) but here elementary: translation $w=z+b$, rotation $w=e^{i\alpha}z$, enlargement $w=kz$, reflection $w=\bar{z}$, shear $w=z+\alpha\bar{z}$.
\item Matrix representation: translation not linear; rotation/enlargement/reflection/shear via $2\times2$ matrices.
\item Similitude matrix: $a+ib$ corresponds to $\begin{pmatrix}a&-b\\b&a\end{pmatrix}$ for direct; $\begin{pmatrix}a&b\\b&-a\end{pmatrix}$ for opposite.
\item Modulus inequalities: $|z-z_1|<|z-z_2|$ defines half-plane; $|z-z_1|+|z-z_2|<k$ interior of ellipse.
\end{itemize}

\textbf{Methods}
\begin{itemize}
\item \textbf{Convert forms}: Cartesian $\leftrightarrow$ Polar $\leftrightarrow$ Exponential.
\item \textbf{De Moivre}: expand $(\cos\theta+i\sin\theta)^n$, equate real/imaginary parts to get $\cos n\theta$, $\sin n\theta$; solve $z^n = w$ by taking $n$th roots.
\item \textbf{Loci}: translate condition into Cartesian $(x,y)$ or geometric description; sketch on Argand diagram.
\item \textbf{Transformations $z\to w$}: find image of lines/circles; use $w=f(z)$ and substitute $z=x+iy$; identify type (line, circle, etc.).
\item \textbf{Matrix representation}: write transformation as $\begin{pmatrix}x'\\y'\end{pmatrix}=M\begin{pmatrix}x\\y\end{pmatrix}+\mathbf{t}$; identify $M$ and $\mathbf{t}$.
\item \textbf{Similitudes}: express as $w=az+b$ (direct) or $w=a\bar{z}+b$ (opposite); find $a,b$ from given conditions; matrix form.
\item \textbf{Modulus inequalities}: square both sides carefully (preserve inequality direction); convert to Cartesian; identify region (half-plane, disc, ellipse, hyperbola).
\item \textbf{Rigid motions}: check orientation by determinant of matrix ($+1$ preserving, $-1$ reversing); composition of reflections yields rotation/translation.
\end{itemize}

\textbf{Common Mistakes}
\begin{itemize}
\item Forgetting to add $2k\pi$ when finding arguments of roots.
\item Confusing principal argument with general argument.
\item Incorrectly squaring modulus inequalities without considering sign (e.g., $|z-1|<|z+1|$ leads to $x>0$ after squaring).
\item Misidentifying transformation type: $w=1/z$ maps lines not through origin to circles through origin.
\item Mixing up direct and opposite similitudes: direct uses $z$, opposite uses $\bar{z}$.
\item Matrix representation: translation is not linear; must use homogeneous coordinates or separate vector.
\item Orientation: determinant of $2\times2$ matrix $>0$ preserves, $<0$ reverses.
\end{itemize}

\textbf{GCE Tips}
\begin{itemize}
\item Always show the conversion between forms explicitly; examiners award marks for correct modulus/argument.
\item For De Moivre expansions, write binomial expansion and separate real/imaginary parts clearly.
\item When sketching loci, indicate centre, radius, asymptotes, and shade region for inequalities.
\item For $w=f(z)$, find image of a generic point $z=x+iy$; eliminate $x,y$ to get equation in $u,v$ where $w=u+iv$.
\item In matrix questions, write the matrix in the form $\begin{pmatrix}a&-b\\b&a\end{pmatrix}$ for rotation+enlargement; for reflection use $\begin{pmatrix}\cos2\theta&\sin2\theta\\\sin2\theta&-\cos2\theta\end{pmatrix}$.
\item Similitude: if $w=az+b$, then $a=|a|e^{i\alpha}$ gives rotation by $\alpha$ and enlargement by $|a|$; $b$ gives translation.
\item Modulus inequalities: square only when both sides are non-negative (always true for moduli); then simplify to Cartesian.
\item Rigid motions: composition of two reflections in intersecting lines gives rotation; in parallel lines gives translation.
\item Time management: allocate ~12 minutes per 10-mark question; practice past papers.
\end{itemize}
\end{bilan}

\findecours

\end{document}
